% 逐条呈现账本事件；不把共享年序坐标当作具体行动。
\section{金与北方社会的逐年入口}
每一项均以其出处、行动者、空间线索和可追踪后果展开。材料没有说明地点、人数或地方反应时，正文保留这一限度。
\begin{textwindow}{本节材料怎样阅读}
《金史》卷一称“金之先，出靺鞨氏”；本纪和交聘表可核验金廷日期、名号与处置。它从中枢和王朝结局组织材料，迁徙者、守城者和乡村征发仍须同宋方、碑志和遗址材料分列。
\end{textwindow}

\subsection{1114年}
\noindent\eventanchor{JIN-1114-EVENT-078}{完颜阿骨打起兵反辽}\textbf{女真与辽。}
女真与辽 在 1114年 的材料痕迹是〔完颜阿骨打起兵反辽〕，出处为 《金史》相关本纪；《宋史》或《三朝北盟会编》《建炎以来系年要录》对应年条；金代碑志与蒙古、高丽材料（据事核）。材料未将地点细化到城址，不能用中枢所在地替它补足。

前一层压力可由以下线索辨认：“完颜阿骨打起兵反辽”发生在女真起兵、金辽宋联盟重组的具体压力与机会之中，相关行动者的选择不能脱离此前事件链解释。 这一处置留下的后续条件是：“完颜阿骨打起兵反辽”为后续政治、边防或资源配置留下条件；其社会影响与实际覆盖范围仍须以异类材料限定。

《金史》为元修，宋金战报和官修合法性语言各有宣传性；军数、俘获、控制范围和地方执行须以对方及物质材料限缩。 因此，金代国家形成、宋金边境、都城与金蒙转换研究 只能扩展问题，不能代替未保存的细节。

\subsection{1115年（金收国元年、辽天庆五年、宋政和五年；干支、月日待核；公元换算据各方本纪年号对照）}
\noindent\eventanchor{JIN-1115-EVENT-079}{完颜阿骨打称帝建金}\textbf{金。}
从 《金史》相关本纪；《宋史》或《三朝北盟会编》《建炎以来系年要录》对应年条；金代碑志与蒙古、高丽材料（据事核） 的 1115年（金收国元年、辽天庆五年、宋政和五年；干支、月日待核；公元换算据各方本纪年号对照） 条目看，〔完颜阿骨打称帝建金〕使 金 的一个具体行动进入叙事；材料未将地点细化到城址，不能用中枢所在地替它补足。

女真起兵、金辽宋联盟重组中前任统治的结束与宫廷、部族或官僚的权力安排相互交织，促成“完颜阿骨打称帝建金”这一继承节点。 记录所见的结果则是：继承并不自动消除既有矛盾；“完颜阿骨打称帝建金”后的任官、军权和对外策略需由后续条目追踪。 日期相邻不等于因果已经证实。

关于可说范围，须保留：《金史》为元修，宋金战报和官修合法性语言各有宣传性；军数、俘获、控制范围和地方执行须以对方及物质材料限缩。 进一步讨论可参照 金代国家形成、宋金边境、都城与金蒙转换研究

\subsection{1118年}
\noindent\eventanchor{JIN-1118-EVENT-080}{金与宋讨论共同攻辽}\textbf{金与宋。}
1118年 的记载将 金与宋 与〔金与宋讨论共同攻辽〕相连。此处可确认的是这项被记录的行动；材料未将地点细化到城址，不能用中枢所在地替它补足。

把本项放回事件链，能够看到的前提为：女真起兵、金辽宋联盟重组中既有的联盟、交通与补给压力，为“金与宋讨论共同攻辽”提供了直接的军事或动员背景。 而它所打开或收紧的下一步是：“金与宋讨论共同攻辽”改变了相关城镇、道路或政权间的军事选择；伤亡、俘获和控制范围仅可按各方材料分别确认。

证据的边界是：《金史》为元修，宋金战报和官修合法性语言各有宣传性；军数、俘获、控制范围和地方执行须以对方及物质材料限缩。 金代国家形成、宋金边境、都城与金蒙转换研究 可用于继续检验这一结论。

\subsection{1119年（宋宣和元年、金天辅3年；公元换算据两国年号对照）}
\noindent\eventanchor{SYD-1119-JIN-001}{丁巳，金人使李善慶來，遣趙有開報聘，至登州而還}\textbf{宋与金。}
〔丁巳，金人使李善慶來，遣趙有開報聘，至登州而還〕见于 《宋史》卷022〈宣和元年〉；《金史》本纪同年条（互校） 的 1119年（宋宣和元年、金天辅3年；公元换算据两国年号对照） 记录。宋与金 的选择因此有了可回查的时间位置；题名保留的城、路、水道或机构名称，是目前可以落地的空间线索。

它承接的条件是：本纪年条提供日期和中枢可见行动；前因不据单条补定。 随后可追踪的变化为：可回查同年及后续条目，地方影响仍须另证。

《宋史》为元末官修，本纪保存宋廷纪年与处置；战果、数字、地方执行及对方动机须以编年、会要、对方史料或地方材料互校。 因此，宋金战事、金代地方治理与碑志研究 只能扩展问题，不能代替未保存的细节。

\subsection{1120年（金天辅四年、宋宣和二年、辽天庆十年；干支、月日待核；公元换算据各方本纪年号对照）}
\noindent\eventanchor{JIN-1120-EVENT-081}{金宋达成海上之盟}\textbf{金与宋。}
1120年（金天辅四年、宋宣和二年、辽天庆十年；干支、月日待核；公元换算据各方本纪年号对照），《金史》相关本纪；《宋史》或《三朝北盟会编》《建炎以来系年要录》对应年条；金代碑志与蒙古、高丽材料（据事核） 所见的〔金宋达成海上之盟〕把 金与宋 的处置留在可复核的年序中；题名保留的城、路、水道或机构名称，是目前可以落地的空间线索。

前一层压力可由以下线索辨认：女真起兵、金辽宋联盟重组里的军事消耗、边境供给与使节往来构成谈判条件，促成“金宋达成海上之盟”所涉的协议或名义关系。 这一处置留下的后续条件是：“金宋达成海上之盟”重排了贡赐、使节或停战的制度接口；条款是否执行及边地实际控制不能由盟约文字直接推定。

关于可说范围，须保留：《金史》为元修，宋金战报和官修合法性语言各有宣传性；军数、俘获、控制范围和地方执行须以对方及物质材料限缩。 进一步讨论可参照 金代国家形成、宋金边境、都城与金蒙转换研究

\subsection{1120年（宋宣和二年、金天辅4年；公元换算据两国年号对照）}
\noindent\eventanchor{SYD-1120-JIN-002}{二月乙亥，遣趙良嗣使金國}\textbf{宋与金。}
宋与金 在 1120年（宋宣和二年、金天辅4年；公元换算据两国年号对照） 的材料痕迹是〔二月乙亥，遣趙良嗣使金國〕，出处为 《宋史》卷022〈宣和二年〉；《金史》本纪同年条（互校）。材料未将地点细化到城址，不能用中枢所在地替它补足。

本纪记录使节、贡赐或往来，不等于稳定统属关系。 记录所见的结果则是：可与对方编年和交通、文书材料核对其后续落实。 日期相邻不等于因果已经证实。

证据的边界是：《宋史》为元末官修，本纪保存宋廷纪年与处置；战果、数字、地方执行及对方动机须以编年、会要、对方史料或地方材料互校。 宋金战事、金代地方治理与碑志研究 可用于继续检验这一结论。

\noindent\eventanchor{SYD-1120-JIN-103}{九月壬寅，金人遣勃堇等來}\textbf{宋与金。}
从 《宋史》卷022〈宣和二年〉；《金史》本纪同年条（互校） 的 1120年（宋宣和二年、金天辅4年；公元换算据两国年号对照） 条目看，〔九月壬寅，金人遣勃堇等來〕使 宋与金 的一个具体行动进入叙事；材料未将地点细化到城址，不能用中枢所在地替它补足。

把本项放回事件链，能够看到的前提为：本纪年条提供日期和中枢可见行动；前因不据单条补定。 而它所打开或收紧的下一步是：可回查同年及后续条目，地方影响仍须另证。

《宋史》为元末官修，本纪保存宋廷纪年与处置；战果、数字、地方执行及对方动机须以编年、会要、对方史料或地方材料互校。 因此，宋金战事、金代地方治理与碑志研究 只能扩展问题，不能代替未保存的细节。

\subsection{1121年（宋宣和三年、金天辅5年；公元换算据两国年号对照）}
\noindent\eventanchor{SYD-1121-JIN-003}{丙午，金人再遣曷魯等來}\textbf{宋与金。}
1121年（宋宣和三年、金天辅5年；公元换算据两国年号对照） 的记载将 宋与金 与〔丙午，金人再遣曷魯等來〕相连。此处可确认的是这项被记录的行动；材料未将地点细化到城址，不能用中枢所在地替它补足。

它承接的条件是：本纪年条提供日期和中枢可见行动；前因不据单条补定。 随后可追踪的变化为：可回查同年及后续条目，地方影响仍须另证。

关于可说范围，须保留：《宋史》为元末官修，本纪保存宋廷纪年与处置；战果、数字、地方执行及对方动机须以编年、会要、对方史料或地方材料互校。 进一步讨论可参照 宋金战事、金代地方治理与碑志研究

\subsection{1122年}
\noindent\eventanchor{JIN-1122-EVENT-082}{金攻取辽南京燕京}\textbf{金与辽。}
〔金攻取辽南京燕京〕见于 《金史》相关本纪；《宋史》或《三朝北盟会编》《建炎以来系年要录》对应年条；金代碑志与蒙古、高丽材料（据事核） 的 1122年 记录。金与辽 的选择因此有了可回查的时间位置；题名保留的城、路、水道或机构名称，是目前可以落地的空间线索。

前一层压力可由以下线索辨认：女真起兵、金辽宋联盟重组中既有的联盟、交通与补给压力，为“金攻取辽南京燕京”提供了直接的军事或动员背景。 这一处置留下的后续条件是：“金攻取辽南京燕京”改变了相关城镇、道路或政权间的军事选择；伤亡、俘获和控制范围仅可按各方材料分别确认。

证据的边界是：《金史》为元修，宋金战报和官修合法性语言各有宣传性；军数、俘获、控制范围和地方执行须以对方及物质材料限缩。 金代国家形成、宋金边境、都城与金蒙转换研究 可用于继续检验这一结论。

\subsection{1122年（宋宣和四年、金天辅6年；公元换算据两国年号对照）}
\noindent\eventanchor{SYD-1122-JIN-004}{己未，金人遣徒孤且烏歇等來議師期}\textbf{宋与金。}
1122年（宋宣和四年、金天辅6年；公元换算据两国年号对照），《宋史》卷022〈宣和四年〉；《金史》本纪同年条（互校） 所见的〔己未，金人遣徒孤且烏歇等來議師期〕把 宋与金 的处置留在可复核的年序中；材料未将地点细化到城址，不能用中枢所在地替它补足。

本纪年条提供日期和中枢可见行动；前因不据单条补定。 记录所见的结果则是：可回查同年及后续条目，地方影响仍须另证。 日期相邻不等于因果已经证实。

《宋史》为元末官修，本纪保存宋廷纪年与处置；战果、数字、地方执行及对方动机须以编年、会要、对方史料或地方材料互校。 因此，宋金战事、金代地方治理与碑志研究 只能扩展问题，不能代替未保存的细节。

\noindent\eventanchor{SYD-1122-JIN-104}{甲戌，遣趙良嗣報聘于金國}\textbf{宋与金。}
宋与金 在 1122年（宋宣和四年、金天辅6年；公元换算据两国年号对照） 的材料痕迹是〔甲戌，遣趙良嗣報聘于金國〕，出处为 《宋史》卷022〈宣和四年〉；《金史》本纪同年条（互校）。材料未将地点细化到城址，不能用中枢所在地替它补足。

把本项放回事件链，能够看到的前提为：本纪年条提供日期和中枢可见行动；前因不据单条补定。 而它所打开或收紧的下一步是：可回查同年及后续条目，地方影响仍须另证。

关于可说范围，须保留：《宋史》为元末官修，本纪保存宋廷纪年与处置；战果、数字、地方执行及对方动机须以编年、会要、对方史料或地方材料互校。 进一步讨论可参照 宋金战事、金代地方治理与碑志研究

\noindent\eventanchor{SYD-1122-JIN-158}{金人來約夾攻，命童貫為河北、河東路宣撫使，屯兵于邊以應之，且招諭幽、燕}\textbf{宋与金。}
从 《宋史》卷022〈宣和四年〉；《金史》本纪同年条（互校） 的 1122年（宋宣和四年、金天辅6年；公元换算据两国年号对照） 条目看，〔金人來約夾攻，命童貫為河北、河東路宣撫使，屯兵于邊以應之，且招諭幽、燕〕使 宋与金 的一个具体行动进入叙事；题名保留的城、路、水道或机构名称，是目前可以落地的空间线索。

它承接的条件是：本纪年条记录政令或机构处置；实施背景须参照制度材料。 随后可追踪的变化为：可在同年及后续政令中追踪；条文不单独证明地方执行。

证据的边界是：《宋史》为元末官修，本纪保存宋廷纪年与处置；战果、数字、地方执行及对方动机须以编年、会要、对方史料或地方材料互校。 宋金战事、金代地方治理与碑志研究 可用于继续检验这一结论。

\noindent\eventanchor{SYD-1122-JIN-195}{戊寅，金人遣李靖等來許山前六州}\textbf{宋与金。}
1122年（宋宣和四年、金天辅6年；公元换算据两国年号对照） 的记载将 宋与金 与〔戊寅，金人遣李靖等來許山前六州〕相连。此处可确认的是这项被记录的行动；题名保留的城、路、水道或机构名称，是目前可以落地的空间线索。

前一层压力可由以下线索辨认：本纪年条提供日期和中枢可见行动；前因不据单条补定。 这一处置留下的后续条件是：可回查同年及后续条目，地方影响仍须另证。

《宋史》为元末官修，本纪保存宋廷纪年与处置；战果、数字、地方执行及对方动机须以编年、会要、对方史料或地方材料互校。 因此，宋金战事、金代地方治理与碑志研究 只能扩展问题，不能代替未保存的细节。

\noindent\eventanchor{SYD-1122-JIN-228}{戊子，遣趙良嗣報聘于金國}\textbf{宋与金。}
〔戊子，遣趙良嗣報聘于金國〕见于 《宋史》卷022〈宣和四年〉；《金史》本纪同年条（互校） 的 1122年（宋宣和四年、金天辅6年；公元换算据两国年号对照） 记录。宋与金 的选择因此有了可回查的时间位置；材料未将地点细化到城址，不能用中枢所在地替它补足。

本纪年条提供日期和中枢可见行动；前因不据单条补定。 记录所见的结果则是：可回查同年及后续条目，地方影响仍须另证。 日期相邻不等于因果已经证实。

关于可说范围，须保留：《宋史》为元末官修，本纪保存宋廷纪年与处置；战果、数字、地方执行及对方动机须以编年、会要、对方史料或地方材料互校。 进一步讨论可参照 宋金战事、金代地方治理与碑志研究

\subsection{1123年}
\noindent\eventanchor{JIN-1123-EVENT-083}{金将燕京部分地区交予宋并要求岁币}\textbf{金与宋。}
1123年，《金史》相关本纪；《宋史》或《三朝北盟会编》《建炎以来系年要录》对应年条；金代碑志与蒙古、高丽材料（据事核） 所见的〔金将燕京部分地区交予宋并要求岁币〕把 金与宋 的处置留在可复核的年序中；题名保留的城、路、水道或机构名称，是目前可以落地的空间线索。

把本项放回事件链，能够看到的前提为：“金将燕京部分地区交予宋并要求岁币”发生在女真起兵、金辽宋联盟重组的具体压力与机会之中，相关行动者的选择不能脱离此前事件链解释。 而它所打开或收紧的下一步是：“金将燕京部分地区交予宋并要求岁币”为后续政治、边防或资源配置留下条件；其社会影响与实际覆盖范围仍须以异类材料限定。

证据的边界是：《金史》为元修，宋金战报和官修合法性语言各有宣传性；军数、俘获、控制范围和地方执行须以对方及物质材料限缩。 金代国家形成、宋金边境、都城与金蒙转换研究 可用于继续检验这一结论。

\subsection{1123年（宋宣和五年、金天会1年；公元换算据两国年号对照）}
\noindent\eventanchor{SYD-1123-JIN-005}{庚子，童貫、蔡攸入燕，時燕之職官、富民、金帛、子女先為金人盡掠而去}\textbf{宋与金。}
宋与金 在 1123年（宋宣和五年、金天会1年；公元换算据两国年号对照） 的材料痕迹是〔庚子，童貫、蔡攸入燕，時燕之職官、富民、金帛、子女先為金人盡掠而去〕，出处为 《宋史》卷022〈宣和五年〉；《金史》本纪同年条（互校）。题名保留的城、路、水道或机构名称，是目前可以落地的空间线索。

它承接的条件是：本纪年条提供日期和中枢可见行动；前因不据单条补定。 随后可追踪的变化为：可回查同年及后续条目，地方影响仍须另证。

《宋史》为元末官修，本纪保存宋廷纪年与处置；战果、数字、地方执行及对方动机须以编年、会要、对方史料或地方材料互校。 因此，宋金战事、金代地方治理与碑志研究 只能扩展问题，不能代替未保存的细节。

\noindent\eventanchor{SYD-1123-JIN-105}{金主阿骨打殂，弟吳乞買立}\textbf{宋与金。}
从 《宋史》卷022〈宣和五年〉；《金史》本纪同年条（互校） 的 1123年（宋宣和五年、金天会1年；公元换算据两国年号对照） 条目看，〔金主阿骨打殂，弟吳乞買立〕使 宋与金 的一个具体行动进入叙事；材料未将地点细化到城址，不能用中枢所在地替它补足。

前一层压力可由以下线索辨认：本纪年条提供日期和中枢可见行动；前因不据单条补定。 这一处置留下的后续条件是：可回查同年及后续条目，地方影响仍须另证。

关于可说范围，须保留：《宋史》为元末官修，本纪保存宋廷纪年与处置；战果、数字、地方执行及对方动机须以编年、会要、对方史料或地方材料互校。 进一步讨论可参照 宋金战事、金代地方治理与碑志研究

\noindent\eventanchor{SYD-1123-JIN-159}{三月乙卯，金人再遣寧術割等來}\textbf{宋与金。}
1123年（宋宣和五年、金天会1年；公元换算据两国年号对照） 的记载将 宋与金 与〔三月乙卯，金人再遣寧術割等來〕相连。此处可确认的是这项被记录的行动；材料未将地点细化到城址，不能用中枢所在地替它补足。

本纪年条提供日期和中枢可见行动；前因不据单条补定。 记录所见的结果则是：可回查同年及后续条目，地方影响仍须另证。 日期相邻不等于因果已经证实。

证据的边界是：《宋史》为元末官修，本纪保存宋廷纪年与处置；战果、数字、地方执行及对方动机须以编年、会要、对方史料或地方材料互校。 宋金战事、金代地方治理与碑志研究 可用于继续检验这一结论。

\noindent\eventanchor{SYD-1123-JIN-196}{十二月乙巳，金人遣高居慶等來賀正旦}\textbf{宋与金。}
〔十二月乙巳，金人遣高居慶等來賀正旦〕见于 《宋史》卷022〈宣和五年〉；《金史》本纪同年条（互校） 的 1123年（宋宣和五年、金天会1年；公元换算据两国年号对照） 记录。宋与金 的选择因此有了可回查的时间位置；材料未将地点细化到城址，不能用中枢所在地替它补足。

把本项放回事件链，能够看到的前提为：本纪年条提供日期和中枢可见行动；前因不据单条补定。 而它所打开或收紧的下一步是：可回查同年及后续条目，地方影响仍须另证。

《宋史》为元末官修，本纪保存宋廷纪年与处置；战果、数字、地方执行及对方动机须以编年、会要、对方史料或地方材料互校。 因此，宋金战事、金代地方治理与碑志研究 只能扩展问题，不能代替未保存的细节。

\noindent\eventanchor{SYD-1123-JIN-229}{是月，金人取平州，張覺走燕山，金人索之甚急，命王安中縊殺，函其首送之}\textbf{宋与金。}
1123年（宋宣和五年、金天会1年；公元换算据两国年号对照），《宋史》卷022〈宣和五年〉；《金史》本纪同年条（互校） 所见的〔是月，金人取平州，張覺走燕山，金人索之甚急，命王安中縊殺，函其首送之〕把 宋与金 的处置留在可复核的年序中；题名保留的城、路、水道或机构名称，是目前可以落地的空间线索。

它承接的条件是：本纪年条记录政令或机构处置；实施背景须参照制度材料。 随后可追踪的变化为：可在同年及后续政令中追踪；条文不单独证明地方执行。

关于可说范围，须保留：《宋史》为元末官修，本纪保存宋廷纪年与处置；战果、数字、地方执行及对方动机须以编年、会要、对方史料或地方材料互校。 进一步讨论可参照 宋金战事、金代地方治理与碑志研究

\subsection{1124年（宋宣和六年、金天会2年；公元换算据两国年号对照）}
\noindent\eventanchor{SYD-1124-JIN-006}{庚子，金人遣富謨弼等以遺留物來獻}\textbf{宋与金。}
宋与金 在 1124年（宋宣和六年、金天会2年；公元换算据两国年号对照） 的材料痕迹是〔庚子，金人遣富謨弼等以遺留物來獻〕，出处为 《宋史》卷022〈宣和六年〉；《金史》本纪同年条（互校）。材料未将地点细化到城址，不能用中枢所在地替它补足。

前一层压力可由以下线索辨认：本纪年条提供日期和中枢可见行动；前因不据单条补定。 这一处置留下的后续条件是：可回查同年及后续条目，地方影响仍须另证。

证据的边界是：《宋史》为元末官修，本纪保存宋廷纪年与处置；战果、数字、地方执行及对方动机须以编年、会要、对方史料或地方材料互校。 宋金战事、金代地方治理与碑志研究 可用于继续检验这一结论。

\noindent\eventanchor{SYD-1124-JIN-106}{癸卯，金人遣使來告嗣位}\textbf{宋与金。}
从 《宋史》卷022〈宣和六年〉；《金史》本纪同年条（互校） 的 1124年（宋宣和六年、金天会2年；公元换算据两国年号对照） 条目看，〔癸卯，金人遣使來告嗣位〕使 宋与金 的一个具体行动进入叙事；材料未将地点细化到城址，不能用中枢所在地替它补足。

本纪记录使节、贡赐或往来，不等于稳定统属关系。 记录所见的结果则是：可与对方编年和交通、文书材料核对其后续落实。 日期相邻不等于因果已经证实。

《宋史》为元末官修，本纪保存宋廷纪年与处置；战果、数字、地方执行及对方动机须以编年、会要、对方史料或地方材料互校。 因此，宋金战事、金代地方治理与碑志研究 只能扩展问题，不能代替未保存的细节。

\noindent\eventanchor{SYD-1124-JIN-160}{六年春正月乙卯，為金主輟朝}\textbf{宋与金。}
1124年（宋宣和六年、金天会2年；公元换算据两国年号对照） 的记载将 宋与金 与〔六年春正月乙卯，為金主輟朝〕相连。此处可确认的是这项被记录的行动；材料未将地点细化到城址，不能用中枢所在地替它补足。

把本项放回事件链，能够看到的前提为：本纪年条提供日期和中枢可见行动；前因不据单条补定。 而它所打开或收紧的下一步是：可回查同年及后续条目，地方影响仍须另证。

关于可说范围，须保留：《宋史》为元末官修，本纪保存宋廷纪年与处置；战果、数字、地方执行及对方动机须以编年、会要、对方史料或地方材料互校。 进一步讨论可参照 宋金战事、金代地方治理与碑志研究

\noindent\eventanchor{SYD-1124-JIN-197}{秋七月戊子，遣許亢宗賀金國嗣位}\textbf{宋与金。}
〔秋七月戊子，遣許亢宗賀金國嗣位〕见于 《宋史》卷022〈宣和六年〉；《金史》本纪同年条（互校） 的 1124年（宋宣和六年、金天会2年；公元换算据两国年号对照） 记录。宋与金 的选择因此有了可回查的时间位置；材料未将地点细化到城址，不能用中枢所在地替它补足。

它承接的条件是：本纪年条提供日期和中枢可见行动；前因不据单条补定。 随后可追踪的变化为：可回查同年及后续条目，地方影响仍须另证。

证据的边界是：《宋史》为元末官修，本纪保存宋廷纪年与处置；战果、数字、地方执行及对方动机须以编年、会要、对方史料或地方材料互校。 宋金战事、金代地方治理与碑志研究 可用于继续检验这一结论。

\noindent\eventanchor{SYD-1124-JIN-230}{戊寅，遣連南夫弔祭金國}\textbf{宋与金。}
1124年（宋宣和六年、金天会2年；公元换算据两国年号对照），《宋史》卷022〈宣和六年〉；《金史》本纪同年条（互校） 所见的〔戊寅，遣連南夫弔祭金國〕把 宋与金 的处置留在可复核的年序中；材料未将地点细化到城址，不能用中枢所在地替它补足。

前一层压力可由以下线索辨认：本纪年条提供日期和中枢可见行动；前因不据单条补定。 这一处置留下的后续条件是：可回查同年及后续条目，地方影响仍须另证。

《宋史》为元末官修，本纪保存宋廷纪年与处置；战果、数字、地方执行及对方动机须以编年、会要、对方史料或地方材料互校。 因此，宋金战事、金代地方治理与碑志研究 只能扩展问题，不能代替未保存的细节。

\subsection{1125年（金天会三年、辽保大五年、宋宣和七年；干支、月日待核；公元换算据各方本纪年号对照）}
\noindent\eventanchor{JIN-1125-EVENT-084}{金军俘获辽天祚帝，辽朝覆亡}\textbf{金与辽。}
金与辽 在 1125年（金天会三年、辽保大五年、宋宣和七年；干支、月日待核；公元换算据各方本纪年号对照） 的材料痕迹是〔金军俘获辽天祚帝，辽朝覆亡〕，出处为 《金史》相关本纪；《宋史》或《三朝北盟会编》《建炎以来系年要录》对应年条；金代碑志与蒙古、高丽材料（据事核）。材料未将地点细化到城址，不能用中枢所在地替它补足。

女真起兵、金辽宋联盟重组中既有的联盟、交通与补给压力，为“金军俘获辽天祚帝，辽朝覆亡”提供了直接的军事或动员背景。 记录所见的结果则是：“金军俘获辽天祚帝，辽朝覆亡”改变了相关城镇、道路或政权间的军事选择；伤亡、俘获和控制范围仅可按各方材料分别确认。 日期相邻不等于因果已经证实。

关于可说范围，须保留：《金史》为元修，宋金战报和官修合法性语言各有宣传性；军数、俘获、控制范围和地方执行须以对方及物质材料限缩。 进一步讨论可参照 金代国家形成、宋金边境、都城与金蒙转换研究

\subsection{1125年（金天会三年、宋宣和七年；干支、月日待核；公元换算据双方本纪年号对照）}
\noindent\eventanchor{JIN-1125-EVENT-085}{金军分路南侵宋境}\textbf{金与宋。}
从 《金史》相关本纪；《宋史》或《三朝北盟会编》《建炎以来系年要录》对应年条；金代碑志与蒙古、高丽材料（据事核） 的 1125年（金天会三年、宋宣和七年；干支、月日待核；公元换算据双方本纪年号对照） 条目看，〔金军分路南侵宋境〕使 金与宋 的一个具体行动进入叙事；题名保留的城、路、水道或机构名称，是目前可以落地的空间线索。

把本项放回事件链，能够看到的前提为：女真起兵、金辽宋联盟重组中既有的联盟、交通与补给压力，为“金军分路南侵宋境”提供了直接的军事或动员背景。 而它所打开或收紧的下一步是：“金军分路南侵宋境”改变了相关城镇、道路或政权间的军事选择；伤亡、俘获和控制范围仅可按各方材料分别确认。

证据的边界是：《金史》为元修，宋金战报和官修合法性语言各有宣传性；军数、俘获、控制范围和地方执行须以对方及物质材料限缩。 金代国家形成、宋金边境、都城与金蒙转换研究 可用于继续检验这一结论。

\subsection{1125年（宋宣和七年、金天会3年；公元换算据两国年号对照）}
\noindent\eventanchor{SYD-1125-JIN-007}{己酉，中山奏金人斡離不、粘罕分兩道入攻}\textbf{宋与金。}
1125年（宋宣和七年、金天会3年；公元换算据两国年号对照） 的记载将 宋与金 与〔己酉，中山奏金人斡離不、粘罕分兩道入攻〕相连。此处可确认的是这项被记录的行动；题名保留的城、路、水道或机构名称，是目前可以落地的空间线索。

它承接的条件是：本纪从朝廷视角记录军政行动；出兵与战果需分辨。 随后可追踪的变化为：后续路线、控制与损失应与对方记录和遗址材料复核。

《宋史》为元末官修，本纪保存宋廷纪年与处置；战果、数字、地方执行及对方动机须以编年、会要、对方史料或地方材料互校。 因此，宋金战事、金代地方治理与碑志研究 只能扩展问题，不能代替未保存的细节。

\subsection{1126年（金天会四年、宋靖康元年；干支、月日待核；公元换算据双方本纪年号对照）}
\noindent\eventanchor{JIN-1126-EVENT-086}{金军第一次围汴京并迫使宋廷议和}\textbf{金与宋。}
〔金军第一次围汴京并迫使宋廷议和〕见于 《金史》相关本纪；《宋史》或《三朝北盟会编》《建炎以来系年要录》对应年条；金代碑志与蒙古、高丽材料（据事核） 的 1126年（金天会四年、宋靖康元年；干支、月日待核；公元换算据双方本纪年号对照） 记录。金与宋 的选择因此有了可回查的时间位置；题名保留的城、路、水道或机构名称，是目前可以落地的空间线索。

前一层压力可由以下线索辨认：女真起兵、金辽宋联盟重组里的军事消耗、边境供给与使节往来构成谈判条件，促成“金军第一次围汴京并迫使宋廷议和”所涉的协议或名义关系。 这一处置留下的后续条件是：“金军第一次围汴京并迫使宋廷议和”重排了贡赐、使节或停战的制度接口；条款是否执行及边地实际控制不能由盟约文字直接推定。

关于可说范围，须保留：《金史》为元修，宋金战报和官修合法性语言各有宣传性；军数、俘获、控制范围和地方执行须以对方及物质材料限缩。 进一步讨论可参照 金代国家形成、宋金边境、都城与金蒙转换研究

\subsection{1126年（宋靖康元年、金天会4年；公元换算据两国年号对照）}
\noindent\eventanchor{SYD-1126-JIN-008}{丙辰，妖人郭京用六甲法，盡令守禦人下城，大啟宣化門出攻金人，兵大敗}\textbf{宋与金。}
1126年（宋靖康元年、金天会4年；公元换算据两国年号对照），《宋史》卷023〈靖康元年〉；《金史》本纪同年条（互校） 所见的〔丙辰，妖人郭京用六甲法，盡令守禦人下城，大啟宣化門出攻金人，兵大敗〕把 宋与金 的处置留在可复核的年序中；题名保留的城、路、水道或机构名称，是目前可以落地的空间线索。

本纪从朝廷视角记录军政行动；出兵与战果需分辨。 记录所见的结果则是：后续路线、控制与损失应与对方记录和遗址材料复核。 日期相邻不等于因果已经证实。

证据的边界是：《宋史》为元末官修，本纪保存宋廷纪年与处置；战果、数字、地方执行及对方动机须以编年、会要、对方史料或地方材料互校。 宋金战事、金代地方治理与碑志研究 可用于继续检验这一结论。

\noindent\eventanchor{SYD-1126-JIN-107}{丙辰，詔有告奸人妄言金人復至以恐動居民者，賞之}\textbf{宋与金。}
宋与金 在 1126年（宋靖康元年、金天会4年；公元换算据两国年号对照） 的材料痕迹是〔丙辰，詔有告奸人妄言金人復至以恐動居民者，賞之〕，出处为 《宋史》卷023〈靖康元年〉；《金史》本纪同年条（互校）。材料未将地点细化到城址，不能用中枢所在地替它补足。

把本项放回事件链，能够看到的前提为：本纪年条记录政令或机构处置；实施背景须参照制度材料。 而它所打开或收紧的下一步是：可在同年及后续政令中追踪；条文不单独证明地方执行。

《宋史》为元末官修，本纪保存宋廷纪年与处置；战果、数字、地方执行及对方动机须以编年、会要、对方史料或地方材料互校。 因此，宋金战事、金代地方治理与碑志研究 只能扩展问题，不能代替未保存的细节。

\noindent\eventanchor{SYD-1126-JIN-161}{丁酉，金人陷真定，都鈐轄劉竧死之}\textbf{宋与金。}
从 《宋史》卷023〈靖康元年〉；《金史》本纪同年条（互校） 的 1126年（宋靖康元年、金天会4年；公元换算据两国年号对照） 条目看，〔丁酉，金人陷真定，都鈐轄劉竧死之〕使 宋与金 的一个具体行动进入叙事；题名保留的城、路、水道或机构名称，是目前可以落地的空间线索。

它承接的条件是：本纪年条提供日期和中枢可见行动；前因不据单条补定。 随后可追踪的变化为：可回查同年及后续条目，地方影响仍须另证。

关于可说范围，须保留：《宋史》为元末官修，本纪保存宋廷纪年与处置；战果、数字、地方执行及对方动机须以编年、会要、对方史料或地方材料互校。 进一步讨论可参照 宋金战事、金代地方治理与碑志研究

\noindent\eventanchor{SYD-1126-JIN-198}{東兵正將占沆與金人戰於交城縣，死之}\textbf{宋与金。}
1126年（宋靖康元年、金天会4年；公元换算据两国年号对照） 的记载将 宋与金 与〔東兵正將占沆與金人戰於交城縣，死之〕相连。此处可确认的是这项被记录的行动；题名保留的城、路、水道或机构名称，是目前可以落地的空间线索。

前一层压力可由以下线索辨认：本纪从朝廷视角记录军政行动；出兵与战果需分辨。 这一处置留下的后续条件是：后续路线、控制与损失应与对方记录和遗址材料复核。

证据的边界是：《宋史》为元末官修，本纪保存宋廷纪年与处置；战果、数字、地方执行及对方动机须以编年、会要、对方史料或地方材料互校。 宋金战事、金代地方治理与碑志研究 可用于继续检验这一结论。

\noindent\eventanchor{SYD-1126-JIN-231}{二月丁酉朔，命都統制姚平仲將兵夜襲金人軍，不克而奔}\textbf{宋与金。}
〔二月丁酉朔，命都統制姚平仲將兵夜襲金人軍，不克而奔〕见于 《宋史》卷023〈靖康元年〉；《金史》本纪同年条（互校） 的 1126年（宋靖康元年、金天会4年；公元换算据两国年号对照） 记录。宋与金 的选择因此有了可回查的时间位置；题名保留的城、路、水道或机构名称，是目前可以落地的空间线索。

本纪年条记录政令或机构处置；实施背景须参照制度材料。 记录所见的结果则是：可在同年及后续政令中追踪；条文不单独证明地方执行。 日期相邻不等于因果已经证实。

《宋史》为元末官修，本纪保存宋廷纪年与处置；战果、数字、地方执行及对方动机须以编年、会要、对方史料或地方材料互校。 因此，宋金战事、金代地方治理与碑志研究 只能扩展问题，不能代替未保存的细节。

\subsection{1127年（金天会五年、宋靖康二年；干支、月日待核；公元换算据双方本纪年号对照）}
\noindent\eventanchor{JIN-1127-EVENT-087}{金军攻陷汴京，俘徽、钦二帝}\textbf{金与宋。}
1127年（金天会五年、宋靖康二年；干支、月日待核；公元换算据双方本纪年号对照），《金史》相关本纪；《宋史》或《三朝北盟会编》《建炎以来系年要录》对应年条；金代碑志与蒙古、高丽材料（据事核） 所见的〔金军攻陷汴京，俘徽、钦二帝〕把 金与宋 的处置留在可复核的年序中；题名保留的城、路、水道或机构名称，是目前可以落地的空间线索。

把本项放回事件链，能够看到的前提为：女真起兵、金辽宋联盟重组中既有的联盟、交通与补给压力，为“金军攻陷汴京，俘徽、钦二帝”提供了直接的军事或动员背景。 而它所打开或收紧的下一步是：“金军攻陷汴京，俘徽、钦二帝”改变了相关城镇、道路或政权间的军事选择；伤亡、俘获和控制范围仅可按各方材料分别确认。

关于可说范围，须保留：《金史》为元修，宋金战报和官修合法性语言各有宣传性；军数、俘获、控制范围和地方执行须以对方及物质材料限缩。 进一步讨论可参照 金代国家形成、宋金边境、都城与金蒙转换研究

\subsection{1127年（南宋靖康二年、金天会5年；公元换算据两国年号对照）}
\noindent\eventanchor{SYD-1127-JIN-009}{丁酉，金人立張邦昌為楚帝}\textbf{南宋与金。}
南宋与金 在 1127年（南宋靖康二年、金天会5年；公元换算据两国年号对照） 的材料痕迹是〔丁酉，金人立張邦昌為楚帝〕，出处为 《宋史》卷023〈靖康二年〉；《金史》本纪同年条（互校）。材料未将地点细化到城址，不能用中枢所在地替它补足。

它承接的条件是：本纪年条提供日期和中枢可见行动；前因不据单条补定。 随后可追踪的变化为：可回查同年及后续条目，地方影响仍须另证。

证据的边界是：《宋史》为元末官修，本纪保存宋廷纪年与处置；战果、数字、地方执行及对方动机须以编年、会要、对方史料或地方材料互校。 宋金战事、金代地方治理与碑志研究 可用于继续检验这一结论。

\noindent\eventanchor{SYD-1127-JIN-108}{二年春正月辛卯朔，命濟王栩、景王杞出賀金軍，金人亦遣使入賀}\textbf{南宋与金。}
从 《宋史》卷023〈靖康二年〉；《金史》本纪同年条（互校） 的 1127年（南宋靖康二年、金天会5年；公元换算据两国年号对照） 条目看，〔二年春正月辛卯朔，命濟王栩、景王杞出賀金軍，金人亦遣使入賀〕使 南宋与金 的一个具体行动进入叙事；材料未将地点细化到城址，不能用中枢所在地替它补足。

前一层压力可由以下线索辨认：本纪年条记录政令或机构处置；实施背景须参照制度材料。 这一处置留下的后续条件是：可在同年及后续政令中追踪；条文不单独证明地方执行。

《宋史》为元末官修，本纪保存宋廷纪年与处置；战果、数字、地方执行及对方动机须以编年、会要、对方史料或地方材料互校。 因此，宋金战事、金代地方治理与碑志研究 只能扩展问题，不能代替未保存的细节。

\noindent\eventanchor{SYD-1127-JIN-162}{庚子，金人來取宗室，開封尹徐秉哲令民結保，毋藏匿}\textbf{南宋与金。}
1127年（南宋靖康二年、金天会5年；公元换算据两国年号对照） 的记载将 南宋与金 与〔庚子，金人來取宗室，開封尹徐秉哲令民結保，毋藏匿〕相连。此处可确认的是这项被记录的行动；材料未将地点细化到城址，不能用中枢所在地替它补足。

本纪年条提供日期和中枢可见行动；前因不据单条补定。 记录所见的结果则是：可回查同年及后续条目，地方影响仍须另证。 日期相邻不等于因果已经证实。

关于可说范围，须保留：《宋史》为元末官修，本纪保存宋廷纪年与处置；战果、数字、地方执行及对方动机须以编年、会要、对方史料或地方材料互校。 进一步讨论可参照 宋金战事、金代地方治理与碑志研究

\noindent\eventanchor{SYD-1127-JIN-199}{金人取至軍中，揆抗論，為所殺}\textbf{南宋与金。}
〔金人取至軍中，揆抗論，為所殺〕见于 《宋史》卷023〈靖康二年〉；《金史》本纪同年条（互校） 的 1127年（南宋靖康二年、金天会5年；公元换算据两国年号对照） 记录。南宋与金 的选择因此有了可回查的时间位置；材料未将地点细化到城址，不能用中枢所在地替它补足。

把本项放回事件链，能够看到的前提为：本纪从朝廷视角记录军政行动；出兵与战果需分辨。 而它所打开或收紧的下一步是：后续路线、控制与损失应与对方记录和遗址材料复核。

证据的边界是：《宋史》为元末官修，本纪保存宋廷纪年与处置；战果、数字、地方执行及对方动机须以编年、会要、对方史料或地方材料互校。 宋金战事、金代地方治理与碑志研究 可用于继续检验这一结论。

\noindent\eventanchor{SYD-1127-JIN-232}{辛未，金人逼上皇召皇后、皇太子入青城}\textbf{南宋与金。}
1127年（南宋靖康二年、金天会5年；公元换算据两国年号对照），《宋史》卷023〈靖康二年〉；《金史》本纪同年条（互校） 所见的〔辛未，金人逼上皇召皇后、皇太子入青城〕把 南宋与金 的处置留在可复核的年序中；题名保留的城、路、水道或机构名称，是目前可以落地的空间线索。

它承接的条件是：本纪年条提供日期和中枢可见行动；前因不据单条补定。 随后可追踪的变化为：可回查同年及后续条目，地方影响仍须另证。

《宋史》为元末官修，本纪保存宋廷纪年与处置；战果、数字、地方执行及对方动机须以编年、会要、对方史料或地方材料互校。 因此，宋金战事、金代地方治理与碑志研究 只能扩展问题，不能代替未保存的细节。

\subsection{1128年（南宋建炎二年、金天会6年；公元换算据两国年号对照）}
\noindent\eventanchor{SYD-1128-JIN-010}{丙申，復命宇文虛中為資政殿大學士，充金國祈請使}\textbf{南宋与金。}
南宋与金 在 1128年（南宋建炎二年、金天会6年；公元换算据两国年号对照） 的材料痕迹是〔丙申，復命宇文虛中為資政殿大學士，充金國祈請使〕，出处为 《宋史》卷025〈建炎二年〉；《金史》本纪同年条（互校）。材料未将地点细化到城址，不能用中枢所在地替它补足。

前一层压力可由以下线索辨认：本纪年条记录政令或机构处置；实施背景须参照制度材料。 这一处置留下的后续条件是：可在同年及后续政令中追踪；条文不单独证明地方执行。

关于可说范围，须保留：《宋史》为元末官修，本纪保存宋廷纪年与处置；战果、数字、地方执行及对方动机须以编年、会要、对方史料或地方材料互校。 进一步讨论可参照 宋金战事、金代地方治理与碑志研究

\noindent\eventanchor{SYD-1128-JIN-109}{丙子，金人陷淮寧府，守臣向子韶死之}\textbf{南宋与金。}
从 《宋史》卷025〈建炎二年〉；《金史》本纪同年条（互校） 的 1128年（南宋建炎二年、金天会6年；公元换算据两国年号对照） 条目看，〔丙子，金人陷淮寧府，守臣向子韶死之〕使 南宋与金 的一个具体行动进入叙事；题名保留的城、路、水道或机构名称，是目前可以落地的空间线索。

本纪年条提供日期和中枢可见行动；前因不据单条补定。 记录所见的结果则是：可回查同年及后续条目，地方影响仍须另证。 日期相邻不等于因果已经证实。

证据的边界是：《宋史》为元末官修，本纪保存宋廷纪年与处置；战果、数字、地方执行及对方动机须以编年、会要、对方史料或地方材料互校。 宋金战事、金代地方治理与碑志研究 可用于继续检验这一结论。

\noindent\eventanchor{SYD-1128-JIN-163}{丁丑，遣王貺等充金國軍前通問使}\textbf{南宋与金。}
1128年（南宋建炎二年、金天会6年；公元换算据两国年号对照） 的记载将 南宋与金 与〔丁丑，遣王貺等充金國軍前通問使〕相连。此处可确认的是这项被记录的行动；材料未将地点细化到城址，不能用中枢所在地替它补足。

把本项放回事件链，能够看到的前提为：本纪记录使节、贡赐或往来，不等于稳定统属关系。 而它所打开或收紧的下一步是：可与对方编年和交通、文书材料核对其后续落实。

《宋史》为元末官修，本纪保存宋廷纪年与处置；战果、数字、地方执行及对方动机须以编年、会要、对方史料或地方材料互校。 因此，宋金战事、金代地方治理与碑志研究 只能扩展问题，不能代替未保存的细节。

\noindent\eventanchor{SYD-1128-JIN-200}{丁未，東京留守統制官薛廣及金人戰於相州，敗死}\textbf{南宋与金。}
〔丁未，東京留守統制官薛廣及金人戰於相州，敗死〕见于 《宋史》卷025〈建炎二年〉；《金史》本纪同年条（互校） 的 1128年（南宋建炎二年、金天会6年；公元换算据两国年号对照） 记录。南宋与金 的选择因此有了可回查的时间位置；题名保留的城、路、水道或机构名称，是目前可以落地的空间线索。

它承接的条件是：本纪从朝廷视角记录军政行动；出兵与战果需分辨。 随后可追踪的变化为：后续路线、控制与损失应与对方记录和遗址材料复核。

关于可说范围，须保留：《宋史》为元末官修，本纪保存宋廷纪年与处置；战果、数字、地方执行及对方动机须以编年、会要、对方史料或地方材料互校。 进一步讨论可参照 宋金战事、金代地方治理与碑志研究

\noindent\eventanchor{SYD-1128-JIN-233}{二月丙辰，金人再犯東京，宗澤遣統制閻中立等拒之，中立戰死}\textbf{南宋与金。}
1128年（南宋建炎二年、金天会6年；公元换算据两国年号对照），《宋史》卷025〈建炎二年〉；《金史》本纪同年条（互校） 所见的〔二月丙辰，金人再犯東京，宗澤遣統制閻中立等拒之，中立戰死〕把 南宋与金 的处置留在可复核的年序中；题名保留的城、路、水道或机构名称，是目前可以落地的空间线索。

前一层压力可由以下线索辨认：本纪从朝廷视角记录军政行动；出兵与战果需分辨。 这一处置留下的后续条件是：后续路线、控制与损失应与对方记录和遗址材料复核。

证据的边界是：《宋史》为元末官修，本纪保存宋廷纪年与处置；战果、数字、地方执行及对方动机须以编年、会要、对方史料或地方材料互校。 宋金战事、金代地方治理与碑志研究 可用于继续检验这一结论。

\subsection{1129年}
\noindent\eventanchor{JIN-1129-EVENT-088}{完颜宗弼等军深入江淮}\textbf{金与南宋。}
金与南宋 在 1129年 的材料痕迹是〔完颜宗弼等军深入江淮〕，出处为 《金史》相关本纪；《宋史》或《三朝北盟会编》《建炎以来系年要录》对应年条；金代碑志与蒙古、高丽材料（据事核）。题名保留的城、路、水道或机构名称，是目前可以落地的空间线索。

“完颜宗弼等军深入江淮”发生在金在华北经营与宋金南北对峙的具体压力与机会之中，相关行动者的选择不能脱离此前事件链解释。 记录所见的结果则是：“完颜宗弼等军深入江淮”为后续政治、边防或资源配置留下条件；其社会影响与实际覆盖范围仍须以异类材料限定。 日期相邻不等于因果已经证实。

《金史》为元修，宋金战报和官修合法性语言各有宣传性；军数、俘获、控制范围和地方执行须以对方及物质材料限缩。 因此，金代国家形成、宋金边境、都城与金蒙转换研究 只能扩展问题，不能代替未保存的细节。

\subsection{1129年（南宋建炎三年、金天会7年；公元换算据两国年号对照）}
\noindent\eventanchor{SYD-1129-JIN-011}{丙辰，遣張邵等充金國軍前通問使}\textbf{南宋与金。}
从 《宋史》卷025〈建炎三年〉；《金史》本纪同年条（互校） 的 1129年（南宋建炎三年、金天会7年；公元换算据两国年号对照） 条目看，〔丙辰，遣張邵等充金國軍前通問使〕使 南宋与金 的一个具体行动进入叙事；材料未将地点细化到城址，不能用中枢所在地替它补足。

把本项放回事件链，能够看到的前提为：本纪记录使节、贡赐或往来，不等于稳定统属关系。 而它所打开或收紧的下一步是：可与对方编年和交通、文书材料核对其后续落实。

关于可说范围，须保留：《宋史》为元末官修，本纪保存宋廷纪年与处置；战果、数字、地方执行及对方动机须以编年、会要、对方史料或地方材料互校。 进一步讨论可参照 宋金战事、金代地方治理与碑志研究

\noindent\eventanchor{SYD-1129-JIN-110}{諜報金人治舟師，將由海道窺江、浙，遣韓世忠控守圌山、福山}\textbf{南宋与金。}
1129年（南宋建炎三年、金天会7年；公元换算据两国年号对照） 的记载将 南宋与金 与〔諜報金人治舟師，將由海道窺江、浙，遣韓世忠控守圌山、福山〕相连。此处可确认的是这项被记录的行动；题名保留的城、路、水道或机构名称，是目前可以落地的空间线索。

它承接的条件是：本纪年条提供日期和中枢可见行动；前因不据单条补定。 随后可追踪的变化为：可回查同年及后续条目，地方影响仍须另证。

证据的边界是：《宋史》为元末官修，本纪保存宋廷纪年与处置；战果、数字、地方执行及对方动机须以编年、会要、对方史料或地方材料互校。 宋金战事、金代地方治理与碑志研究 可用于继续检验这一结论。

\noindent\eventanchor{SYD-1129-JIN-164}{丁卯，遣杜時亮使金軍前}\textbf{南宋与金。}
〔丁卯，遣杜時亮使金軍前〕见于 《宋史》卷025〈建炎三年〉；《金史》本纪同年条（互校） 的 1129年（南宋建炎三年、金天会7年；公元换算据两国年号对照） 记录。南宋与金 的选择因此有了可回查的时间位置；材料未将地点细化到城址，不能用中枢所在地替它补足。

前一层压力可由以下线索辨认：本纪记录使节、贡赐或往来，不等于稳定统属关系。 这一处置留下的后续条件是：可与对方编年和交通、文书材料核对其后续落实。

《宋史》为元末官修，本纪保存宋廷纪年与处置；战果、数字、地方执行及对方动机须以编年、会要、对方史料或地方材料互校。 因此，宋金战事、金代地方治理与碑志研究 只能扩展问题，不能代替未保存的细节。

\noindent\eventanchor{SYD-1129-JIN-201}{丁巳，金人犯泰州，守臣曾班以城降}\textbf{南宋与金。}
1129年（南宋建炎三年、金天会7年；公元换算据两国年号对照），《宋史》卷025〈建炎三年〉；《金史》本纪同年条（互校） 所见的〔丁巳，金人犯泰州，守臣曾班以城降〕把 南宋与金 的处置留在可复核的年序中；题名保留的城、路、水道或机构名称，是目前可以落地的空间线索。

本纪年条记录政令或机构处置；实施背景须参照制度材料。 记录所见的结果则是：可在同年及后续政令中追踪；条文不单独证明地方执行。 日期相邻不等于因果已经证实。

关于可说范围，须保留：《宋史》为元末官修，本纪保存宋廷纪年与处置；战果、数字、地方执行及对方动机须以编年、会要、对方史料或地方材料互校。 进一步讨论可参照 宋金战事、金代地方治理与碑志研究

\noindent\eventanchor{SYD-1129-JIN-234}{丁酉，遣崔縱使金軍前}\textbf{南宋与金。}
南宋与金 在 1129年（南宋建炎三年、金天会7年；公元换算据两国年号对照） 的材料痕迹是〔丁酉，遣崔縱使金軍前〕，出处为 《宋史》卷025〈建炎三年〉；《金史》本纪同年条（互校）。材料未将地点细化到城址，不能用中枢所在地替它补足。

把本项放回事件链，能够看到的前提为：本纪记录使节、贡赐或往来，不等于稳定统属关系。 而它所打开或收紧的下一步是：可与对方编年和交通、文书材料核对其后续落实。

证据的边界是：《宋史》为元末官修，本纪保存宋廷纪年与处置；战果、数字、地方执行及对方动机须以编年、会要、对方史料或地方材料互校。 宋金战事、金代地方治理与碑志研究 可用于继续检验这一结论。

\subsection{1130年}
\noindent\eventanchor{JIN-1130-EVENT-089}{金扶植刘豫政权经营中原}\textbf{金与南宋。}
从 《金史》相关本纪；《宋史》或《三朝北盟会编》《建炎以来系年要录》对应年条；金代碑志与蒙古、高丽材料（据事核） 的 1130年 条目看，〔金扶植刘豫政权经营中原〕使 金与南宋 的一个具体行动进入叙事；材料未将地点细化到城址，不能用中枢所在地替它补足。

它承接的条件是：“金扶植刘豫政权经营中原”发生在金在华北经营与宋金南北对峙的具体压力与机会之中，相关行动者的选择不能脱离此前事件链解释。 随后可追踪的变化为：“金扶植刘豫政权经营中原”为后续政治、边防或资源配置留下条件；其社会影响与实际覆盖范围仍须以异类材料限定。

《金史》为元修，宋金战报和官修合法性语言各有宣传性；军数、俘获、控制范围和地方执行须以对方及物质材料限缩。 因此，金代国家形成、宋金边境、都城与金蒙转换研究 只能扩展问题，不能代替未保存的细节。

\subsection{1130年（南宋建炎四年、金天会8年；公元换算据两国年号对照）}
\noindent\eventanchor{SYD-1130-JIN-012}{丙辰，金人犯終南縣，經略使鄭恩戰敗，死之}\textbf{南宋与金。}
1130年（南宋建炎四年、金天会8年；公元换算据两国年号对照） 的记载将 南宋与金 与〔丙辰，金人犯終南縣，經略使鄭恩戰敗，死之〕相连。此处可确认的是这项被记录的行动；材料未将地点细化到城址，不能用中枢所在地替它补足。

前一层压力可由以下线索辨认：本纪从朝廷视角记录军政行动；出兵与战果需分辨。 这一处置留下的后续条件是：后续路线、控制与损失应与对方记录和遗址材料复核。

关于可说范围，须保留：《宋史》为元末官修，本纪保存宋廷纪年与处置；战果、数字、地方执行及对方动机须以编年、会要、对方史料或地方材料互校。 进一步讨论可参照 宋金战事、金代地方治理与碑志研究

\noindent\eventanchor{SYD-1130-JIN-111}{丙申，詔劉光世節制諸鎮，守禦通、泰州，伺便襲金人過淮}\textbf{南宋与金。}
〔丙申，詔劉光世節制諸鎮，守禦通、泰州，伺便襲金人過淮〕见于 《宋史》卷026〈建炎四年〉；《金史》本纪同年条（互校） 的 1130年（南宋建炎四年、金天会8年；公元换算据两国年号对照） 记录。南宋与金 的选择因此有了可回查的时间位置；题名保留的城、路、水道或机构名称，是目前可以落地的空间线索。

本纪年条记录政令或机构处置；实施背景须参照制度材料。 记录所见的结果则是：可在同年及后续政令中追踪；条文不单独证明地方执行。 日期相邻不等于因果已经证实。

证据的边界是：《宋史》为元末官修，本纪保存宋廷纪年与处置；战果、数字、地方执行及对方动机须以编年、会要、对方史料或地方材料互校。 宋金战事、金代地方治理与碑志研究 可用于继续检验这一结论。

\noindent\eventanchor{SYD-1130-JIN-165}{丙戌，金人自臨安退兵，命劉光世率兵追之}\textbf{南宋与金。}
1130年（南宋建炎四年、金天会8年；公元换算据两国年号对照），《宋史》卷026〈建炎四年〉；《金史》本纪同年条（互校） 所见的〔丙戌，金人自臨安退兵，命劉光世率兵追之〕把 南宋与金 的处置留在可复核的年序中；材料未将地点细化到城址，不能用中枢所在地替它补足。

把本项放回事件链，能够看到的前提为：本纪年条记录政令或机构处置；实施背景须参照制度材料。 而它所打开或收紧的下一步是：可在同年及后续政令中追踪；条文不单独证明地方执行。

《宋史》为元末官修，本纪保存宋廷纪年与处置；战果、数字、地方执行及对方动机须以编年、会要、对方史料或地方材料互校。 因此，宋金战事、金代地方治理与碑志研究 只能扩展问题，不能代替未保存的细节。

\noindent\eventanchor{SYD-1130-JIN-202}{丁亥，金人陷汴京，權留守上官悟出奔，為盜所殺}\textbf{南宋与金。}
南宋与金 在 1130年（南宋建炎四年、金天会8年；公元换算据两国年号对照） 的材料痕迹是〔丁亥，金人陷汴京，權留守上官悟出奔，為盜所殺〕，出处为 《宋史》卷026〈建炎四年〉；《金史》本纪同年条（互校）。题名保留的城、路、水道或机构名称，是目前可以落地的空间线索。

它承接的条件是：本纪从朝廷视角记录军政行动；出兵与战果需分辨。 随后可追踪的变化为：后续路线、控制与损失应与对方记录和遗址材料复核。

关于可说范围，须保留：《宋史》为元末官修，本纪保存宋廷纪年与处置；战果、数字、地方执行及对方动机须以编年、会要、对方史料或地方材料互校。 进一步讨论可参照 宋金战事、金代地方治理与碑志研究

\noindent\eventanchor{SYD-1130-JIN-235}{丁卯，金人立劉豫為帝，國號齊}\textbf{南宋与金。}
从 《宋史》卷026〈建炎四年〉；《金史》本纪同年条（互校） 的 1130年（南宋建炎四年、金天会8年；公元换算据两国年号对照） 条目看，〔丁卯，金人立劉豫為帝，國號齊〕使 南宋与金 的一个具体行动进入叙事；材料未将地点细化到城址，不能用中枢所在地替它补足。

前一层压力可由以下线索辨认：本纪年条提供日期和中枢可见行动；前因不据单条补定。 这一处置留下的后续条件是：可回查同年及后续条目，地方影响仍须另证。

证据的边界是：《宋史》为元末官修，本纪保存宋廷纪年与处置；战果、数字、地方执行及对方动机须以编年、会要、对方史料或地方材料互校。 宋金战事、金代地方治理与碑志研究 可用于继续检验这一结论。

\subsection{1131年（南宋紹興元年、金天会9年；公元换算据两国年号对照）}
\noindent\eventanchor{SYD-1131-JIN-013}{金人犯和尚原，吳玠擊敗之}\textbf{南宋与金。}
1131年（南宋紹興元年、金天会9年；公元换算据两国年号对照） 的记载将 南宋与金 与〔金人犯和尚原，吳玠擊敗之〕相连。此处可确认的是这项被记录的行动；材料未将地点细化到城址，不能用中枢所在地替它补足。

本纪从朝廷视角记录军政行动；出兵与战果需分辨。 记录所见的结果则是：后续路线、控制与损失应与对方记录和遗址材料复核。 日期相邻不等于因果已经证实。

《宋史》为元末官修，本纪保存宋廷纪年与处置；战果、数字、地方执行及对方动机须以编年、会要、对方史料或地方材料互校。 因此，宋金战事、金代地方治理与碑志研究 只能扩展问题，不能代替未保存的细节。

\noindent\eventanchor{SYD-1131-JIN-112}{金人犯秦州，吳玠擊敗之}\textbf{南宋与金。}
〔金人犯秦州，吳玠擊敗之〕见于 《宋史》卷026〈紹興元年〉；《金史》本纪同年条（互校） 的 1131年（南宋紹興元年、金天会9年；公元换算据两国年号对照） 记录。南宋与金 的选择因此有了可回查的时间位置；题名保留的城、路、水道或机构名称，是目前可以落地的空间线索。

把本项放回事件链，能够看到的前提为：本纪从朝廷视角记录军政行动；出兵与战果需分辨。 而它所打开或收紧的下一步是：后续路线、控制与损失应与对方记录和遗址材料复核。

关于可说范围，须保留：《宋史》为元末官修，本纪保存宋廷纪年与处置；战果、数字、地方执行及对方动机须以编年、会要、对方史料或地方材料互校。 进一步讨论可参照 宋金战事、金代地方治理与碑志研究

\noindent\eventanchor{SYD-1131-JIN-166}{金人迫興州，張浚退保閬州，以端明殿學士張深為四川制置使，及參議軍事劉子羽趨益昌}\textbf{南宋与金。}
1131年（南宋紹興元年、金天会9年；公元换算据两国年号对照），《宋史》卷026〈紹興元年〉；《金史》本纪同年条（互校） 所见的〔金人迫興州，張浚退保閬州，以端明殿學士張深為四川制置使，及參議軍事劉子羽趨益昌〕把 南宋与金 的处置留在可复核的年序中；题名保留的城、路、水道或机构名称，是目前可以落地的空间线索。

它承接的条件是：本纪年条记录政令或机构处置；实施背景须参照制度材料。 随后可追踪的变化为：可在同年及后续政令中追踪；条文不单独证明地方执行。

证据的边界是：《宋史》为元末官修，本纪保存宋廷纪年与处置；战果、数字、地方执行及对方动机须以编年、会要、对方史料或地方材料互校。 宋金战事、金代地方治理与碑志研究 可用于继续检验这一结论。

\noindent\eventanchor{SYD-1131-JIN-203}{是月，金人攻張榮縮頭湖水砦，榮擊敗之，來告捷，劉光世以榮知泰州}\textbf{南宋与金。}
南宋与金 在 1131年（南宋紹興元年、金天会9年；公元换算据两国年号对照） 的材料痕迹是〔是月，金人攻張榮縮頭湖水砦，榮擊敗之，來告捷，劉光世以榮知泰州〕，出处为 《宋史》卷026〈紹興元年〉；《金史》本纪同年条（互校）。题名保留的城、路、水道或机构名称，是目前可以落地的空间线索。

前一层压力可由以下线索辨认：本纪保留灾害或工程的报告地点与朝廷处置；灾情范围需另证。 这一处置留下的后续条件是：可与地方文书、河工和环境材料比较，不将单点记录外推为全域因果。

《宋史》为元末官修，本纪保存宋廷纪年与处置；战果、数字、地方执行及对方动机须以编年、会要、对方史料或地方材料互校。 因此，宋金战事、金代地方治理与碑志研究 只能扩展问题，不能代替未保存的细节。

\noindent\eventanchor{SYD-1131-JIN-236}{趙彬及金人合兵圍慶陽府，守臣楊可升擊敗之}\textbf{南宋与金。}
从 《宋史》卷026〈紹興元年〉；《金史》本纪同年条（互校） 的 1131年（南宋紹興元年、金天会9年；公元换算据两国年号对照） 条目看，〔趙彬及金人合兵圍慶陽府，守臣楊可升擊敗之〕使 南宋与金 的一个具体行动进入叙事；题名保留的城、路、水道或机构名称，是目前可以落地的空间线索。

本纪年条记录政令或机构处置；实施背景须参照制度材料。 记录所见的结果则是：可在同年及后续政令中追踪；条文不单独证明地方执行。 日期相邻不等于因果已经证实。

关于可说范围，须保留：《宋史》为元末官修，本纪保存宋廷纪年与处置；战果、数字、地方执行及对方动机须以编年、会要、对方史料或地方材料互校。 进一步讨论可参照 宋金战事、金代地方治理与碑志研究

\subsection{1132年（南宋紹興二年、金天会10年；公元换算据两国年号对照）}
\noindent\eventanchor{SYD-1132-JIN-014}{庚子，金人攻方山原，陝西統制楊政援之，金兵引去}\textbf{南宋与金。}
1132年（南宋紹興二年、金天会10年；公元换算据两国年号对照） 的记载将 南宋与金 与〔庚子，金人攻方山原，陝西統制楊政援之，金兵引去〕相连。此处可确认的是这项被记录的行动；题名保留的城、路、水道或机构名称，是目前可以落地的空间线索。

把本项放回事件链，能够看到的前提为：本纪从朝廷视角记录军政行动；出兵与战果需分辨。 而它所打开或收紧的下一步是：后续路线、控制与损失应与对方记录和遗址材料复核。

证据的边界是：《宋史》为元末官修，本纪保存宋廷纪年与处置；战果、数字、地方执行及对方动机须以编年、会要、对方史料或地方材料互校。 宋金战事、金代地方治理与碑志研究 可用于继续检验这一结论。

\noindent\eventanchor{SYD-1132-JIN-113}{甲寅，金人復自水洛城來攻，楊政等又敗之}\textbf{南宋与金。}
〔甲寅，金人復自水洛城來攻，楊政等又敗之〕见于 《宋史》卷027〈紹興二年〉；《金史》本纪同年条（互校） 的 1132年（南宋紹興二年、金天会10年；公元换算据两国年号对照） 记录。南宋与金 的选择因此有了可回查的时间位置；题名保留的城、路、水道或机构名称，是目前可以落地的空间线索。

它承接的条件是：本纪年条记录政令或机构处置；实施背景须参照制度材料。 随后可追踪的变化为：可在同年及后续政令中追踪；条文不单独证明地方执行。

《宋史》为元末官修，本纪保存宋廷纪年与处置；战果、数字、地方执行及对方动机须以编年、会要、对方史料或地方材料互校。 因此，宋金战事、金代地方治理与碑志研究 只能扩展问题，不能代替未保存的细节。

\noindent\eventanchor{SYD-1132-JIN-167}{金人侵熙、秦，關師古擊敗之}\textbf{南宋与金。}
1132年（南宋紹興二年、金天会10年；公元换算据两国年号对照），《宋史》卷027〈紹興二年〉；《金史》本纪同年条（互校） 所见的〔金人侵熙、秦，關師古擊敗之〕把 南宋与金 的处置留在可复核的年序中；题名保留的城、路、水道或机构名称，是目前可以落地的空间线索。

前一层压力可由以下线索辨认：本纪从朝廷视角记录军政行动；出兵与战果需分辨。 这一处置留下的后续条件是：后续路线、控制与损失应与对方记录和遗址材料复核。

关于可说范围，须保留：《宋史》为元末官修，本纪保存宋廷纪年与处置；战果、数字、地方执行及对方动机须以编年、会要、对方史料或地方材料互校。 进一步讨论可参照 宋金战事、金代地方治理与碑志研究

\noindent\eventanchor{SYD-1132-JIN-204}{金人陷慶陽府，執楊可升，降之}\textbf{南宋与金。}
南宋与金 在 1132年（南宋紹興二年、金天会10年；公元换算据两国年号对照） 的材料痕迹是〔金人陷慶陽府，執楊可升，降之〕，出处为 《宋史》卷027〈紹興二年〉；《金史》本纪同年条（互校）。题名保留的城、路、水道或机构名称，是目前可以落地的空间线索。

本纪年条记录政令或机构处置；实施背景须参照制度材料。 记录所见的结果则是：可在同年及后续政令中追踪；条文不单独证明地方执行。 日期相邻不等于因果已经证实。

证据的边界是：《宋史》为元末官修，本纪保存宋廷纪年与处置；战果、数字、地方执行及对方动机须以编年、会要、对方史料或地方材料互校。 宋金战事、金代地方治理与碑志研究 可用于继续检验这一结论。

\noindent\eventanchor{SYD-1132-JIN-237}{遣潘致堯等為金國軍前通問使，附茶藥金幣進兩宮}\textbf{南宋与金。}
从 《宋史》卷027〈紹興二年〉；《金史》本纪同年条（互校） 的 1132年（南宋紹興二年、金天会10年；公元换算据两国年号对照） 条目看，〔遣潘致堯等為金國軍前通問使，附茶藥金幣進兩宮〕使 南宋与金 的一个具体行动进入叙事；题名保留的城、路、水道或机构名称，是目前可以落地的空间线索。

把本项放回事件链，能够看到的前提为：本纪记录使节、贡赐或往来，不等于稳定统属关系。 而它所打开或收紧的下一步是：可与对方编年和交通、文书材料核对其后续落实。

《宋史》为元末官修，本纪保存宋廷纪年与处置；战果、数字、地方执行及对方动机须以编年、会要、对方史料或地方材料互校。 因此，宋金战事、金代地方治理与碑志研究 只能扩展问题，不能代替未保存的细节。

\subsection{1133年（南宋紹興三年、金天会11年；公元换算据两国年号对照）}
\noindent\eventanchor{SYD-1133-JIN-015}{丙子，王彥復金州，金兵棄均、房去}\textbf{南宋与金。}
1133年（南宋紹興三年、金天会11年；公元换算据两国年号对照） 的记载将 南宋与金 与〔丙子，王彥復金州，金兵棄均、房去〕相连。此处可确认的是这项被记录的行动；题名保留的城、路、水道或机构名称，是目前可以落地的空间线索。

它承接的条件是：本纪年条记录政令或机构处置；实施背景须参照制度材料。 随后可追踪的变化为：可在同年及后续政令中追踪；条文不单独证明地方执行。

关于可说范围，须保留：《宋史》为元末官修，本纪保存宋廷纪年与处置；战果、数字、地方执行及对方动机须以编年、会要、对方史料或地方材料互校。 进一步讨论可参照 宋金战事、金代地方治理与碑志研究

\noindent\eventanchor{SYD-1133-JIN-114}{己巳，金人遣兵援劉豫，李橫敗走，潁昌復陷}\textbf{南宋与金。}
〔己巳，金人遣兵援劉豫，李橫敗走，潁昌復陷〕见于 《宋史》卷027〈紹興三年〉；《金史》本纪同年条（互校） 的 1133年（南宋紹興三年、金天会11年；公元换算据两国年号对照） 记录。南宋与金 的选择因此有了可回查的时间位置；材料未将地点细化到城址，不能用中枢所在地替它补足。

前一层压力可由以下线索辨认：本纪年条记录政令或机构处置；实施背景须参照制度材料。 这一处置留下的后续条件是：可在同年及后续政令中追踪；条文不单独证明地方执行。

证据的边界是：《宋史》为元末官修，本纪保存宋廷纪年与处置；战果、数字、地方执行及对方动机须以编年、会要、对方史料或地方材料互校。 宋金战事、金代地方治理与碑志研究 可用于继续检验这一结论。

\noindent\eventanchor{SYD-1133-JIN-168}{己酉，金國元帥府遣李永壽、王翊來見}\textbf{南宋与金。}
1133年（南宋紹興三年、金天会11年；公元换算据两国年号对照），《宋史》卷027〈紹興三年〉；《金史》本纪同年条（互校） 所见的〔己酉，金國元帥府遣李永壽、王翊來見〕把 南宋与金 的处置留在可复核的年序中；题名保留的城、路、水道或机构名称，是目前可以落地的空间线索。

本纪年条提供日期和中枢可见行动；前因不据单条补定。 记录所见的结果则是：可回查同年及后续条目，地方影响仍须另证。 日期相邻不等于因果已经证实。

《宋史》为元末官修，本纪保存宋廷纪年与处置；战果、数字、地方执行及对方动机须以编年、会要、对方史料或地方材料互校。 因此，宋金战事、金代地方治理与碑志研究 只能扩展问题，不能代替未保存的细节。

\noindent\eventanchor{SYD-1133-JIN-205}{時金兵深入至金牛鎮，疑有伏，由褒斜谷引兵還興元，吳玠、劉子羽追擊其後，殺獲甚眾}\textbf{南宋与金。}
南宋与金 在 1133年（南宋紹興三年、金天会11年；公元换算据两国年号对照） 的材料痕迹是〔時金兵深入至金牛鎮，疑有伏，由褒斜谷引兵還興元，吳玠、劉子羽追擊其後，殺獲甚眾〕，出处为 《宋史》卷027〈紹興三年〉；《金史》本纪同年条（互校）。材料未将地点细化到城址，不能用中枢所在地替它补足。

把本项放回事件链，能够看到的前提为：本纪从朝廷视角记录军政行动；出兵与战果需分辨。 而它所打开或收紧的下一步是：后续路线、控制与损失应与对方记录和遗址材料复核。

关于可说范围，须保留：《宋史》为元末官修，本纪保存宋廷纪年与处置；战果、数字、地方执行及对方动机须以编年、会要、对方史料或地方材料互校。 进一步讨论可参照 宋金战事、金代地方治理与碑志研究

\noindent\eventanchor{SYD-1133-JIN-238}{是月，金人圍方山原，王似命吳玠發兵救之}\textbf{南宋与金。}
从 《宋史》卷027〈紹興三年〉；《金史》本纪同年条（互校） 的 1133年（南宋紹興三年、金天会11年；公元换算据两国年号对照） 条目看，〔是月，金人圍方山原，王似命吳玠發兵救之〕使 南宋与金 的一个具体行动进入叙事；题名保留的城、路、水道或机构名称，是目前可以落地的空间线索。

它承接的条件是：本纪年条记录政令或机构处置；实施背景须参照制度材料。 随后可追踪的变化为：可在同年及后续政令中追踪；条文不单独证明地方执行。

证据的边界是：《宋史》为元末官修，本纪保存宋廷纪年与处置；战果、数字、地方执行及对方动机须以编年、会要、对方史料或地方材料互校。 宋金战事、金代地方治理与碑志研究 可用于继续检验这一结论。

\subsection{1134年（南宋紹興四年、金天会12年；公元换算据两国年号对照）}
\noindent\eventanchor{SYD-1134-JIN-016}{丙戌，吳玠敗金兵，復鳳、秦、隴州}\textbf{南宋与金。}
1134年（南宋紹興四年、金天会12年；公元换算据两国年号对照） 的记载将 南宋与金 与〔丙戌，吳玠敗金兵，復鳳、秦、隴州〕相连。此处可确认的是这项被记录的行动；题名保留的城、路、水道或机构名称，是目前可以落地的空间线索。

前一层压力可由以下线索辨认：本纪年条记录政令或机构处置；实施背景须参照制度材料。 这一处置留下的后续条件是：可在同年及后续政令中追踪；条文不单独证明地方执行。

《宋史》为元末官修，本纪保存宋廷纪年与处置；战果、数字、地方执行及对方动机须以编年、会要、对方史料或地方材料互校。 因此，宋金战事、金代地方治理与碑志研究 只能扩展问题，不能代替未保存的细节。

\noindent\eventanchor{SYD-1134-JIN-115}{庚戌，賞承州水砦首領徐康等要擊金兵之功，轉官有差，仍蠲承、楚、泰州水砦民兵賦役十年}\textbf{南宋与金。}
〔庚戌，賞承州水砦首領徐康等要擊金兵之功，轉官有差，仍蠲承、楚、泰州水砦民兵賦役十年〕见于 《宋史》卷027〈紹興四年〉；《金史》本纪同年条（互校） 的 1134年（南宋紹興四年、金天会12年；公元换算据两国年号对照） 记录。南宋与金 的选择因此有了可回查的时间位置；题名保留的城、路、水道或机构名称，是目前可以落地的空间线索。

本纪所记财用、赈济或征发属于特定报告场景。 记录所见的结果则是：可与食货志、文书和物质材料比较，不外推为全境总量。 日期相邻不等于因果已经证实。

关于可说范围，须保留：《宋史》为元末官修，本纪保存宋廷纪年与处置；战果、数字、地方执行及对方动机须以编年、会要、对方史料或地方材料互校。 进一步讨论可参照 宋金战事、金代地方治理与碑志研究

\noindent\eventanchor{SYD-1134-JIN-169}{癸亥，劉光世遣統制王德擊金人于滁州之桑根，敗之}\textbf{南宋与金。}
1134年（南宋紹興四年、金天会12年；公元换算据两国年号对照），《宋史》卷027〈紹興四年〉；《金史》本纪同年条（互校） 所见的〔癸亥，劉光世遣統制王德擊金人于滁州之桑根，敗之〕把 南宋与金 的处置留在可复核的年序中；题名保留的城、路、水道或机构名称，是目前可以落地的空间线索。

把本项放回事件链，能够看到的前提为：本纪从朝廷视角记录军政行动；出兵与战果需分辨。 而它所打开或收紧的下一步是：后续路线、控制与损失应与对方记录和遗址材料复核。

证据的边界是：《宋史》为元末官修，本纪保存宋廷纪年与处置；战果、数字、地方执行及对方动机须以编年、会要、对方史料或地方材料互校。 宋金战事、金代地方治理与碑志研究 可用于继续检验这一结论。

\noindent\eventanchor{SYD-1134-JIN-206}{己丑，金人攻承州，韓世忠遣將成閔、解元合兵擊於北門，敗之}\textbf{南宋与金。}
南宋与金 在 1134年（南宋紹興四年、金天会12年；公元换算据两国年号对照） 的材料痕迹是〔己丑，金人攻承州，韓世忠遣將成閔、解元合兵擊於北門，敗之〕，出处为 《宋史》卷027〈紹興四年〉；《金史》本纪同年条（互校）。题名保留的城、路、水道或机构名称，是目前可以落地的空间线索。

它承接的条件是：本纪从朝廷视角记录军政行动；出兵与战果需分辨。 随后可追踪的变化为：后续路线、控制与损失应与对方记录和遗址材料复核。

《宋史》为元末官修，本纪保存宋廷纪年与处置；战果、数字、地方执行及对方动机须以编年、会要、对方史料或地方材料互校。 因此，宋金战事、金代地方治理与碑志研究 只能扩展问题，不能代替未保存的细节。

\noindent\eventanchor{SYD-1134-JIN-239}{己巳，劉光世遣統制王師晟等率兵夜入南壽春府襲金人，敗之，執偽齊知府王靖}\textbf{南宋与金。}
从 《宋史》卷027〈紹興四年〉；《金史》本纪同年条（互校） 的 1134年（南宋紹興四年、金天会12年；公元换算据两国年号对照） 条目看，〔己巳，劉光世遣統制王師晟等率兵夜入南壽春府襲金人，敗之，執偽齊知府王靖〕使 南宋与金 的一个具体行动进入叙事；题名保留的城、路、水道或机构名称，是目前可以落地的空间线索。

前一层压力可由以下线索辨认：本纪从朝廷视角记录军政行动；出兵与战果需分辨。 这一处置留下的后续条件是：后续路线、控制与损失应与对方记录和遗址材料复核。

关于可说范围，须保留：《宋史》为元末官修，本纪保存宋廷纪年与处置；战果、数字、地方执行及对方动机须以编年、会要、对方史料或地方材料互校。 进一步讨论可参照 宋金战事、金代地方治理与碑志研究

\subsection{1135年（南宋紹興五年、金天会13年；公元换算据两国年号对照）}
\noindent\eventanchor{SYD-1135-JIN-017}{丁未，戒諸軍戰陳毋殺中原民籍充金兵者}\textbf{南宋与金。}
1135年（南宋紹興五年、金天会13年；公元换算据两国年号对照） 的记载将 南宋与金 与〔丁未，戒諸軍戰陳毋殺中原民籍充金兵者〕相连。此处可确认的是这项被记录的行动；材料未将地点细化到城址，不能用中枢所在地替它补足。

本纪从朝廷视角记录军政行动；出兵与战果需分辨。 记录所见的结果则是：后续路线、控制与损失应与对方记录和遗址材料复核。 日期相邻不等于因果已经证实。

证据的边界是：《宋史》为元末官修，本纪保存宋廷纪年与处置；战果、数字、地方执行及对方动机须以编年、会要、对方史料或地方材料互校。 宋金战事、金代地方治理与碑志研究 可用于继续检验这一结论。

\noindent\eventanchor{SYD-1135-JIN-116}{庚戌，張俊遣統領楊忠閔、王進夾擊金人於淮南岸，敗之，降其將程師回、張延壽}\textbf{南宋与金。}
〔庚戌，張俊遣統領楊忠閔、王進夾擊金人於淮南岸，敗之，降其將程師回、張延壽〕见于 《宋史》卷028〈紹興五年〉；《金史》本纪同年条（互校） 的 1135年（南宋紹興五年、金天会13年；公元换算据两国年号对照） 记录。南宋与金 的选择因此有了可回查的时间位置；题名保留的城、路、水道或机构名称，是目前可以落地的空间线索。

把本项放回事件链，能够看到的前提为：本纪从朝廷视角记录军政行动；出兵与战果需分辨。 而它所打开或收紧的下一步是：后续路线、控制与损失应与对方记录和遗址材料复核。

《宋史》为元末官修，本纪保存宋廷纪年与处置；战果、数字、地方执行及对方动机须以编年、会要、对方史料或地方材料互校。 因此，宋金战事、金代地方治理与碑志研究 只能扩展问题，不能代替未保存的细节。

\noindent\eventanchor{SYD-1135-JIN-170}{遣何蘚等奉使金國，通問二帝}\textbf{南宋与金。}
1135年（南宋紹興五年、金天会13年；公元换算据两国年号对照），《宋史》卷028〈紹興五年〉；《金史》本纪同年条（互校） 所见的〔遣何蘚等奉使金國，通問二帝〕把 南宋与金 的处置留在可复核的年序中；材料未将地点细化到城址，不能用中枢所在地替它补足。

它承接的条件是：本纪记录使节、贡赐或往来，不等于稳定统属关系。 随后可追踪的变化为：可与对方编年和交通、文书材料核对其后续落实。

关于可说范围，须保留：《宋史》为元末官修，本纪保存宋廷纪年与处置；战果、数字、地方执行及对方动机须以编年、会要、对方史料或地方材料互校。 进一步讨论可参照 宋金战事、金代地方治理与碑志研究

\noindent\eventanchor{SYD-1135-JIN-207}{是月，金主晟殂，旻之孫亶立}\textbf{南宋与金。}
南宋与金 在 1135年（南宋紹興五年、金天会13年；公元换算据两国年号对照） 的材料痕迹是〔是月，金主晟殂，旻之孫亶立〕，出处为 《宋史》卷028〈紹興五年〉；《金史》本纪同年条（互校）。材料未将地点细化到城址，不能用中枢所在地替它补足。

前一层压力可由以下线索辨认：本纪年条提供日期和中枢可见行动；前因不据单条补定。 这一处置留下的后续条件是：可回查同年及后续条目，地方影响仍须另证。

证据的边界是：《宋史》为元末官修，本纪保存宋廷纪年与处置；战果、数字、地方执行及对方动机须以编年、会要、对方史料或地方材料互校。 宋金战事、金代地方治理与碑志研究 可用于继续检验这一结论。

\noindent\eventanchor{SYD-1135-JIN-240}{辛亥，淮東統制崔德明襲敗金兵於盱眙}\textbf{南宋与金。}
从 《宋史》卷028〈紹興五年〉；《金史》本纪同年条（互校） 的 1135年（南宋紹興五年、金天会13年；公元换算据两国年号对照） 条目看，〔辛亥，淮東統制崔德明襲敗金兵於盱眙〕使 南宋与金 的一个具体行动进入叙事；题名保留的城、路、水道或机构名称，是目前可以落地的空间线索。

本纪从朝廷视角记录军政行动；出兵与战果需分辨。 记录所见的结果则是：后续路线、控制与损失应与对方记录和遗址材料复核。 日期相邻不等于因果已经证实。

《宋史》为元末官修，本纪保存宋廷纪年与处置；战果、数字、地方执行及对方动机须以编年、会要、对方史料或地方材料互校。 因此，宋金战事、金代地方治理与碑志研究 只能扩展问题，不能代替未保存的细节。

\subsection{1136年（南宋紹興六年、金天会14年；公元换算据两国年号对照）}
\noindent\eventanchor{SYD-1136-JIN-018}{是月，劉豫聞親征，告急于金主亶求援，亶不許，豫自起兵三十萬，命子麟趣合肥，侄猊出渦口，引兵分道入寇}\textbf{南宋与金。}
1136年（南宋紹興六年、金天会14年；公元换算据两国年号对照） 的记载将 南宋与金 与〔是月，劉豫聞親征，告急于金主亶求援，亶不許，豫自起兵三十萬，命子麟趣合肥，侄猊出渦口，引兵分道入寇〕相连。此处可确认的是这项被记录的行动；材料未将地点细化到城址，不能用中枢所在地替它补足。

把本项放回事件链，能够看到的前提为：本纪年条记录政令或机构处置；实施背景须参照制度材料。 而它所打开或收紧的下一步是：可在同年及后续政令中追踪；条文不单独证明地方执行。

关于可说范围，须保留：《宋史》为元末官修，本纪保存宋廷纪年与处置；战果、数字、地方执行及对方动机须以编年、会要、对方史料或地方材料互校。 进一步讨论可参照 宋金战事、金代地方治理与碑志研究

\noindent\eventanchor{SYD-1136-JIN-117}{戊戌，韓世忠攻淮陽軍，及金人戰，敗之}\textbf{南宋与金。}
〔戊戌，韓世忠攻淮陽軍，及金人戰，敗之〕见于 《宋史》卷028〈紹興六年〉；《金史》本纪同年条（互校） 的 1136年（南宋紹興六年、金天会14年；公元换算据两国年号对照） 记录。南宋与金 的选择因此有了可回查的时间位置；题名保留的城、路、水道或机构名称，是目前可以落地的空间线索。

它承接的条件是：本纪从朝廷视角记录军政行动；出兵与战果需分辨。 随后可追踪的变化为：后续路线、控制与损失应与对方记录和遗址材料复核。

证据的边界是：《宋史》为元末官修，本纪保存宋廷纪年与处置；战果、数字、地方执行及对方动机须以编年、会要、对方史料或地方材料互校。 宋金战事、金代地方治理与碑志研究 可用于继续检验这一结论。

\noindent\eventanchor{SYD-1136-JIN-171}{乙卯，韓世忠引兵攻宿遷縣，統制呼延通與金兵戰，敗之，禽其將孛堇牙合}\textbf{南宋与金。}
1136年（南宋紹興六年、金天会14年；公元换算据两国年号对照），《宋史》卷028〈紹興六年〉；《金史》本纪同年条（互校） 所见的〔乙卯，韓世忠引兵攻宿遷縣，統制呼延通與金兵戰，敗之，禽其將孛堇牙合〕把 南宋与金 的处置留在可复核的年序中；材料未将地点细化到城址，不能用中枢所在地替它补足。

前一层压力可由以下线索辨认：本纪从朝廷视角记录军政行动；出兵与战果需分辨。 这一处置留下的后续条件是：后续路线、控制与损失应与对方记录和遗址材料复核。

《宋史》为元末官修，本纪保存宋廷纪年与处置；战果、数字、地方执行及对方动机须以编年、会要、对方史料或地方材料互校。 因此，宋金战事、金代地方治理与碑志研究 只能扩展问题，不能代替未保存的细节。

\subsection{1137年}
\noindent\eventanchor{JIN-1137-EVENT-090}{金废除刘豫政权，改为直接经营中原}\textbf{金与南宋。}
金与南宋 在 1137年 的材料痕迹是〔金废除刘豫政权，改为直接经营中原〕，出处为 《金史》相关本纪；《宋史》或《三朝北盟会编》《建炎以来系年要录》对应年条；金代碑志与蒙古、高丽材料（据事核）。材料未将地点细化到城址，不能用中枢所在地替它补足。

“金废除刘豫政权，改为直接经营中原”发生在金在华北经营与宋金南北对峙的具体压力与机会之中，相关行动者的选择不能脱离此前事件链解释。 记录所见的结果则是：“金废除刘豫政权，改为直接经营中原”为后续政治、边防或资源配置留下条件；其社会影响与实际覆盖范围仍须以异类材料限定。 日期相邻不等于因果已经证实。

关于可说范围，须保留：《金史》为元修，宋金战报和官修合法性语言各有宣传性；军数、俘获、控制范围和地方执行须以对方及物质材料限缩。 进一步讨论可参照 金代国家形成、宋金边境、都城与金蒙转换研究

\subsection{1137年（南宋紹興七年、金天会15年；公元换算据两国年号对照）}
\noindent\eventanchor{SYD-1137-JIN-019}{庚辰，韓世忠引兵渡淮，逆擊金人于劉冷莊，敗之}\textbf{南宋与金。}
从 《宋史》卷028〈紹興七年〉；《金史》本纪同年条（互校） 的 1137年（南宋紹興七年、金天会15年；公元换算据两国年号对照） 条目看，〔庚辰，韓世忠引兵渡淮，逆擊金人于劉冷莊，敗之〕使 南宋与金 的一个具体行动进入叙事；题名保留的城、路、水道或机构名称，是目前可以落地的空间线索。

把本项放回事件链，能够看到的前提为：本纪从朝廷视角记录军政行动；出兵与战果需分辨。 而它所打开或收紧的下一步是：后续路线、控制与损失应与对方记录和遗址材料复核。

证据的边界是：《宋史》为元末官修，本纪保存宋廷纪年与处置；战果、数字、地方执行及对方动机须以编年、会要、对方史料或地方材料互校。 宋金战事、金代地方治理与碑志研究 可用于继续检验这一结论。

\noindent\eventanchor{SYD-1137-JIN-118}{庚子，遣王倫等使金國迎奉梓宮}\textbf{南宋与金。}
1137年（南宋紹興七年、金天会15年；公元换算据两国年号对照） 的记载将 南宋与金 与〔庚子，遣王倫等使金國迎奉梓宮〕相连。此处可确认的是这项被记录的行动；题名保留的城、路、水道或机构名称，是目前可以落地的空间线索。

它承接的条件是：本纪记录使节、贡赐或往来，不等于稳定统属关系。 随后可追踪的变化为：可与对方编年和交通、文书材料核对其后续落实。

《宋史》为元末官修，本纪保存宋廷纪年与处置；战果、数字、地方执行及对方动机须以编年、会要、对方史料或地方材料互校。 因此，宋金战事、金代地方治理与碑志研究 只能扩展问题，不能代替未保存的细节。

\noindent\eventanchor{SYD-1137-JIN-172}{癸未，王倫等使還，入見，言金國許還梓宮及皇太后，又許還河南諸州}\textbf{南宋与金。}
〔癸未，王倫等使還，入見，言金國許還梓宮及皇太后，又許還河南諸州〕见于 《宋史》卷028〈紹興七年〉；《金史》本纪同年条（互校） 的 1137年（南宋紹興七年、金天会15年；公元换算据两国年号对照） 记录。南宋与金 的选择因此有了可回查的时间位置；题名保留的城、路、水道或机构名称，是目前可以落地的空间线索。

前一层压力可由以下线索辨认：本纪保留灾害或工程的报告地点与朝廷处置；灾情范围需另证。 这一处置留下的后续条件是：可与地方文书、河工和环境材料比较，不将单点记录外推为全域因果。

关于可说范围，须保留：《宋史》为元末官修，本纪保存宋廷纪年与处置；战果、数字、地方执行及对方动机须以编年、会要、对方史料或地方材料互校。 进一步讨论可参照 宋金战事、金代地方治理与碑志研究

\noindent\eventanchor{SYD-1137-JIN-208}{壬子，統制呼延通、王權等襲擊金人于淮陽軍，敗之，丁巳，詔六參日，輪行在百官一員轉對}\textbf{南宋与金。}
1137年（南宋紹興七年、金天会15年；公元换算据两国年号对照），《宋史》卷028〈紹興七年〉；《金史》本纪同年条（互校） 所见的〔壬子，統制呼延通、王權等襲擊金人于淮陽軍，敗之，丁巳，詔六參日，輪行在百官一員轉對〕把 南宋与金 的处置留在可复核的年序中；题名保留的城、路、水道或机构名称，是目前可以落地的空间线索。

本纪年条记录政令或机构处置；实施背景须参照制度材料。 记录所见的结果则是：可在同年及后续政令中追踪；条文不单独证明地方执行。 日期相邻不等于因果已经证实。

证据的边界是：《宋史》为元末官修，本纪保存宋廷纪年与处置；战果、数字、地方执行及对方动机须以编年、会要、对方史料或地方材料互校。 宋金战事、金代地方治理与碑志研究 可用于继续检验这一结论。

\subsection{1138年（南宋紹興八年、金天眷1年；公元换算据两国年号对照）}
\noindent\eventanchor{SYD-1138-JIN-020}{帝曰：「朕嗣守祖宗基業，豈受金人封冊}\textbf{南宋与金。}
南宋与金 在 1138年（南宋紹興八年、金天眷1年；公元换算据两国年号对照） 的材料痕迹是〔帝曰：「朕嗣守祖宗基業，豈受金人封冊〕，出处为 《宋史》卷029〈紹興八年〉；《金史》本纪同年条（互校）。材料未将地点细化到城址，不能用中枢所在地替它补足。

把本项放回事件链，能够看到的前提为：本纪年条提供日期和中枢可见行动；前因不据单条补定。 而它所打开或收紧的下一步是：可回查同年及后续条目，地方影响仍须另证。

《宋史》为元末官修，本纪保存宋廷纪年与处置；战果、数字、地方执行及对方动机须以编年、会要、对方史料或地方材料互校。 因此，宋金战事、金代地方治理与碑志研究 只能扩展问题，不能代替未保存的细节。

\noindent\eventanchor{SYD-1138-JIN-119}{丁-{丑}-，金國使張通古、蕭哲與王倫皆來}\textbf{南宋与金。}
从 《宋史》卷029〈紹興八年〉；《金史》本纪同年条（互校） 的 1138年（南宋紹興八年、金天眷1年；公元换算据两国年号对照） 条目看，〔丁-{丑}-，金國使張通古、蕭哲與王倫皆來〕使 南宋与金 的一个具体行动进入叙事；材料未将地点细化到城址，不能用中枢所在地替它补足。

它承接的条件是：本纪年条提供日期和中枢可见行动；前因不据单条补定。 随后可追踪的变化为：可回查同年及后续条目，地方影响仍须另证。

关于可说范围，须保留：《宋史》为元末官修，本纪保存宋廷纪年与处置；战果、数字、地方执行及对方动机须以编年、会要、对方史料或地方材料互校。 进一步讨论可参照 宋金战事、金代地方治理与碑志研究

\noindent\eventanchor{SYD-1138-JIN-173}{丁-{丑}-，詔：「金國使來，盡割河南、陝西故地，通好于我，許還梓宮及母兄親族，余無需索}\textbf{南宋与金。}
1138年（南宋紹興八年、金天眷1年；公元换算据两国年号对照） 的记载将 南宋与金 与〔丁-{丑}-，詔：「金國使來，盡割河南、陝西故地，通好于我，許還梓宮及母兄親族，余無需索〕相连。此处可确认的是这项被记录的行动；题名保留的城、路、水道或机构名称，是目前可以落地的空间线索。

前一层压力可由以下线索辨认：本纪年条记录政令或机构处置；实施背景须参照制度材料。 这一处置留下的后续条件是：可在同年及后续政令中追踪；条文不单独证明地方执行。

证据的边界是：《宋史》为元末官修，本纪保存宋廷纪年与处置；战果、数字、地方执行及对方动机须以编年、会要、对方史料或地方材料互校。 宋金战事、金代地方治理与碑志研究 可用于继续检验这一结论。

\noindent\eventanchor{SYD-1138-JIN-209}{丁未，金國使烏陵思謀、石慶充與王倫等偕來}\textbf{南宋与金。}
〔丁未，金國使烏陵思謀、石慶充與王倫等偕來〕见于 《宋史》卷029〈紹興八年〉；《金史》本纪同年条（互校） 的 1138年（南宋紹興八年、金天眷1年；公元换算据两国年号对照） 记录。南宋与金 的选择因此有了可回查的时间位置；材料未将地点细化到城址，不能用中枢所在地替它补足。

本纪年条提供日期和中枢可见行动；前因不据单条补定。 记录所见的结果则是：可回查同年及后续条目，地方影响仍须另证。 日期相邻不等于因果已经证实。

《宋史》为元末官修，本纪保存宋廷纪年与处置；战果、数字、地方执行及对方动机须以编年、会要、对方史料或地方材料互校。 因此，宋金战事、金代地方治理与碑志研究 只能扩展问题，不能代替未保存的细节。

\noindent\eventanchor{SYD-1138-JIN-241}{辛-{丑}-，詔：「金國遣使入境，欲朕屈己就和，命侍從、台諫詳思條奏}\textbf{南宋与金。}
1138年（南宋紹興八年、金天眷1年；公元换算据两国年号对照），《宋史》卷029〈紹興八年〉；《金史》本纪同年条（互校） 所见的〔辛-{丑}-，詔：「金國遣使入境，欲朕屈己就和，命侍從、台諫詳思條奏〕把 南宋与金 的处置留在可复核的年序中；材料未将地点细化到城址，不能用中枢所在地替它补足。

把本项放回事件链，能够看到的前提为：本纪年条记录政令或机构处置；实施背景须参照制度材料。 而它所打开或收紧的下一步是：可在同年及后续政令中追踪；条文不单独证明地方执行。

关于可说范围，须保留：《宋史》为元末官修，本纪保存宋廷纪年与处置；战果、数字、地方执行及对方动机须以编年、会要、对方史料或地方材料互校。 进一步讨论可参照 宋金战事、金代地方治理与碑志研究

\subsection{1139年（南宋紹興九年、金天眷2年；公元换算据两国年号对照）}
\noindent\eventanchor{SYD-1139-JIN-021}{丙辰，金國以撻懶主和割地，疑其二心，殺之}\textbf{南宋与金。}
南宋与金 在 1139年（南宋紹興九年、金天眷2年；公元换算据两国年号对照） 的材料痕迹是〔丙辰，金國以撻懶主和割地，疑其二心，殺之〕，出处为 《宋史》卷029〈紹興九年〉；《金史》本纪同年条（互校）。材料未将地点细化到城址，不能用中枢所在地替它补足。

它承接的条件是：本纪从朝廷视角记录军政行动；出兵与战果需分辨。 随后可追踪的变化为：后续路线、控制与损失应与对方记录和遗址材料复核。

证据的边界是：《宋史》为元末官修，本纪保存宋廷纪年与处置；战果、数字、地方执行及对方动机须以编年、会要、对方史料或地方材料互校。 宋金战事、金代地方治理与碑志研究 可用于继续检验这一结论。

\noindent\eventanchor{SYD-1139-JIN-120}{丙申，金主詔諭河南諸州以割地歸我之意}\textbf{南宋与金。}
从 《宋史》卷029〈紹興九年〉；《金史》本纪同年条（互校） 的 1139年（南宋紹興九年、金天眷2年；公元换算据两国年号对照） 条目看，〔丙申，金主詔諭河南諸州以割地歸我之意〕使 南宋与金 的一个具体行动进入叙事；题名保留的城、路、水道或机构名称，是目前可以落地的空间线索。

前一层压力可由以下线索辨认：本纪年条记录政令或机构处置；实施背景须参照制度材料。 这一处置留下的后续条件是：可在同年及后续政令中追踪；条文不单独证明地方执行。

《宋史》为元末官修，本纪保存宋廷纪年与处置；战果、数字、地方执行及对方动机须以编年、会要、对方史料或地方材料互校。 因此，宋金战事、金代地方治理与碑志研究 只能扩展问题，不能代替未保存的细节。

\noindent\eventanchor{SYD-1139-JIN-174}{王倫見金主于-{御}-林子，被拘于河間，遣其副藍公佐先歸}\textbf{南宋与金。}
1139年（南宋紹興九年、金天眷2年；公元换算据两国年号对照） 的记载将 南宋与金 与〔王倫見金主于-{御}-林子，被拘于河間，遣其副藍公佐先歸〕相连。此处可确认的是这项被记录的行动；题名保留的城、路、水道或机构名称，是目前可以落地的空间线索。

本纪保留灾害或工程的报告地点与朝廷处置；灾情范围需另证。 记录所见的结果则是：可与地方文书、河工和环境材料比较，不将单点记录外推为全域因果。 日期相邻不等于因果已经证实。

关于可说范围，须保留：《宋史》为元末官修，本纪保存宋廷纪年与处置；战果、数字、地方执行及对方动机须以编年、会要、对方史料或地方材料互校。 进一步讨论可参照 宋金战事、金代地方治理与碑志研究

\noindent\eventanchor{SYD-1139-JIN-210}{乙亥，遣前知宿州趙榮、知壽州王威俱還金國}\textbf{南宋与金。}
〔乙亥，遣前知宿州趙榮、知壽州王威俱還金國〕见于 《宋史》卷029〈紹興九年〉；《金史》本纪同年条（互校） 的 1139年（南宋紹興九年、金天眷2年；公元换算据两国年号对照） 记录。南宋与金 的选择因此有了可回查的时间位置；题名保留的城、路、水道或机构名称，是目前可以落地的空间线索。

把本项放回事件链，能够看到的前提为：本纪年条提供日期和中枢可见行动；前因不据单条补定。 而它所打开或收紧的下一步是：可回查同年及后续条目，地方影响仍须另证。

证据的边界是：《宋史》为元末官修，本纪保存宋廷纪年与处置；战果、数字、地方执行及对方动机须以编年、会要、对方史料或地方材料互校。 宋金战事、金代地方治理与碑志研究 可用于继续检验这一结论。

\subsection{1140年}
\noindent\eventanchor{JIN-1140-EVENT-091}{金宋在河南、淮西和陕西再度大规模作战}\textbf{金与南宋。}
1140年，《金史》相关本纪；《宋史》或《三朝北盟会编》《建炎以来系年要录》对应年条；金代碑志与蒙古、高丽材料（据事核） 所见的〔金宋在河南、淮西和陕西再度大规模作战〕把 金与南宋 的处置留在可复核的年序中；题名保留的城、路、水道或机构名称，是目前可以落地的空间线索。

它承接的条件是：金在华北经营与宋金南北对峙中既有的联盟、交通与补给压力，为“金宋在河南、淮西和陕西再度大规模作战”提供了直接的军事或动员背景。 随后可追踪的变化为：“金宋在河南、淮西和陕西再度大规模作战”改变了相关城镇、道路或政权间的军事选择；伤亡、俘获和控制范围仅可按各方材料分别确认。

《金史》为元修，宋金战报和官修合法性语言各有宣传性；军数、俘获、控制范围和地方执行须以对方及物质材料限缩。 因此，金代国家形成、宋金边境、都城与金蒙转换研究 只能扩展问题，不能代替未保存的细节。

\subsection{1140年（南宋紹興十年、金天眷3年；公元换算据两国年号对照）}
\noindent\eventanchor{SYD-1140-JIN-022}{-{岳}-飛領兵援劉錡，與金人戰于蔡州，敗之，復蔡州}\textbf{南宋与金。}
南宋与金 在 1140年（南宋紹興十年、金天眷3年；公元换算据两国年号对照） 的材料痕迹是〔-{岳}-飛領兵援劉錡，與金人戰于蔡州，敗之，復蔡州〕，出处为 《宋史》卷029〈紹興十年〉；《金史》本纪同年条（互校）。题名保留的城、路、水道或机构名称，是目前可以落地的空间线索。

前一层压力可由以下线索辨认：本纪年条记录政令或机构处置；实施背景须参照制度材料。 这一处置留下的后续条件是：可在同年及后续政令中追踪；条文不单独证明地方执行。

关于可说范围，须保留：《宋史》为元末官修，本纪保存宋廷纪年与处置；战果、数字、地方执行及对方动机须以编年、会要、对方史料或地方材料互校。 进一步讨论可参照 宋金战事、金代地方治理与碑志研究

\noindent\eventanchor{SYD-1140-JIN-121}{-{岳}-飛遣統制郝晸等與金人戰于鄭州北，復鄭州}\textbf{南宋与金。}
从 《宋史》卷029〈紹興十年〉；《金史》本纪同年条（互校） 的 1140年（南宋紹興十年、金天眷3年；公元换算据两国年号对照） 条目看，〔-{岳}-飛遣統制郝晸等與金人戰于鄭州北，復鄭州〕使 南宋与金 的一个具体行动进入叙事；题名保留的城、路、水道或机构名称，是目前可以落地的空间线索。

本纪年条记录政令或机构处置；实施背景须参照制度材料。 记录所见的结果则是：可在同年及后续政令中追踪；条文不单独证明地方执行。 日期相邻不等于因果已经证实。

证据的边界是：《宋史》为元末官修，本纪保存宋廷纪年与处置；战果、数字、地方执行及对方动机须以编年、会要、对方史料或地方材料互校。 宋金战事、金代地方治理与碑志研究 可用于继续检验这一结论。

\noindent\eventanchor{SYD-1140-JIN-175}{丙辰，-{岳}-飛將牛皋及金人戰于京西，敗之}\textbf{南宋与金。}
1140年（南宋紹興十年、金天眷3年；公元换算据两国年号对照） 的记载将 南宋与金 与〔丙辰，-{岳}-飛將牛皋及金人戰于京西，敗之〕相连。此处可确认的是这项被记录的行动；题名保留的城、路、水道或机构名称，是目前可以落地的空间线索。

把本项放回事件链，能够看到的前提为：本纪从朝廷视角记录军政行动；出兵与战果需分辨。 而它所打开或收紧的下一步是：后续路线、控制与损失应与对方记录和遗址材料复核。

《宋史》为元末官修，本纪保存宋廷纪年与处置；战果、数字、地方执行及对方动机须以编年、会要、对方史料或地方材料互校。 因此，宋金战事、金代地方治理与碑志研究 只能扩展问题，不能代替未保存的细节。

\noindent\eventanchor{SYD-1140-JIN-211}{丙戌，金人陷拱州，守臣王慥死之}\textbf{南宋与金。}
〔丙戌，金人陷拱州，守臣王慥死之〕见于 《宋史》卷029〈紹興十年〉；《金史》本纪同年条（互校） 的 1140年（南宋紹興十年、金天眷3年；公元换算据两国年号对照） 记录。南宋与金 的选择因此有了可回查的时间位置；题名保留的城、路、水道或机构名称，是目前可以落地的空间线索。

它承接的条件是：本纪年条提供日期和中枢可见行动；前因不据单条补定。 随后可追踪的变化为：可回查同年及后续条目，地方影响仍须另证。

关于可说范围，须保留：《宋史》为元末官修，本纪保存宋廷纪年与处置；战果、数字、地方执行及对方动机须以编年、会要、对方史料或地方材料互校。 进一步讨论可参照 宋金战事、金代地方治理与碑志研究

\noindent\eventanchor{SYD-1140-JIN-242}{川、陝宣撫副使胡世將屢言金人必渝盟，宜爲備}\textbf{南宋与金。}
1140年（南宋紹興十年、金天眷3年；公元换算据两国年号对照），《宋史》卷029〈紹興十年〉；《金史》本纪同年条（互校） 所见的〔川、陝宣撫副使胡世將屢言金人必渝盟，宜爲備〕把 南宋与金 的处置留在可复核的年序中；题名保留的城、路、水道或机构名称，是目前可以落地的空间线索。

前一层压力可由以下线索辨认：本纪记录使节、贡赐或往来，不等于稳定统属关系。 这一处置留下的后续条件是：可与对方编年和交通、文书材料核对其后续落实。

证据的边界是：《宋史》为元末官修，本纪保存宋廷纪年与处置；战果、数字、地方执行及对方动机须以编年、会要、对方史料或地方材料互校。 宋金战事、金代地方治理与碑志研究 可用于继续检验这一结论。

\subsection{1141年（金皇统元年、宋绍兴十一年；干支、月日待核；公元换算据双方本纪年号对照）}
\noindent\eventanchor{JIN-1141-EVENT-092}{金宋议定绍兴和约框架}\textbf{金与南宋。}
金与南宋 在 1141年（金皇统元年、宋绍兴十一年；干支、月日待核；公元换算据双方本纪年号对照） 的材料痕迹是〔金宋议定绍兴和约框架〕，出处为 《金史》相关本纪；《宋史》或《三朝北盟会编》《建炎以来系年要录》对应年条；金代碑志与蒙古、高丽材料（据事核）。材料未将地点细化到城址，不能用中枢所在地替它补足。

“金宋议定绍兴和约框架”发生在金在华北经营与宋金南北对峙的具体压力与机会之中，相关行动者的选择不能脱离此前事件链解释。 记录所见的结果则是：“金宋议定绍兴和约框架”为后续政治、边防或资源配置留下条件；其社会影响与实际覆盖范围仍须以异类材料限定。 日期相邻不等于因果已经证实。

《金史》为元修，宋金战报和官修合法性语言各有宣传性；军数、俘获、控制范围和地方执行须以对方及物质材料限缩。 因此，金代国家形成、宋金边境、都城与金蒙转换研究 只能扩展问题，不能代替未保存的细节。

\subsection{1141年（南宋紹興十一年、金皇统1年；公元换算据两国年号对照）}
\noindent\eventanchor{SYD-1141-JIN-023}{丙申，遣劉光遠等充金國通問使}\textbf{南宋与金。}
从 《宋史》卷029〈紹興十一年〉；《金史》本纪同年条（互校） 的 1141年（南宋紹興十一年、金皇统1年；公元换算据两国年号对照） 条目看，〔丙申，遣劉光遠等充金國通問使〕使 南宋与金 的一个具体行动进入叙事；材料未将地点细化到城址，不能用中枢所在地替它补足。

把本项放回事件链，能够看到的前提为：本纪记录使节、贡赐或往来，不等于稳定统属关系。 而它所打开或收紧的下一步是：可与对方编年和交通、文书材料核对其后续落实。

关于可说范围，须保留：《宋史》为元末官修，本纪保存宋廷纪年与处置；战果、数字、地方执行及对方动机须以编年、会要、对方史料或地方材料互校。 进一步讨论可参照 宋金战事、金代地方治理与碑志研究

\noindent\eventanchor{SYD-1141-JIN-122}{己卯，統制<u>關師古</u>、<u>李橫</u>擊敗金人于<u>巢縣</u>，復之}\textbf{南宋与金。}
1141年（南宋紹興十一年、金皇统1年；公元换算据两国年号对照） 的记载将 南宋与金 与〔己卯，統制<u>關師古</u>、<u>李橫</u>擊敗金人于<u>巢縣</u>，復之〕相连。此处可确认的是这项被记录的行动；题名保留的城、路、水道或机构名称，是目前可以落地的空间线索。

它承接的条件是：本纪年条记录政令或机构处置；实施背景须参照制度材料。 随后可追踪的变化为：可在同年及后续政令中追踪；条文不单独证明地方执行。

证据的边界是：《宋史》为元末官修，本纪保存宋廷纪年与处置；战果、数字、地方执行及对方动机须以编年、会要、对方史料或地方材料互校。 宋金战事、金代地方治理与碑志研究 可用于继续检验这一结论。

\noindent\eventanchor{SYD-1141-JIN-176}{壬午，遣魏良臣、王公亮爲金國稟議使}\textbf{南宋与金。}
〔壬午，遣魏良臣、王公亮爲金國稟議使〕见于 《宋史》卷029〈紹興十一年〉；《金史》本纪同年条（互校） 的 1141年（南宋紹興十一年、金皇统1年；公元换算据两国年号对照） 记录。南宋与金 的选择因此有了可回查的时间位置；材料未将地点细化到城址，不能用中枢所在地替它补足。

前一层压力可由以下线索辨认：本纪记录使节、贡赐或往来，不等于稳定统属关系。 这一处置留下的后续条件是：可与对方编年和交通、文书材料核对其后续落实。

《宋史》为元末官修，本纪保存宋廷纪年与处置；战果、数字、地方执行及对方动机须以编年、会要、对方史料或地方材料互校。 因此，宋金战事、金代地方治理与碑志研究 只能扩展问题，不能代替未保存的细节。

\noindent\eventanchor{SYD-1141-JIN-212}{壬子，璘率姚仲及金人戰于丁劉圈，敗之}\textbf{南宋与金。}
1141年（南宋紹興十一年、金皇统1年；公元换算据两国年号对照），《宋史》卷029〈紹興十一年〉；《金史》本纪同年条（互校） 所见的〔壬子，璘率姚仲及金人戰于丁劉圈，敗之〕把 南宋与金 的处置留在可复核的年序中；材料未将地点细化到城址，不能用中枢所在地替它补足。

本纪从朝廷视角记录军政行动；出兵与战果需分辨。 记录所见的结果则是：后续路线、控制与损失应与对方记录和遗址材料复核。 日期相邻不等于因果已经证实。

关于可说范围，须保留：《宋史》为元末官修，本纪保存宋廷纪年与处置；战果、数字、地方执行及对方动机须以编年、会要、对方史料或地方材料互校。 进一步讨论可参照 宋金战事、金代地方治理与碑志研究

\noindent\eventanchor{SYD-1141-JIN-243}{是月，金人陷濠州，邵隆復陝州}\textbf{南宋与金。}
南宋与金 在 1141年（南宋紹興十一年、金皇统1年；公元换算据两国年号对照） 的材料痕迹是〔是月，金人陷濠州，邵隆復陝州〕，出处为 《宋史》卷029〈紹興十一年〉；《金史》本纪同年条（互校）。题名保留的城、路、水道或机构名称，是目前可以落地的空间线索。

把本项放回事件链，能够看到的前提为：本纪年条记录政令或机构处置；实施背景须参照制度材料。 而它所打开或收紧的下一步是：可在同年及后续政令中追踪；条文不单独证明地方执行。

证据的边界是：《宋史》为元末官修，本纪保存宋廷纪年与处置；战果、数字、地方执行及对方动机须以编年、会要、对方史料或地方材料互校。 宋金战事、金代地方治理与碑志研究 可用于继续检验这一结论。

\subsection{1142年（南宋紹興十二年、金皇统2年；公元换算据两国年号对照）}
\noindent\eventanchor{SYD-1142-JIN-024}{癸巳，金主許歸梓宮及皇太后，遣何鑄等還}\textbf{南宋与金。}
从 《宋史》卷030〈紹興十二年〉；《金史》本纪同年条（互校） 的 1142年（南宋紹興十二年、金皇统2年；公元换算据两国年号对照） 条目看，〔癸巳，金主許歸梓宮及皇太后，遣何鑄等還〕使 南宋与金 的一个具体行动进入叙事；题名保留的城、路、水道或机构名称，是目前可以落地的空间线索。

它承接的条件是：本纪年条提供日期和中枢可见行动；前因不据单条补定。 随后可追踪的变化为：可回查同年及后续条目，地方影响仍须另证。

《宋史》为元末官修，本纪保存宋廷纪年与处置；战果、数字、地方执行及对方动机须以编年、会要、对方史料或地方材料互校。 因此，宋金战事、金代地方治理与碑志研究 只能扩展问题，不能代替未保存的细节。

\noindent\eventanchor{SYD-1142-JIN-123}{是月，鄭剛中分畫陝西地界，割商、秦之半畀金國，存上津、豐陽、天水三縣及隴西成紀餘地，棄和尚、方山二原，以大散關為界}\textbf{南宋与金。}
1142年（南宋紹興十二年、金皇统2年；公元换算据两国年号对照） 的记载将 南宋与金 与〔是月，鄭剛中分畫陝西地界，割商、秦之半畀金國，存上津、豐陽、天水三縣及隴西成紀餘地，棄和尚、方山二原，以大散關為界〕相连。此处可确认的是这项被记录的行动；题名保留的城、路、水道或机构名称，是目前可以落地的空间线索。

前一层压力可由以下线索辨认：本纪保留灾害或工程的报告地点与朝廷处置；灾情范围需另证。 这一处置留下的后续条件是：可与地方文书、河工和环境材料比较，不将单点记录外推为全域因果。

关于可说范围，须保留：《宋史》为元末官修，本纪保存宋廷纪年与处置；战果、数字、地方执行及对方动机须以编年、会要、对方史料或地方材料互校。 进一步讨论可参照 宋金战事、金代地方治理与碑志研究

\noindent\eventanchor{SYD-1142-JIN-177}{辛亥，遣張中孚、中彥還金國}\textbf{南宋与金。}
〔辛亥，遣張中孚、中彥還金國〕见于 《宋史》卷030〈紹興十二年〉；《金史》本纪同年条（互校） 的 1142年（南宋紹興十二年、金皇统2年；公元换算据两国年号对照） 记录。南宋与金 的选择因此有了可回查的时间位置；材料未将地点细化到城址，不能用中枢所在地替它补足。

本纪年条提供日期和中枢可见行动；前因不据单条补定。 记录所见的结果则是：可回查同年及后续条目，地方影响仍须另证。 日期相邻不等于因果已经证实。

证据的边界是：《宋史》为元末官修，本纪保存宋廷纪年与处置；战果、数字、地方执行及对方动机须以编年、会要、对方史料或地方材料互校。 宋金战事、金代地方治理与碑志研究 可用于继续检验这一结论。

\noindent\eventanchor{SYD-1142-JIN-213}{乙未，遣沈昭遠等賀金主生辰}\textbf{南宋与金。}
1142年（南宋紹興十二年、金皇统2年；公元换算据两国年号对照），《宋史》卷030〈紹興十二年〉；《金史》本纪同年条（互校） 所见的〔乙未，遣沈昭遠等賀金主生辰〕把 南宋与金 的处置留在可复核的年序中；材料未将地点细化到城址，不能用中枢所在地替它补足。

把本项放回事件链，能够看到的前提为：本纪年条提供日期和中枢可见行动；前因不据单条补定。 而它所打开或收紧的下一步是：可回查同年及后续条目，地方影响仍须另证。

《宋史》为元末官修，本纪保存宋廷纪年与处置；战果、数字、地方执行及对方动机须以编年、会要、对方史料或地方材料互校。 因此，宋金战事、金代地方治理与碑志研究 只能扩展问题，不能代替未保存的细节。

\subsection{1143年（南宋紹興十三年、金皇统3年；公元换算据两国年号对照）}
\noindent\eventanchor{SYD-1143-JIN-025}{己亥，遣鄭朴等使金賀正旦，王師心等賀金主生辰}\textbf{南宋与金。}
南宋与金 在 1143年（南宋紹興十三年、金皇统3年；公元换算据两国年号对照） 的材料痕迹是〔己亥，遣鄭朴等使金賀正旦，王師心等賀金主生辰〕，出处为 《宋史》卷030〈紹興十三年〉；《金史》本纪同年条（互校）。材料未将地点细化到城址，不能用中枢所在地替它补足。

它承接的条件是：本纪记录使节、贡赐或往来，不等于稳定统属关系。 随后可追踪的变化为：可与对方编年和交通、文书材料核对其后续落实。

关于可说范围，须保留：《宋史》为元末官修，本纪保存宋廷纪年与处置；战果、数字、地方执行及对方动机须以编年、会要、对方史料或地方材料互校。 进一步讨论可参照 宋金战事、金代地方治理与碑志研究

\noindent\eventanchor{SYD-1143-JIN-124}{戊戌，洪皓至自金國，入見}\textbf{南宋与金。}
从 《宋史》卷030〈紹興十三年〉；《金史》本纪同年条（互校） 的 1143年（南宋紹興十三年、金皇统3年；公元换算据两国年号对照） 条目看，〔戊戌，洪皓至自金國，入見〕使 南宋与金 的一个具体行动进入叙事；材料未将地点细化到城址，不能用中枢所在地替它补足。

前一层压力可由以下线索辨认：本纪年条提供日期和中枢可见行动；前因不据单条补定。 这一处置留下的后续条件是：可回查同年及后续条目，地方影响仍须另证。

证据的边界是：《宋史》为元末官修，本纪保存宋廷纪年与处置；战果、数字、地方执行及对方动机须以编年、会要、对方史料或地方材料互校。 宋金战事、金代地方治理与碑志研究 可用于继续检验这一结论。

\subsection{1144年（南宋紹興十四年、金皇统4年；公元换算据两国年号对照）}
\noindent\eventanchor{SYD-1144-JIN-026}{甲午，金人來求淮北人之在南者，詔原者聽還}\textbf{南宋与金。}
1144年（南宋紹興十四年、金皇统4年；公元换算据两国年号对照） 的记载将 南宋与金 与〔甲午，金人來求淮北人之在南者，詔原者聽還〕相连。此处可确认的是这项被记录的行动；题名保留的城、路、水道或机构名称，是目前可以落地的空间线索。

本纪年条记录政令或机构处置；实施背景须参照制度材料。 记录所见的结果则是：可在同年及后续政令中追踪；条文不单独证明地方执行。 日期相邻不等于因果已经证实。

《宋史》为元末官修，本纪保存宋廷纪年与处置；战果、数字、地方执行及对方动机须以编年、会要、对方史料或地方材料互校。 因此，宋金战事、金代地方治理与碑志研究 只能扩展问题，不能代替未保存的细节。

\noindent\eventanchor{SYD-1144-JIN-125}{秋七月戊午，金人殺王倫於河間府}\textbf{南宋与金。}
〔秋七月戊午，金人殺王倫於河間府〕见于 《宋史》卷030〈紹興十四年〉；《金史》本纪同年条（互校） 的 1144年（南宋紹興十四年、金皇统4年；公元换算据两国年号对照） 记录。南宋与金 的选择因此有了可回查的时间位置；题名保留的城、路、水道或机构名称，是目前可以落地的空间线索。

把本项放回事件链，能够看到的前提为：本纪保留灾害或工程的报告地点与朝廷处置；灾情范围需另证。 而它所打开或收紧的下一步是：可与地方文书、河工和环境材料比较，不将单点记录外推为全域因果。

关于可说范围，须保留：《宋史》为元末官修，本纪保存宋廷纪年与处置；战果、数字、地方执行及对方动机须以编年、会要、对方史料或地方材料互校。 进一步讨论可参照 宋金战事、金代地方治理与碑志研究

\noindent\eventanchor{SYD-1144-JIN-178}{十四年春正月丁巳，遣羅汝楫等報謝金國}\textbf{南宋与金。}
1144年（南宋紹興十四年、金皇统4年；公元换算据两国年号对照），《宋史》卷030〈紹興十四年〉；《金史》本纪同年条（互校） 所见的〔十四年春正月丁巳，遣羅汝楫等報謝金國〕把 南宋与金 的处置留在可复核的年序中；材料未将地点细化到城址，不能用中枢所在地替它补足。

它承接的条件是：本纪年条提供日期和中枢可见行动；前因不据单条补定。 随后可追踪的变化为：可回查同年及后续条目，地方影响仍须另证。

证据的边界是：《宋史》为元末官修，本纪保存宋廷纪年与处置；战果、数字、地方执行及对方动机须以编年、会要、对方史料或地方材料互校。 宋金战事、金代地方治理与碑志研究 可用于继续检验这一结论。

\noindent\eventanchor{SYD-1144-JIN-214}{乙未，遣林保使金賀正旦，宋之才賀金主生辰}\textbf{南宋与金。}
南宋与金 在 1144年（南宋紹興十四年、金皇统4年；公元换算据两国年号对照） 的材料痕迹是〔乙未，遣林保使金賀正旦，宋之才賀金主生辰〕，出处为 《宋史》卷030〈紹興十四年〉；《金史》本纪同年条（互校）。材料未将地点细化到城址，不能用中枢所在地替它补足。

前一层压力可由以下线索辨认：本纪记录使节、贡赐或往来，不等于稳定统属关系。 这一处置留下的后续条件是：可与对方编年和交通、文书材料核对其后续落实。

《宋史》为元末官修，本纪保存宋廷纪年与处置；战果、数字、地方执行及对方动机须以编年、会要、对方史料或地方材料互校。 因此，宋金战事、金代地方治理与碑志研究 只能扩展问题，不能代替未保存的细节。

\subsection{1145年（南宋紹興十五年、金皇统5年；公元换算据两国年号对照）}
\noindent\eventanchor{SYD-1145-JIN-027}{九月辛酉，遣錢周材使金賀正旦，嚴抑賀金主生辰}\textbf{南宋与金。}
从 《宋史》卷030〈紹興十五年〉；《金史》本纪同年条（互校） 的 1145年（南宋紹興十五年、金皇统5年；公元换算据两国年号对照） 条目看，〔九月辛酉，遣錢周材使金賀正旦，嚴抑賀金主生辰〕使 南宋与金 的一个具体行动进入叙事；材料未将地点细化到城址，不能用中枢所在地替它补足。

本纪所记财用、赈济或征发属于特定报告场景。 记录所见的结果则是：可与食货志、文书和物质材料比较，不外推为全境总量。 日期相邻不等于因果已经证实。

关于可说范围，须保留：《宋史》为元末官修，本纪保存宋廷纪年与处置；战果、数字、地方执行及对方动机须以编年、会要、对方史料或地方材料互校。 进一步讨论可参照 宋金战事、金代地方治理与碑志研究

\noindent\eventanchor{SYD-1145-JIN-126}{三月甲子，遣敷文閣待制周襟、馬觀國、史願、諸將程師回、馬欽、白常皆還金國}\textbf{南宋与金。}
1145年（南宋紹興十五年、金皇统5年；公元换算据两国年号对照） 的记载将 南宋与金 与〔三月甲子，遣敷文閣待制周襟、馬觀國、史願、諸將程師回、馬欽、白常皆還金國〕相连。此处可确认的是这项被记录的行动；材料未将地点细化到城址，不能用中枢所在地替它补足。

把本项放回事件链，能够看到的前提为：本纪年条提供日期和中枢可见行动；前因不据单条补定。 而它所打开或收紧的下一步是：可回查同年及后续条目，地方影响仍须另证。

证据的边界是：《宋史》为元末官修，本纪保存宋廷纪年与处置；战果、数字、地方执行及对方动机须以编年、会要、对方史料或地方材料互校。 宋金战事、金代地方治理与碑志研究 可用于继续检验这一结论。

\subsection{1146年（南宋紹興十六年、金皇统6年；公元换算据两国年号对照）}
\noindent\eventanchor{SYD-1146-JIN-028}{九月甲戌，命何鑄等為金國祈請使，請國族}\textbf{南宋与金。}
〔九月甲戌，命何鑄等為金國祈請使，請國族〕见于 《宋史》卷030〈紹興十六年〉；《金史》本纪同年条（互校） 的 1146年（南宋紹興十六年、金皇统6年；公元换算据两国年号对照） 记录。南宋与金 的选择因此有了可回查的时间位置；材料未将地点细化到城址，不能用中枢所在地替它补足。

它承接的条件是：本纪年条记录政令或机构处置；实施背景须参照制度材料。 随后可追踪的变化为：可在同年及后续政令中追踪；条文不单独证明地方执行。

《宋史》为元末官修，本纪保存宋廷纪年与处置；战果、数字、地方执行及对方动机须以编年、会要、对方史料或地方材料互校。 因此，宋金战事、金代地方治理与碑志研究 只能扩展问题，不能代替未保存的细节。

\noindent\eventanchor{SYD-1146-JIN-127}{壬子，遣邊知白使金賀正旦，周執羔賀金主生辰}\textbf{南宋与金。}
1146年（南宋紹興十六年、金皇统6年；公元换算据两国年号对照），《宋史》卷030〈紹興十六年〉；《金史》本纪同年条（互校） 所见的〔壬子，遣邊知白使金賀正旦，周執羔賀金主生辰〕把 南宋与金 的处置留在可复核的年序中；材料未将地点细化到城址，不能用中枢所在地替它补足。

前一层压力可由以下线索辨认：本纪记录使节、贡赐或往来，不等于稳定统属关系。 这一处置留下的后续条件是：可与对方编年和交通、文书材料核对其后续落实。

关于可说范围，须保留：《宋史》为元末官修，本纪保存宋廷纪年与处置；战果、数字、地方执行及对方动机须以编年、会要、对方史料或地方材料互校。 进一步讨论可参照 宋金战事、金代地方治理与碑志研究

\subsection{1147年（南宋紹興十七年、金皇统7年；公元换算据两国年号对照）}
\noindent\eventanchor{SYD-1147-JIN-029}{戊申，遣沈該使金賀正旦，詹大方賀金主生辰}\textbf{南宋与金。}
南宋与金 在 1147年（南宋紹興十七年、金皇统7年；公元换算据两国年号对照） 的材料痕迹是〔戊申，遣沈該使金賀正旦，詹大方賀金主生辰〕，出处为 《宋史》卷030〈紹興十七年〉；《金史》本纪同年条（互校）。材料未将地点细化到城址，不能用中枢所在地替它补足。

本纪记录使节、贡赐或往来，不等于稳定统属关系。 记录所见的结果则是：可与对方编年和交通、文书材料核对其后续落实。 日期相邻不等于因果已经证实。

证据的边界是：《宋史》为元末官修，本纪保存宋廷纪年与处置；战果、数字、地方执行及对方动机须以编年、会要、对方史料或地方材料互校。 宋金战事、金代地方治理与碑志研究 可用于继续检验这一结论。

\subsection{1148年（南宋紹興十八年、金皇统8年；公元换算据两国年号对照）}
\noindent\eventanchor{SYD-1148-JIN-030}{壬申，遣王墨卿使金賀正旦，陳誠之賀金主生辰}\textbf{南宋与金。}
从 《宋史》卷030〈紹興十八年〉；《金史》本纪同年条（互校） 的 1148年（南宋紹興十八年、金皇统8年；公元换算据两国年号对照） 条目看，〔壬申，遣王墨卿使金賀正旦，陳誠之賀金主生辰〕使 南宋与金 的一个具体行动进入叙事；材料未将地点细化到城址，不能用中枢所在地替它补足。

把本项放回事件链，能够看到的前提为：本纪记录使节、贡赐或往来，不等于稳定统属关系。 而它所打开或收紧的下一步是：可与对方编年和交通、文书材料核对其后续落实。

《宋史》为元末官修，本纪保存宋廷纪年与处置；战果、数字、地方执行及对方动机须以编年、会要、对方史料或地方材料互校。 因此，宋金战事、金代地方治理与碑志研究 只能扩展问题，不能代替未保存的细节。

\subsection{1149年}
\noindent\eventanchor{JIN-1149-EVENT-093}{完颜亮发动政变即位}\textbf{金。}
1149年 的记载将 金 与〔完颜亮发动政变即位〕相连。此处可确认的是这项被记录的行动；材料未将地点细化到城址，不能用中枢所在地替它补足。

它承接的条件是：金在华北经营与宋金南北对峙中前任统治的结束与宫廷、部族或官僚的权力安排相互交织，促成“完颜亮发动政变即位”这一继承节点。 随后可追踪的变化为：继承并不自动消除既有矛盾；“完颜亮发动政变即位”后的任官、军权和对外策略需由后续条目追踪。

关于可说范围，须保留：《金史》为元修，宋金战报和官修合法性语言各有宣传性；军数、俘获、控制范围和地方执行须以对方及物质材料限缩。 进一步讨论可参照 金代国家形成、宋金边境、都城与金蒙转换研究

\subsection{1149年（南宋紹興十九年、金天德1年；公元换算据两国年号对照）}
\noindent\eventanchor{SYD-1149-JIN-031}{丙寅，秘閣修撰張邵上秦檜在金國代徽宗與粘罕書稿，詔付史館，以邵為徽猷閣待制}\textbf{南宋与金。}
〔丙寅，秘閣修撰張邵上秦檜在金國代徽宗與粘罕書稿，詔付史館，以邵為徽猷閣待制〕见于 《宋史》卷030〈紹興十九年〉；《金史》本纪同年条（互校） 的 1149年（南宋紹興十九年、金天德1年；公元换算据两国年号对照） 记录。南宋与金 的选择因此有了可回查的时间位置；材料未将地点细化到城址，不能用中枢所在地替它补足。

前一层压力可由以下线索辨认：本纪年条记录政令或机构处置；实施背景须参照制度材料。 这一处置留下的后续条件是：可在同年及后续政令中追踪；条文不单独证明地方执行。

证据的边界是：《宋史》为元末官修，本纪保存宋廷纪年与处置；战果、数字、地方执行及对方动机须以编年、会要、对方史料或地方材料互校。 宋金战事、金代地方治理与碑志研究 可用于继续检验这一结论。

\subsection{1150年（南宋紹興二十年、金天德2年；公元换算据两国年号对照）}
\noindent\eventanchor{SYD-1150-JIN-032}{丙戌，遣堯弼等賀金主即位}\textbf{南宋与金。}
1150年（南宋紹興二十年、金天德2年；公元换算据两国年号对照），《宋史》卷030〈紹興二十年〉；《金史》本纪同年条（互校） 所见的〔丙戌，遣堯弼等賀金主即位〕把 南宋与金 的处置留在可复核的年序中；材料未将地点细化到城址，不能用中枢所在地替它补足。

本纪年条提供日期和中枢可见行动；前因不据单条补定。 记录所见的结果则是：可回查同年及后续条目，地方影响仍须另证。 日期相邻不等于因果已经证实。

《宋史》为元末官修，本纪保存宋廷纪年与处置；战果、数字、地方执行及对方动机须以编年、会要、对方史料或地方材料互校。 因此，宋金战事、金代地方治理与碑志研究 只能扩展问题，不能代替未保存的细节。

\noindent\eventanchor{SYD-1150-JIN-128}{辛酉，遣陳誠之使金賀正旦，王〈日严〉賀金主生辰}\textbf{南宋与金。}
南宋与金 在 1150年（南宋紹興二十年、金天德2年；公元换算据两国年号对照） 的材料痕迹是〔辛酉，遣陳誠之使金賀正旦，王〈日严〉賀金主生辰〕，出处为 《宋史》卷030〈紹興二十年〉；《金史》本纪同年条（互校）。材料未将地点细化到城址，不能用中枢所在地替它补足。

把本项放回事件链，能够看到的前提为：本纪记录使节、贡赐或往来，不等于稳定统属关系。 而它所打开或收紧的下一步是：可与对方编年和交通、文书材料核对其后续落实。

关于可说范围，须保留：《宋史》为元末官修，本纪保存宋廷纪年与处置；战果、数字、地方执行及对方动机须以编年、会要、对方史料或地方材料互校。 进一步讨论可参照 宋金战事、金代地方治理与碑志研究

\subsection{1151年（南宋紹興二十一年、金天德3年；公元换算据两国年号对照）}
\noindent\eventanchor{SYD-1151-JIN-033}{甲申，遣陳夔使金賀正旦，陳相賀金主生辰}\textbf{南宋与金。}
从 《宋史》卷030〈紹興二十一年〉；《金史》本纪同年条（互校） 的 1151年（南宋紹興二十一年、金天德3年；公元换算据两国年号对照） 条目看，〔甲申，遣陳夔使金賀正旦，陳相賀金主生辰〕使 南宋与金 的一个具体行动进入叙事；材料未将地点细化到城址，不能用中枢所在地替它补足。

它承接的条件是：本纪记录使节、贡赐或往来，不等于稳定统属关系。 随后可追踪的变化为：可与对方编年和交通、文书材料核对其后续落实。

证据的边界是：《宋史》为元末官修，本纪保存宋廷纪年与处置；战果、数字、地方执行及对方动机须以编年、会要、对方史料或地方材料互校。 宋金战事、金代地方治理与碑志研究 可用于继续检验这一结论。

\noindent\eventanchor{SYD-1151-JIN-129}{壬戌，遣巫伋等為金國祈請使，請歸淵聖皇帝及皇族、增加帝號等事}\textbf{南宋与金。}
1151年（南宋紹興二十一年、金天德3年；公元换算据两国年号对照） 的记载将 南宋与金 与〔壬戌，遣巫伋等為金國祈請使，請歸淵聖皇帝及皇族、增加帝號等事〕相连。此处可确认的是这项被记录的行动；材料未将地点细化到城址，不能用中枢所在地替它补足。

前一层压力可由以下线索辨认：本纪记录使节、贡赐或往来，不等于稳定统属关系。 这一处置留下的后续条件是：可与对方编年和交通、文书材料核对其后续落实。

《宋史》为元末官修，本纪保存宋廷纪年与处置；战果、数字、地方执行及对方动机须以编年、会要、对方史料或地方材料互校。 因此，宋金战事、金代地方治理与碑志研究 只能扩展问题，不能代替未保存的细节。

\subsection{1153年}
\noindent\eventanchor{JIN-1153-EVENT-094}{金正式迁都中都}\textbf{金。}
〔金正式迁都中都〕见于 《金史》相关本纪；《宋史》或《三朝北盟会编》《建炎以来系年要录》对应年条；金代碑志与蒙古、高丽材料（据事核） 的 1153年 记录。金 的选择因此有了可回查的时间位置；题名保留的城、路、水道或机构名称，是目前可以落地的空间线索。

金在华北经营与宋金南北对峙对中枢资源、交通或行政协同的要求，推动了“金正式迁都中都”这一制度或空间安排。 记录所见的结果则是：“金正式迁都中都”提供新的治理框架或资源通道，但颁行、复申与不同地区的实际执行须分开核验。 日期相邻不等于因果已经证实。

关于可说范围，须保留：《金史》为元修，宋金战报和官修合法性语言各有宣传性；军数、俘获、控制范围和地方执行须以对方及物质材料限缩。 进一步讨论可参照 金代国家形成、宋金边境、都城与金蒙转换研究

\subsection{1153年（南宋紹興二十三年、金贞元1年；公元换算据两国年号对照）}
\noindent\eventanchor{SYD-1153-JIN-034}{戊午，遣吳㮚使金賀正旦，施鉅賀金主生辰}\textbf{南宋与金。}
1153年（南宋紹興二十三年、金贞元1年；公元换算据两国年号对照），《宋史》卷031〈紹興二十三年〉；《金史》本纪同年条（互校） 所见的〔戊午，遣吳㮚使金賀正旦，施鉅賀金主生辰〕把 南宋与金 的处置留在可复核的年序中；材料未将地点细化到城址，不能用中枢所在地替它补足。

把本项放回事件链，能够看到的前提为：本纪记录使节、贡赐或往来，不等于稳定统属关系。 而它所打开或收紧的下一步是：可与对方编年和交通、文书材料核对其后续落实。

证据的边界是：《宋史》为元末官修，本纪保存宋廷纪年与处置；战果、数字、地方执行及对方动机须以编年、会要、对方史料或地方材料互校。 宋金战事、金代地方治理与碑志研究 可用于继续检验这一结论。

\subsection{1154年（南宋紹興二十四年、金贞元2年；公元换算据两国年号对照）}
\noindent\eventanchor{SYD-1154-JIN-035}{戊子，遣沈虛中使金賀正旦，張士襄賀金主生辰}\textbf{南宋与金。}
南宋与金 在 1154年（南宋紹興二十四年、金贞元2年；公元换算据两国年号对照） 的材料痕迹是〔戊子，遣沈虛中使金賀正旦，張士襄賀金主生辰〕，出处为 《宋史》卷031〈紹興二十四年〉；《金史》本纪同年条（互校）。材料未将地点细化到城址，不能用中枢所在地替它补足。

它承接的条件是：本纪记录使节、贡赐或往来，不等于稳定统属关系。 随后可追踪的变化为：可与对方编年和交通、文书材料核对其后续落实。

《宋史》为元末官修，本纪保存宋廷纪年与处置；战果、数字、地方执行及对方动机须以编年、会要、对方史料或地方材料互校。 因此，宋金战事、金代地方治理与碑志研究 只能扩展问题，不能代替未保存的细节。

\subsection{1155年（南宋紹興二十五年、金贞元3年；公元换算据两国年号对照）}
\noindent\eventanchor{SYD-1155-JIN-036}{壬午，遣王岷使金賀正旦，鄭柟賀金主生辰}\textbf{南宋与金。}
从 《宋史》卷031〈紹興二十五年〉；《金史》本纪同年条（互校） 的 1155年（南宋紹興二十五年、金贞元3年；公元换算据两国年号对照） 条目看，〔壬午，遣王岷使金賀正旦，鄭柟賀金主生辰〕使 南宋与金 的一个具体行动进入叙事；材料未将地点细化到城址，不能用中枢所在地替它补足。

前一层压力可由以下线索辨认：本纪记录使节、贡赐或往来，不等于稳定统属关系。 这一处置留下的后续条件是：可与对方编年和交通、文书材料核对其后续落实。

关于可说范围，须保留：《宋史》为元末官修，本纪保存宋廷纪年与处置；战果、数字、地方执行及对方动机须以编年、会要、对方史料或地方材料互校。 进一步讨论可参照 宋金战事、金代地方治理与碑志研究

\noindent\eventanchor{SYD-1155-JIN-130}{三月己酉，右司郎中張士襄自金國使還，坐奉使不肅罷官}\textbf{南宋与金。}
1155年（南宋紹興二十五年、金贞元3年；公元换算据两国年号对照） 的记载将 南宋与金 与〔三月己酉，右司郎中張士襄自金國使還，坐奉使不肅罷官〕相连。此处可确认的是这项被记录的行动；材料未将地点细化到城址，不能用中枢所在地替它补足。

本纪年条记录政令或机构处置；实施背景须参照制度材料。 记录所见的结果则是：可在同年及后续政令中追踪；条文不单独证明地方执行。 日期相邻不等于因果已经证实。

证据的边界是：《宋史》为元末官修，本纪保存宋廷纪年与处置；战果、数字、地方执行及对方动机须以编年、会要、对方史料或地方材料互校。 宋金战事、金代地方治理与碑志研究 可用于继续检验这一结论。

\subsection{1156年}
\noindent\eventanchor{JIN-1156-EVENT-095}{完颜亮集结军队与物资准备南侵}\textbf{金。}
〔完颜亮集结军队与物资准备南侵〕见于 《金史》相关本纪；《宋史》或《三朝北盟会编》《建炎以来系年要录》对应年条；金代碑志与蒙古、高丽材料（据事核） 的 1156年 记录。金 的选择因此有了可回查的时间位置；材料未将地点细化到城址，不能用中枢所在地替它补足。

把本项放回事件链，能够看到的前提为：金在华北经营与宋金南北对峙中既有的联盟、交通与补给压力，为“完颜亮集结军队与物资准备南侵”提供了直接的军事或动员背景。 而它所打开或收紧的下一步是：“完颜亮集结军队与物资准备南侵”改变了相关城镇、道路或政权间的军事选择；伤亡、俘获和控制范围仅可按各方材料分别确认。

《金史》为元修，宋金战报和官修合法性语言各有宣传性；军数、俘获、控制范围和地方执行须以对方及物质材料限缩。 因此，金代国家形成、宋金边境、都城与金蒙转换研究 只能扩展问题，不能代替未保存的细节。

\subsection{1156年（南宋紹興二十六年、金正隆1年；公元换算据两国年号对照）}
\noindent\eventanchor{SYD-1156-JIN-037}{庚寅，遣陳誠之等賀金主尊號禮成}\textbf{南宋与金。}
1156年（南宋紹興二十六年、金正隆1年；公元换算据两国年号对照），《宋史》卷031〈紹興二十六年〉；《金史》本纪同年条（互校） 所见的〔庚寅，遣陳誠之等賀金主尊號禮成〕把 南宋与金 的处置留在可复核的年序中；材料未将地点细化到城址，不能用中枢所在地替它补足。

它承接的条件是：本纪年条提供日期和中枢可见行动；前因不据单条补定。 随后可追踪的变化为：可回查同年及后续条目，地方影响仍须另证。

关于可说范围，须保留：《宋史》为元末官修，本纪保存宋廷纪年与处置；战果、数字、地方执行及对方动机须以编年、会要、对方史料或地方材料互校。 进一步讨论可参照 宋金战事、金代地方治理与碑志研究

\noindent\eventanchor{SYD-1156-JIN-131}{遣李琳使金賀正旦，葛立方賀金主生辰}\textbf{南宋与金。}
南宋与金 在 1156年（南宋紹興二十六年、金正隆1年；公元换算据两国年号对照） 的材料痕迹是〔遣李琳使金賀正旦，葛立方賀金主生辰〕，出处为 《宋史》卷031〈紹興二十六年〉；《金史》本纪同年条（互校）。材料未将地点细化到城址，不能用中枢所在地替它补足。

前一层压力可由以下线索辨认：本纪记录使节、贡赐或往来，不等于稳定统属关系。 这一处置留下的后续条件是：可与对方编年和交通、文书材料核对其后续落实。

证据的边界是：《宋史》为元末官修，本纪保存宋廷纪年与处置；战果、数字、地方执行及对方动机须以编年、会要、对方史料或地方材料互校。 宋金战事、金代地方治理与碑志研究 可用于继续检验这一结论。

\subsection{1157年（南宋紹興二十七年、金正隆2年；公元换算据两国年号对照）}
\noindent\eventanchor{SYD-1157-JIN-038}{辛巳，遣劉章賀金主生辰，丁亥，湯鵬舉罷}\textbf{南宋与金。}
从 《宋史》卷031〈紹興二十七年〉；《金史》本纪同年条（互校） 的 1157年（南宋紹興二十七年、金正隆2年；公元换算据两国年号对照） 条目看，〔辛巳，遣劉章賀金主生辰，丁亥，湯鵬舉罷〕使 南宋与金 的一个具体行动进入叙事；材料未将地点细化到城址，不能用中枢所在地替它补足。

本纪年条记录政令或机构处置；实施背景须参照制度材料。 记录所见的结果则是：可在同年及后续政令中追踪；条文不单独证明地方执行。 日期相邻不等于因果已经证实。

《宋史》为元末官修，本纪保存宋廷纪年与处置；战果、数字、地方执行及对方动机须以编年、会要、对方史料或地方材料互校。 因此，宋金战事、金代地方治理与碑志研究 只能扩展问题，不能代替未保存的细节。

\subsection{1158年（南宋紹興二十八年、金正隆3年；公元换算据两国年号对照）}
\noindent\eventanchor{SYD-1158-JIN-039}{冬十月丁亥朔，遣沈介使金賀正旦，黃中賀金主生辰}\textbf{南宋与金。}
1158年（南宋紹興二十八年、金正隆3年；公元换算据两国年号对照） 的记载将 南宋与金 与〔冬十月丁亥朔，遣沈介使金賀正旦，黃中賀金主生辰〕相连。此处可确认的是这项被记录的行动；材料未将地点细化到城址，不能用中枢所在地替它补足。

把本项放回事件链，能够看到的前提为：本纪记录使节、贡赐或往来，不等于稳定统属关系。 而它所打开或收紧的下一步是：可与对方编年和交通、文书材料核对其后续落实。

关于可说范围，须保留：《宋史》为元末官修，本纪保存宋廷纪年与处置；战果、数字、地方执行及对方动机须以编年、会要、对方史料或地方材料互校。 进一步讨论可参照 宋金战事、金代地方治理与碑志研究

\subsection{1159年（南宋紹興二十九年、金正隆4年；公元换算据两国年号对照）}
\noindent\eventanchor{SYD-1159-JIN-040}{癸卯，遣周麟之等為金國奉表哀謝使}\textbf{南宋与金。}
〔癸卯，遣周麟之等為金國奉表哀謝使〕见于 《宋史》卷031〈紹興二十九年〉；《金史》本纪同年条（互校） 的 1159年（南宋紹興二十九年、金正隆4年；公元换算据两国年号对照） 记录。南宋与金 的选择因此有了可回查的时间位置；材料未将地点细化到城址，不能用中枢所在地替它补足。

它承接的条件是：本纪记录使节、贡赐或往来，不等于稳定统属关系。 随后可追踪的变化为：可与对方编年和交通、文书材料核对其后续落实。

证据的边界是：《宋史》为元末官修，本纪保存宋廷纪年与处置；战果、数字、地方执行及对方动机须以编年、会要、对方史料或地方材料互校。 宋金战事、金代地方治理与碑志研究 可用于继续检验这一结论。

\noindent\eventanchor{SYD-1159-JIN-132}{六月甲辰朔，遣王綸等為金國奉表稱謝使}\textbf{南宋与金。}
1159年（南宋紹興二十九年、金正隆4年；公元换算据两国年号对照），《宋史》卷031〈紹興二十九年〉；《金史》本纪同年条（互校） 所见的〔六月甲辰朔，遣王綸等為金國奉表稱謝使〕把 南宋与金 的处置留在可复核的年序中；材料未将地点细化到城址，不能用中枢所在地替它补足。

前一层压力可由以下线索辨认：本纪记录使节、贡赐或往来，不等于稳定统属关系。 这一处置留下的后续条件是：可与对方编年和交通、文书材料核对其后续落实。

《宋史》为元末官修，本纪保存宋廷纪年与处置；战果、数字、地方执行及对方动机须以编年、会要、对方史料或地方材料互校。 因此，宋金战事、金代地方治理与碑志研究 只能扩展问题，不能代替未保存的细节。

\noindent\eventanchor{SYD-1159-JIN-179}{十一月丁亥，遣賀允中等為金國遺留國信使}\textbf{南宋与金。}
南宋与金 在 1159年（南宋紹興二十九年、金正隆4年；公元换算据两国年号对照） 的材料痕迹是〔十一月丁亥，遣賀允中等為金國遺留國信使〕，出处为 《宋史》卷031〈紹興二十九年〉；《金史》本纪同年条（互校）。材料未将地点细化到城址，不能用中枢所在地替它补足。

本纪记录使节、贡赐或往来，不等于稳定统属关系。 记录所见的结果则是：可与对方编年和交通、文书材料核对其后续落实。 日期相邻不等于因果已经证实。

关于可说范围，须保留：《宋史》为元末官修，本纪保存宋廷纪年与处置；战果、数字、地方执行及对方动机须以编年、会要、对方史料或地方材料互校。 进一步讨论可参照 宋金战事、金代地方治理与碑志研究

\noindent\eventanchor{SYD-1159-JIN-215}{是月，金國罷沿邊榷場，惟泗州如舊}\textbf{南宋与金。}
从 《宋史》卷031〈紹興二十九年〉；《金史》本纪同年条（互校） 的 1159年（南宋紹興二十九年、金正隆4年；公元换算据两国年号对照） 条目看，〔是月，金國罷沿邊榷場，惟泗州如舊〕使 南宋与金 的一个具体行动进入叙事；题名保留的城、路、水道或机构名称，是目前可以落地的空间线索。

把本项放回事件链，能够看到的前提为：本纪年条记录政令或机构处置；实施背景须参照制度材料。 而它所打开或收紧的下一步是：可在同年及后续政令中追踪；条文不单独证明地方执行。

证据的边界是：《宋史》为元末官修，本纪保存宋廷纪年与处置；战果、数字、地方执行及对方动机须以编年、会要、对方史料或地方材料互校。 宋金战事、金代地方治理与碑志研究 可用于继续检验这一结论。

\noindent\eventanchor{SYD-1159-JIN-244}{乙酉，王綸使還入見，言金國和好無他}\textbf{南宋与金。}
1159年（南宋紹興二十九年、金正隆4年；公元换算据两国年号对照） 的记载将 南宋与金 与〔乙酉，王綸使還入見，言金國和好無他〕相连。此处可确认的是这项被记录的行动；材料未将地点细化到城址，不能用中枢所在地替它补足。

它承接的条件是：本纪年条提供日期和中枢可见行动；前因不据单条补定。 随后可追踪的变化为：可回查同年及后续条目，地方影响仍须另证。

《宋史》为元末官修，本纪保存宋廷纪年与处置；战果、数字、地方执行及对方动机须以编年、会要、对方史料或地方材料互校。 因此，宋金战事、金代地方治理与碑志研究 只能扩展问题，不能代替未保存的细节。

\subsection{1160年（南宋紹興三十年、金正隆5年；公元换算据两国年号对照）}
\noindent\eventanchor{SYD-1160-JIN-041}{丁未，遣虞允文使金賀正旦，徐度賀金主生辰}\textbf{南宋与金。}
〔丁未，遣虞允文使金賀正旦，徐度賀金主生辰〕见于 《宋史》卷031〈紹興三十年〉；《金史》本纪同年条（互校） 的 1160年（南宋紹興三十年、金正隆5年；公元换算据两国年号对照） 记录。南宋与金 的选择因此有了可回查的时间位置；材料未将地点细化到城址，不能用中枢所在地替它补足。

前一层压力可由以下线索辨认：本纪记录使节、贡赐或往来，不等于稳定统属关系。 这一处置留下的后续条件是：可与对方编年和交通、文书材料核对其后续落实。

关于可说范围，须保留：《宋史》为元末官修，本纪保存宋廷纪年与处置；战果、数字、地方执行及对方动机须以编年、会要、对方史料或地方材料互校。 进一步讨论可参照 宋金战事、金代地方治理与碑志研究

\noindent\eventanchor{SYD-1160-JIN-133}{壬申，淮東總管許世安奏，金主亮至汴京，起重兵五十余萬，屯宿、泗州，謀來攻}\textbf{南宋与金。}
1160年（南宋紹興三十年、金正隆5年；公元换算据两国年号对照），《宋史》卷031〈紹興三十年〉；《金史》本纪同年条（互校） 所见的〔壬申，淮東總管許世安奏，金主亮至汴京，起重兵五十余萬，屯宿、泗州，謀來攻〕把 南宋与金 的处置留在可复核的年序中；题名保留的城、路、水道或机构名称，是目前可以落地的空间线索。

本纪从朝廷视角记录军政行动；出兵与战果需分辨。 记录所见的结果则是：后续路线、控制与损失应与对方记录和遗址材料复核。 日期相邻不等于因果已经证实。

证据的边界是：《宋史》为元末官修，本纪保存宋廷纪年与处置；战果、数字、地方执行及对方动机须以编年、会要、对方史料或地方材料互校。 宋金战事、金代地方治理与碑志研究 可用于继续检验这一结论。

\noindent\eventanchor{SYD-1160-JIN-180}{壬子，賀允中使還，言金人必叛盟，宜為之備}\textbf{南宋与金。}
南宋与金 在 1160年（南宋紹興三十年、金正隆5年；公元换算据两国年号对照） 的材料痕迹是〔壬子，賀允中使還，言金人必叛盟，宜為之備〕，出处为 《宋史》卷031〈紹興三十年〉；《金史》本纪同年条（互校）。材料未将地点细化到城址，不能用中枢所在地替它补足。

把本项放回事件链，能够看到的前提为：本纪从朝廷视角记录军政行动；出兵与战果需分辨。 而它所打开或收紧的下一步是：后续路线、控制与损失应与对方记录和遗址材料复核。

《宋史》为元末官修，本纪保存宋廷纪年与处置；战果、数字、地方执行及对方动机须以编年、会要、对方史料或地方材料互校。 因此，宋金战事、金代地方治理与碑志研究 只能扩展问题，不能代替未保存的细节。

\noindent\eventanchor{SYD-1160-JIN-216}{戊午，遣葉義問為金國報謝使}\textbf{南宋与金。}
从 《宋史》卷031〈紹興三十年〉；《金史》本纪同年条（互校） 的 1160年（南宋紹興三十年、金正隆5年；公元换算据两国年号对照） 条目看，〔戊午，遣葉義問為金國報謝使〕使 南宋与金 的一个具体行动进入叙事；材料未将地点细化到城址，不能用中枢所在地替它补足。

它承接的条件是：本纪记录使节、贡赐或往来，不等于稳定统属关系。 随后可追踪的变化为：可与对方编年和交通、文书材料核对其后续落实。

关于可说范围，须保留：《宋史》为元末官修，本纪保存宋廷纪年与处置；战果、数字、地方执行及对方动机须以编年、会要、对方史料或地方材料互校。 进一步讨论可参照 宋金战事、金代地方治理与碑志研究

\subsection{1161年（金正隆六年、宋绍兴三十一年；干支、月日待核；公元换算据双方本纪年号对照）}
\noindent\eventanchor{JIN-1161-EVENT-096}{金军南侵在采石等战场受挫，完颜亮被杀}\textbf{金与南宋。}
1161年（金正隆六年、宋绍兴三十一年；干支、月日待核；公元换算据双方本纪年号对照） 的记载将 金与南宋 与〔金军南侵在采石等战场受挫，完颜亮被杀〕相连。此处可确认的是这项被记录的行动；材料未将地点细化到城址，不能用中枢所在地替它补足。

前一层压力可由以下线索辨认：金在华北经营与宋金南北对峙中既有的联盟、交通与补给压力，为“金军南侵在采石等战场受挫，完颜亮被杀”提供了直接的军事或动员背景。 这一处置留下的后续条件是：“金军南侵在采石等战场受挫，完颜亮被杀”改变了相关城镇、道路或政权间的军事选择；伤亡、俘获和控制范围仅可按各方材料分别确认。

证据的边界是：《金史》为元修，宋金战报和官修合法性语言各有宣传性；军数、俘获、控制范围和地方执行须以对方及物质材料限缩。 金代国家形成、宋金边境、都城与金蒙转换研究 可用于继续检验这一结论。

\subsection{1161年}
\noindent\eventanchor{JIN-1161-EVENT-097}{完颜雍即位为金世宗}\textbf{金。}
〔完颜雍即位为金世宗〕见于 《金史》相关本纪；《宋史》或《三朝北盟会编》《建炎以来系年要录》对应年条；金代碑志与蒙古、高丽材料（据事核） 的 1161年 记录。金 的选择因此有了可回查的时间位置；材料未将地点细化到城址，不能用中枢所在地替它补足。

金在华北经营与宋金南北对峙中前任统治的结束与宫廷、部族或官僚的权力安排相互交织，促成“完颜雍即位为金世宗”这一继承节点。 记录所见的结果则是：继承并不自动消除既有矛盾；“完颜雍即位为金世宗”后的任官、军权和对外策略需由后续条目追踪。 日期相邻不等于因果已经证实。

《金史》为元修，宋金战报和官修合法性语言各有宣传性；军数、俘获、控制范围和地方执行须以对方及物质材料限缩。 因此，金代国家形成、宋金边境、都城与金蒙转换研究 只能扩展问题，不能代替未保存的细节。

\subsection{1161年（南宋紹興三十一年、金大定1年；公元换算据两国年号对照）}
\noindent\eventanchor{SYD-1161-JIN-042}{丙辰，金主亮入廬州，王權自昭關遁，金人追至尉子橋，破敵軍統制姚興戰死，權退保和州}\textbf{南宋与金。}
1161年（南宋紹興三十一年、金大定1年；公元换算据两国年号对照），《宋史》卷032〈紹興三十一年〉；《金史》本纪同年条（互校） 所见的〔丙辰，金主亮入廬州，王權自昭關遁，金人追至尉子橋，破敵軍統制姚興戰死，權退保和州〕把 南宋与金 的处置留在可复核的年序中；题名保留的城、路、水道或机构名称，是目前可以落地的空间线索。

把本项放回事件链，能够看到的前提为：本纪从朝廷视角记录军政行动；出兵与战果需分辨。 而它所打开或收紧的下一步是：后续路线、控制与损失应与对方记录和遗址材料复核。

关于可说范围，须保留：《宋史》为元末官修，本纪保存宋廷纪年与处置；战果、数字、地方执行及对方动机须以编年、会要、对方史料或地方材料互校。 进一步讨论可参照 宋金战事、金代地方治理与碑志研究

\noindent\eventanchor{SYD-1161-JIN-134}{丙子，虞允文督建康諸軍統制官張振、王琪、時俊、戴皋等以舟師拒金主亮於東採石，戰勝，卻之}\textbf{南宋与金。}
南宋与金 在 1161年（南宋紹興三十一年、金大定1年；公元换算据两国年号对照） 的材料痕迹是〔丙子，虞允文督建康諸軍統制官張振、王琪、時俊、戴皋等以舟師拒金主亮於東採石，戰勝，卻之〕，出处为 《宋史》卷032〈紹興三十一年〉；《金史》本纪同年条（互校）。材料未将地点细化到城址，不能用中枢所在地替它补足。

它承接的条件是：本纪从朝廷视角记录军政行动；出兵与战果需分辨。 随后可追踪的变化为：后续路线、控制与损失应与对方记录和遗址材料复核。

证据的边界是：《宋史》为元末官修，本纪保存宋廷纪年与处置；战果、数字、地方执行及对方动机须以编年、会要、对方史料或地方材料互校。 宋金战事、金代地方治理与碑志研究 可用于继续检验这一结论。

\noindent\eventanchor{SYD-1161-JIN-181}{丁丑，虞允文遣水軍統制盛新以舟師擊金人于楊林河口，又敗之}\textbf{南宋与金。}
从 《宋史》卷032〈紹興三十一年〉；《金史》本纪同年条（互校） 的 1161年（南宋紹興三十一年、金大定1年；公元换算据两国年号对照） 条目看，〔丁丑，虞允文遣水軍統制盛新以舟師擊金人于楊林河口，又敗之〕使 南宋与金 的一个具体行动进入叙事；题名保留的城、路、水道或机构名称，是目前可以落地的空间线索。

前一层压力可由以下线索辨认：本纪保留灾害或工程的报告地点与朝廷处置；灾情范围需另证。 这一处置留下的后续条件是：可与地方文书、河工和环境材料比较，不将单点记录外推为全域因果。

《宋史》为元末官修，本纪保存宋廷纪年与处置；战果、数字、地方执行及对方动机须以编年、会要、对方史料或地方材料互校。 因此，宋金战事、金代地方治理与碑志研究 只能扩展问题，不能代替未保存的细节。

\noindent\eventanchor{SYD-1161-JIN-217}{丁未，命宣撫制置司傳檄契丹、西夏、高麗、渤海諸國及河北、河東、陝西、京東、河南諸路，諭出師共討金人}\textbf{南宋与金。}
1161年（南宋紹興三十一年、金大定1年；公元换算据两国年号对照） 的记载将 南宋与金 与〔丁未，命宣撫制置司傳檄契丹、西夏、高麗、渤海諸國及河北、河東、陝西、京東、河南諸路，諭出師共討金人〕相连。此处可确认的是这项被记录的行动；题名保留的城、路、水道或机构名称，是目前可以落地的空间线索。

本纪年条记录政令或机构处置；实施背景须参照制度材料。 记录所见的结果则是：可在同年及后续政令中追踪；条文不单独证明地方执行。 日期相邻不等于因果已经证实。

关于可说范围，须保留：《宋史》为元末官修，本纪保存宋廷纪年与处置；战果、数字、地方执行及对方动机须以编年、会要、对方史料或地方材料互校。 进一步讨论可参照 宋金战事、金代地方治理与碑志研究

\noindent\eventanchor{SYD-1161-JIN-245}{鄂州統制楊欽以舟師追敗金人于洪澤鎮}\textbf{南宋与金。}
〔鄂州統制楊欽以舟師追敗金人于洪澤鎮〕见于 《宋史》卷032〈紹興三十一年〉；《金史》本纪同年条（互校） 的 1161年（南宋紹興三十一年、金大定1年；公元换算据两国年号对照） 记录。南宋与金 的选择因此有了可回查的时间位置；题名保留的城、路、水道或机构名称，是目前可以落地的空间线索。

把本项放回事件链，能够看到的前提为：本纪从朝廷视角记录军政行动；出兵与战果需分辨。 而它所打开或收紧的下一步是：后续路线、控制与损失应与对方记录和遗址材料复核。

证据的边界是：《宋史》为元末官修，本纪保存宋廷纪年与处置；战果、数字、地方执行及对方动机须以编年、会要、对方史料或地方材料互校。 宋金战事、金代地方治理与碑志研究 可用于继续检验这一结论。

\subsection{1162年（南宋紹興三十二年、金大定2年；公元换算据两国年号对照）}
\noindent\eventanchor{SYD-1162-JIN-043}{丁巳，遣洪邁等賀金主即位}\textbf{南宋与金。}
1162年（南宋紹興三十二年、金大定2年；公元换算据两国年号对照），《宋史》卷032〈紹興三十二年〉；《金史》本纪同年条（互校） 所见的〔丁巳，遣洪邁等賀金主即位〕把 南宋与金 的处置留在可复核的年序中；材料未将地点细化到城址，不能用中枢所在地替它补足。

它承接的条件是：本纪年条提供日期和中枢可见行动；前因不据单条补定。 随后可追踪的变化为：可回查同年及后续条目，地方影响仍须另证。

《宋史》为元末官修，本纪保存宋廷纪年与处置；战果、数字、地方执行及对方动机须以编年、会要、对方史料或地方材料互校。 因此，宋金战事、金代地方治理与碑志研究 只能扩展问题，不能代替未保存的细节。

\noindent\eventanchor{SYD-1162-JIN-135}{庚戌，金人全師來攻，宣敗績，棄去}\textbf{南宋与金。}
南宋与金 在 1162年（南宋紹興三十二年、金大定2年；公元换算据两国年号对照） 的材料痕迹是〔庚戌，金人全師來攻，宣敗績，棄去〕，出处为 《宋史》卷032〈紹興三十二年〉；《金史》本纪同年条（互校）。材料未将地点细化到城址，不能用中枢所在地替它补足。

前一层压力可由以下线索辨认：本纪从朝廷视角记录军政行动；出兵与战果需分辨。 这一处置留下的后续条件是：后续路线、控制与损失应与对方记录和遗址材料复核。

关于可说范围，须保留：《宋史》为元末官修，本纪保存宋廷纪年与处置；战果、数字、地方执行及对方动机须以编年、会要、对方史料或地方材料互校。 进一步讨论可参照 宋金战事、金代地方治理与碑志研究

\noindent\eventanchor{SYD-1162-JIN-182}{庚寅，王剛破金人於海州}\textbf{南宋与金。}
从 《宋史》卷032〈紹興三十二年〉；《金史》本纪同年条（互校） 的 1162年（南宋紹興三十二年、金大定2年；公元换算据两国年号对照） 条目看，〔庚寅，王剛破金人於海州〕使 南宋与金 的一个具体行动进入叙事；题名保留的城、路、水道或机构名称，是目前可以落地的空间线索。

本纪从朝廷视角记录军政行动；出兵与战果需分辨。 记录所见的结果则是：后续路线、控制与损失应与对方记录和遗址材料复核。 日期相邻不等于因果已经证实。

证据的边界是：《宋史》为元末官修，本纪保存宋廷纪年与处置；战果、数字、地方执行及对方动机须以编年、会要、对方史料或地方材料互校。 宋金战事、金代地方治理与碑志研究 可用于继续检验这一结论。

\noindent\eventanchor{SYD-1162-JIN-218}{癸丑，金人圍淮寧府，守臣陳亨祖死之}\textbf{南宋与金。}
1162年（南宋紹興三十二年、金大定2年；公元换算据两国年号对照） 的记载将 南宋与金 与〔癸丑，金人圍淮寧府，守臣陳亨祖死之〕相连。此处可确认的是这项被记录的行动；题名保留的城、路、水道或机构名称，是目前可以落地的空间线索。

把本项放回事件链，能够看到的前提为：本纪从朝廷视角记录军政行动；出兵与战果需分辨。 而它所打开或收紧的下一步是：后续路线、控制与损失应与对方记录和遗址材料复核。

《宋史》为元末官修，本纪保存宋廷纪年与处置；战果、数字、地方执行及对方动机须以编年、会要、对方史料或地方材料互校。 因此，宋金战事、金代地方治理与碑志研究 只能扩展问题，不能代替未保存的细节。

\noindent\eventanchor{SYD-1162-JIN-246}{癸卯，成閔遣統制杜彥救淮甯，擊敗金人于項城縣}\textbf{南宋与金。}
〔癸卯，成閔遣統制杜彥救淮甯，擊敗金人于項城縣〕见于 《宋史》卷032〈紹興三十二年〉；《金史》本纪同年条（互校） 的 1162年（南宋紹興三十二年、金大定2年；公元换算据两国年号对照） 记录。南宋与金 的选择因此有了可回查的时间位置；题名保留的城、路、水道或机构名称，是目前可以落地的空间线索。

它承接的条件是：本纪从朝廷视角记录军政行动；出兵与战果需分辨。 随后可追踪的变化为：后续路线、控制与损失应与对方记录和遗址材料复核。

关于可说范围，须保留：《宋史》为元末官修，本纪保存宋廷纪年与处置；战果、数字、地方执行及对方动机须以编年、会要、对方史料或地方材料互校。 进一步讨论可参照 宋金战事、金代地方治理与碑志研究

\subsection{1163年（南宋隆興元年、金大定3年；公元换算据两国年号对照）}
\noindent\eventanchor{SYD-1163-JIN-044}{丙午，復宿州，戮金兵數千人}\textbf{南宋与金。}
1163年（南宋隆興元年、金大定3年；公元换算据两国年号对照），《宋史》卷033〈隆興元年〉；《金史》本纪同年条（互校） 所见的〔丙午，復宿州，戮金兵數千人〕把 南宋与金 的处置留在可复核的年序中；题名保留的城、路、水道或机构名称，是目前可以落地的空间线索。

前一层压力可由以下线索辨认：本纪年条记录政令或机构处置；实施背景须参照制度材料。 这一处置留下的后续条件是：可在同年及后续政令中追踪；条文不单独证明地方执行。

证据的边界是：《宋史》为元末官修，本纪保存宋廷纪年与处置；战果、数字、地方执行及对方动机须以编年、会要、对方史料或地方材料互校。 宋金战事、金代地方治理与碑志研究 可用于继续检验这一结论。

\noindent\eventanchor{SYD-1163-JIN-136}{庚子，遣王子望等爲金國通問使}\textbf{南宋与金。}
南宋与金 在 1163年（南宋隆興元年、金大定3年；公元换算据两国年号对照） 的材料痕迹是〔庚子，遣王子望等爲金國通問使〕，出处为 《宋史》卷033〈隆興元年〉；《金史》本纪同年条（互校）。材料未将地点细化到城址，不能用中枢所在地替它补足。

本纪记录使节、贡赐或往来，不等于稳定统属关系。 记录所见的结果则是：可与对方编年和交通、文书材料核对其后续落实。 日期相邻不等于因果已经证实。

《宋史》为元末官修，本纪保存宋廷纪年与处置；战果、数字、地方执行及对方动机须以编年、会要、对方史料或地方材料互校。 因此，宋金战事、金代地方治理与碑志研究 只能扩展问题，不能代替未保存的细节。

\noindent\eventanchor{SYD-1163-JIN-183}{甲辰，顯忠及宏淵敗金人於宿州}\textbf{南宋与金。}
从 《宋史》卷033〈隆興元年〉；《金史》本纪同年条（互校） 的 1163年（南宋隆興元年、金大定3年；公元换算据两国年号对照） 条目看，〔甲辰，顯忠及宏淵敗金人於宿州〕使 南宋与金 的一个具体行动进入叙事；题名保留的城、路、水道或机构名称，是目前可以落地的空间线索。

把本项放回事件链，能够看到的前提为：本纪从朝廷视角记录军政行动；出兵与战果需分辨。 而它所打开或收紧的下一步是：后续路线、控制与损失应与对方记录和遗址材料复核。

关于可说范围，须保留：《宋史》为元末官修，本纪保存宋廷纪年与处置；战果、数字、地方执行及对方动机须以编年、会要、对方史料或地方材料互校。 进一步讨论可参照 宋金战事、金代地方治理与碑志研究

\noindent\eventanchor{SYD-1163-JIN-219}{金人攻宿州城，顯忠大敗之}\textbf{南宋与金。}
1163年（南宋隆興元年、金大定3年；公元换算据两国年号对照） 的记载将 南宋与金 与〔金人攻宿州城，顯忠大敗之〕相连。此处可确认的是这项被记录的行动；题名保留的城、路、水道或机构名称，是目前可以落地的空间线索。

它承接的条件是：本纪从朝廷视角记录军政行动；出兵与战果需分辨。 随后可追踪的变化为：后续路线、控制与损失应与对方记录和遗址材料复核。

证据的边界是：《宋史》为元末官修，本纪保存宋廷纪年与处置；战果、数字、地方执行及对方动机须以编年、会要、对方史料或地方材料互校。 宋金战事、金代地方治理与碑志研究 可用于继续检验这一结论。

\noindent\eventanchor{SYD-1163-JIN-247}{璘已棄德順，道爲金人所邀，將士死者數萬計}\textbf{南宋与金。}
〔璘已棄德順，道爲金人所邀，將士死者數萬計〕见于 《宋史》卷033〈隆興元年〉；《金史》本纪同年条（互校） 的 1163年（南宋隆興元年、金大定3年；公元换算据两国年号对照） 记录。南宋与金 的选择因此有了可回查的时间位置；材料未将地点细化到城址，不能用中枢所在地替它补足。

前一层压力可由以下线索辨认：本纪年条提供日期和中枢可见行动；前因不据单条补定。 这一处置留下的后续条件是：可回查同年及后续条目，地方影响仍须另证。

《宋史》为元末官修，本纪保存宋廷纪年与处置；战果、数字、地方执行及对方动机须以编年、会要、对方史料或地方材料互校。 因此，宋金战事、金代地方治理与碑志研究 只能扩展问题，不能代替未保存的细节。

\subsection{1164年（金大定四年、宋隆兴二年；干支、月日待核；公元换算据双方本纪年号对照）}
\noindent\eventanchor{JIN-1164-EVENT-098}{金宋缔结隆兴和议}\textbf{金与南宋。}
1164年（金大定四年、宋隆兴二年；干支、月日待核；公元换算据双方本纪年号对照），《金史》相关本纪；《宋史》或《三朝北盟会编》《建炎以来系年要录》对应年条；金代碑志与蒙古、高丽材料（据事核） 所见的〔金宋缔结隆兴和议〕把 金与南宋 的处置留在可复核的年序中；材料未将地点细化到城址，不能用中枢所在地替它补足。

金世宗至章宗时期的制度与边境压力里的军事消耗、边境供给与使节往来构成谈判条件，促成“金宋缔结隆兴和议”所涉的协议或名义关系。 记录所见的结果则是：“金宋缔结隆兴和议”重排了贡赐、使节或停战的制度接口；条款是否执行及边地实际控制不能由盟约文字直接推定。 日期相邻不等于因果已经证实。

关于可说范围，须保留：《金史》为元修，宋金战报和官修合法性语言各有宣传性；军数、俘获、控制范围和地方执行须以对方及物质材料限缩。 进一步讨论可参照 金代国家形成、宋金边境、都城与金蒙转换研究

\subsection{1164年（南宋隆興二年、金大定4年；公元换算据两国年号对照）}
\noindent\eventanchor{SYD-1164-JIN-045}{遣洪適等賀金主生辰}\textbf{南宋与金。}
南宋与金 在 1164年（南宋隆興二年、金大定4年；公元换算据两国年号对照） 的材料痕迹是〔遣洪適等賀金主生辰〕，出处为 《宋史》卷033〈隆興二年〉；《金史》本纪同年条（互校）。材料未将地点细化到城址，不能用中枢所在地替它补足。

把本项放回事件链，能够看到的前提为：本纪年条提供日期和中枢可见行动；前因不据单条补定。 而它所打开或收紧的下一步是：可回查同年及后续条目，地方影响仍须另证。

证据的边界是：《宋史》为元末官修，本纪保存宋廷纪年与处置；战果、数字、地方执行及对方动机须以编年、会要、对方史料或地方材料互校。 宋金战事、金代地方治理与碑志研究 可用于继续检验这一结论。

\noindent\eventanchor{SYD-1164-JIN-137}{初，金帥以昉等不許四郡，械系之，昉等不屈，金主命歸之}\textbf{南宋与金。}
从 《宋史》卷033〈隆興二年〉；《金史》本纪同年条（互校） 的 1164年（南宋隆興二年、金大定4年；公元换算据两国年号对照） 条目看，〔初，金帥以昉等不許四郡，械系之，昉等不屈，金主命歸之〕使 南宋与金 的一个具体行动进入叙事；材料未将地点细化到城址，不能用中枢所在地替它补足。

它承接的条件是：本纪年条记录政令或机构处置；实施背景须参照制度材料。 随后可追踪的变化为：可在同年及后续政令中追踪；条文不单独证明地方执行。

《宋史》为元末官修，本纪保存宋廷纪年与处置；战果、数字、地方执行及对方动机须以编年、会要、对方史料或地方材料互校。 因此，宋金战事、金代地方治理与碑志研究 只能扩展问题，不能代替未保存的细节。

\noindent\eventanchor{SYD-1164-JIN-184}{甲辰，金人犯六合縣，步軍司統制崔皋擊卻之}\textbf{南宋与金。}
1164年（南宋隆興二年、金大定4年；公元换算据两国年号对照） 的记载将 南宋与金 与〔甲辰，金人犯六合縣，步軍司統制崔皋擊卻之〕相连。此处可确认的是这项被记录的行动；材料未将地点细化到城址，不能用中枢所在地替它补足。

前一层压力可由以下线索辨认：本纪从朝廷视角记录军政行动；出兵与战果需分辨。 这一处置留下的后续条件是：后续路线、控制与损失应与对方记录和遗址材料复核。

关于可说范围，须保留：《宋史》为元末官修，本纪保存宋廷纪年与处置；战果、数字、地方执行及对方动机须以编年、会要、对方史料或地方材料互校。 进一步讨论可参照 宋金战事、金代地方治理与碑志研究

\noindent\eventanchor{SYD-1164-JIN-220}{壬午，遣魏杞等爲金國通問使}\textbf{南宋与金。}
〔壬午，遣魏杞等爲金國通問使〕见于 《宋史》卷033〈隆興二年〉；《金史》本纪同年条（互校） 的 1164年（南宋隆興二年、金大定4年；公元换算据两国年号对照） 记录。南宋与金 的选择因此有了可回查的时间位置；材料未将地点细化到城址，不能用中枢所在地替它补足。

本纪记录使节、贡赐或往来，不等于稳定统属关系。 记录所见的结果则是：可与对方编年和交通、文书材料核对其后续落实。 日期相邻不等于因果已经证实。

证据的边界是：《宋史》为元末官修，本纪保存宋廷纪年与处置；战果、数字、地方执行及对方动机须以编年、会要、对方史料或地方材料互校。 宋金战事、金代地方治理与碑志研究 可用于继续检验这一结论。

\noindent\eventanchor{SYD-1164-JIN-248}{十一月乙酉，知楚州魏勝與金人戰，死之，州遂陷，濠州亦陷}\textbf{南宋与金。}
1164年（南宋隆興二年、金大定4年；公元换算据两国年号对照），《宋史》卷033〈隆興二年〉；《金史》本纪同年条（互校） 所见的〔十一月乙酉，知楚州魏勝與金人戰，死之，州遂陷，濠州亦陷〕把 南宋与金 的处置留在可复核的年序中；题名保留的城、路、水道或机构名称，是目前可以落地的空间线索。

把本项放回事件链，能够看到的前提为：本纪从朝廷视角记录军政行动；出兵与战果需分辨。 而它所打开或收紧的下一步是：后续路线、控制与损失应与对方记录和遗址材料复核。

《宋史》为元末官修，本纪保存宋廷纪年与处置；战果、数字、地方执行及对方动机须以编年、会要、对方史料或地方材料互校。 因此，宋金战事、金代地方治理与碑志研究 只能扩展问题，不能代替未保存的细节。

\subsection{1165年（南宋乾道元年、金大定5年；公元换算据两国年号对照）}
\noindent\eventanchor{SYD-1165-JIN-046}{癸未，遣<u>王曮</u>等賀金主生辰}\textbf{南宋与金。}
南宋与金 在 1165年（南宋乾道元年、金大定5年；公元换算据两国年号对照） 的材料痕迹是〔癸未，遣<u>王曮</u>等賀金主生辰〕，出处为 《宋史》卷033〈乾道元年〉；《金史》本纪同年条（互校）。材料未将地点细化到城址，不能用中枢所在地替它补足。

它承接的条件是：本纪年条提供日期和中枢可见行动；前因不据单条补定。 随后可追踪的变化为：可回查同年及后续条目，地方影响仍须另证。

关于可说范围，须保留：《宋史》为元末官修，本纪保存宋廷纪年与处置；战果、数字、地方执行及对方动机须以编年、会要、对方史料或地方材料互校。 进一步讨论可参照 宋金战事、金代地方治理与碑志研究

\subsection{1166年（南宋乾道二年、金大定6年；公元换算据两国年号对照）}
\noindent\eventanchor{SYD-1166-JIN-047}{乙亥，遣梁克家等賀金主生辰}\textbf{南宋与金。}
从 《宋史》卷033〈乾道二年〉；《金史》本纪同年条（互校） 的 1166年（南宋乾道二年、金大定6年；公元换算据两国年号对照） 条目看，〔乙亥，遣梁克家等賀金主生辰〕使 南宋与金 的一个具体行动进入叙事；材料未将地点细化到城址，不能用中枢所在地替它补足。

前一层压力可由以下线索辨认：本纪年条提供日期和中枢可见行动；前因不据单条补定。 这一处置留下的后续条件是：可回查同年及后续条目，地方影响仍须另证。

证据的边界是：《宋史》为元末官修，本纪保存宋廷纪年与处置；战果、数字、地方执行及对方动机须以编年、会要、对方史料或地方材料互校。 宋金战事、金代地方治理与碑志研究 可用于继续检验这一结论。

\subsection{1167年（南宋乾道三年、金大定7年；公元换算据两国年号对照）}
\noindent\eventanchor{SYD-1167-JIN-048}{己亥，遣王瀹賀金主生辰}\textbf{南宋与金。}
1167年（南宋乾道三年、金大定7年；公元换算据两国年号对照） 的记载将 南宋与金 与〔己亥，遣王瀹賀金主生辰〕相连。此处可确认的是这项被记录的行动；材料未将地点细化到城址，不能用中枢所在地替它补足。

本纪年条提供日期和中枢可见行动；前因不据单条补定。 记录所见的结果则是：可回查同年及后续条目，地方影响仍须另证。 日期相邻不等于因果已经证实。

《宋史》为元末官修，本纪保存宋廷纪年与处置；战果、数字、地方执行及对方动机须以编年、会要、对方史料或地方材料互校。 因此，宋金战事、金代地方治理与碑志研究 只能扩展问题，不能代替未保存的细节。

\subsection{1168年（南宋乾道四年、金大定8年；公元换算据两国年号对照）}
\noindent\eventanchor{SYD-1168-JIN-049}{甲午，禁歸正人藏匿金人者}\textbf{南宋与金。}
〔甲午，禁歸正人藏匿金人者〕见于 《宋史》卷034〈乾道四年〉；《金史》本纪同年条（互校） 的 1168年（南宋乾道四年、金大定8年；公元换算据两国年号对照） 记录。南宋与金 的选择因此有了可回查的时间位置；材料未将地点细化到城址，不能用中枢所在地替它补足。

把本项放回事件链，能够看到的前提为：本纪年条记录政令或机构处置；实施背景须参照制度材料。 而它所打开或收紧的下一步是：可在同年及后续政令中追踪；条文不单独证明地方执行。

关于可说范围，须保留：《宋史》为元末官修，本纪保存宋廷纪年与处置；战果、数字、地方执行及对方动机须以编年、会要、对方史料或地方材料互校。 进一步讨论可参照 宋金战事、金代地方治理与碑志研究

\noindent\eventanchor{SYD-1168-JIN-138}{十二月丙申，遣胡元質等賀金主生辰}\textbf{南宋与金。}
1168年（南宋乾道四年、金大定8年；公元换算据两国年号对照），《宋史》卷034〈乾道四年〉；《金史》本纪同年条（互校） 所见的〔十二月丙申，遣胡元質等賀金主生辰〕把 南宋与金 的处置留在可复核的年序中；材料未将地点细化到城址，不能用中枢所在地替它补足。

它承接的条件是：本纪年条提供日期和中枢可见行动；前因不据单条补定。 随后可追踪的变化为：可回查同年及后续条目，地方影响仍须另证。

证据的边界是：《宋史》为元末官修，本纪保存宋廷纪年与处置；战果、数字、地方执行及对方动机须以编年、会要、对方史料或地方材料互校。 宋金战事、金代地方治理与碑志研究 可用于继续检验这一结论。

\subsection{1169年（南宋乾道五年、金大定9年；公元换算据两国年号对照）}
\noindent\eventanchor{SYD-1169-JIN-050}{十二月己-{丑}-，遣司馬伋等賀金主生辰}\textbf{南宋与金。}
南宋与金 在 1169年（南宋乾道五年、金大定9年；公元换算据两国年号对照） 的材料痕迹是〔十二月己-{丑}-，遣司馬伋等賀金主生辰〕，出处为 《宋史》卷034〈乾道五年〉；《金史》本纪同年条（互校）。材料未将地点细化到城址，不能用中枢所在地替它补足。

前一层压力可由以下线索辨认：本纪年条提供日期和中枢可见行动；前因不据单条补定。 这一处置留下的后续条件是：可回查同年及后续条目，地方影响仍须另证。

《宋史》为元末官修，本纪保存宋廷纪年与处置；战果、数字、地方执行及对方动机须以编年、会要、对方史料或地方材料互校。 因此，宋金战事、金代地方治理与碑志研究 只能扩展问题，不能代替未保存的细节。

\noindent\eventanchor{SYD-1169-JIN-139}{先是，海州人時旺聚衆數千來請命，旺尋爲金人所獲，其徒渡淮而南者甚衆，故命滋彈壓之}\textbf{南宋与金。}
从 《宋史》卷034〈乾道五年〉；《金史》本纪同年条（互校） 的 1169年（南宋乾道五年、金大定9年；公元换算据两国年号对照） 条目看，〔先是，海州人時旺聚衆數千來請命，旺尋爲金人所獲，其徒渡淮而南者甚衆，故命滋彈壓之〕使 南宋与金 的一个具体行动进入叙事；题名保留的城、路、水道或机构名称，是目前可以落地的空间线索。

本纪年条记录政令或机构处置；实施背景须参照制度材料。 记录所见的结果则是：可在同年及后续政令中追踪；条文不单独证明地方执行。 日期相邻不等于因果已经证实。

关于可说范围，须保留：《宋史》为元末官修，本纪保存宋廷纪年与处置；战果、数字、地方执行及对方动机须以编年、会要、对方史料或地方材料互校。 进一步讨论可参照 宋金战事、金代地方治理与碑志研究

\subsection{1170年（南宋乾道六年、金大定10年；公元换算据两国年号对照）}
\noindent\eventanchor{SYD-1170-JIN-051}{是月，遣趙雄等賀金主生辰，別函書請更受書之禮}\textbf{南宋与金。}
1170年（南宋乾道六年、金大定10年；公元换算据两国年号对照） 的记载将 南宋与金 与〔是月，遣趙雄等賀金主生辰，別函書請更受書之禮〕相连。此处可确认的是这项被记录的行动；材料未将地点细化到城址，不能用中枢所在地替它补足。

把本项放回事件链，能够看到的前提为：本纪年条提供日期和中枢可见行动；前因不据单条补定。 而它所打开或收紧的下一步是：可回查同年及后续条目，地方影响仍须另证。

证据的边界是：《宋史》为元末官修，本纪保存宋廷纪年与处置；战果、数字、地方执行及对方动机须以编年、会要、对方史料或地方材料互校。 宋金战事、金代地方治理与碑志研究 可用于继续检验这一结论。

\subsection{1171年（南宋乾道七年、金大定11年；公元换算据两国年号对照）}
\noindent\eventanchor{SYD-1171-JIN-052}{壬戌，金遣烏林答天錫等來賀會慶節，天錫要帝降榻問金主起居，虞允文請帝還內，命知閣門事王抃諭天錫以明日見，天錫沮退}\textbf{南宋与金。}
〔壬戌，金遣烏林答天錫等來賀會慶節，天錫要帝降榻問金主起居，虞允文請帝還內，命知閣門事王抃諭天錫以明日見，天錫沮退〕见于 《宋史》卷034〈乾道七年〉；《金史》本纪同年条（互校） 的 1171年（南宋乾道七年、金大定11年；公元换算据两国年号对照） 记录。南宋与金 的选择因此有了可回查的时间位置；材料未将地点细化到城址，不能用中枢所在地替它补足。

它承接的条件是：本纪年条记录政令或机构处置；实施背景须参照制度材料。 随后可追踪的变化为：可在同年及后续政令中追踪；条文不单独证明地方执行。

《宋史》为元末官修，本纪保存宋廷纪年与处置；战果、数字、地方执行及对方动机须以编年、会要、对方史料或地方材料互校。 因此，宋金战事、金代地方治理与碑志研究 只能扩展问题，不能代替未保存的细节。

\noindent\eventanchor{SYD-1171-JIN-140}{十二月丁未，遣翟紱等賀金主生辰}\textbf{南宋与金。}
1171年（南宋乾道七年、金大定11年；公元换算据两国年号对照），《宋史》卷034〈乾道七年〉；《金史》本纪同年条（互校） 所见的〔十二月丁未，遣翟紱等賀金主生辰〕把 南宋与金 的处置留在可复核的年序中；材料未将地点细化到城址，不能用中枢所在地替它补足。

前一层压力可由以下线索辨认：本纪年条提供日期和中枢可见行动；前因不据单条补定。 这一处置留下的后续条件是：可回查同年及后续条目，地方影响仍须另证。

关于可说范围，须保留：《宋史》为元末官修，本纪保存宋廷纪年与处置；战果、数字、地方执行及对方动机须以编年、会要、对方史料或地方材料互校。 进一步讨论可参照 宋金战事、金代地方治理与碑志研究

\subsection{1172年（南宋乾道八年、金大定12年；公元换算据两国年号对照）}
\noindent\eventanchor{SYD-1172-JIN-053}{丁巳，遣韓元吉等賀金主生辰}\textbf{南宋与金。}
南宋与金 在 1172年（南宋乾道八年、金大定12年；公元换算据两国年号对照） 的材料痕迹是〔丁巳，遣韓元吉等賀金主生辰〕，出处为 《宋史》卷034〈乾道八年〉；《金史》本纪同年条（互校）。材料未将地点细化到城址，不能用中枢所在地替它补足。

本纪年条提供日期和中枢可见行动；前因不据单条补定。 记录所见的结果则是：可回查同年及后续条目，地方影响仍须另证。 日期相邻不等于因果已经证实。

证据的边界是：《宋史》为元末官修，本纪保存宋廷纪年与处置；战果、数字、地方执行及对方动机须以编年、会要、对方史料或地方材料互校。 宋金战事、金代地方治理与碑志研究 可用于继续检验这一结论。

\noindent\eventanchor{SYD-1172-JIN-141}{姚憲、曾覿至自金，金人拒其請}\textbf{南宋与金。}
从 《宋史》卷034〈乾道八年〉；《金史》本纪同年条（互校） 的 1172年（南宋乾道八年、金大定12年；公元换算据两国年号对照） 条目看，〔姚憲、曾覿至自金，金人拒其請〕使 南宋与金 的一个具体行动进入叙事；材料未将地点细化到城址，不能用中枢所在地替它补足。

把本项放回事件链，能够看到的前提为：本纪年条提供日期和中枢可见行动；前因不据单条补定。 而它所打开或收紧的下一步是：可回查同年及后续条目，地方影响仍须另证。

《宋史》为元末官修，本纪保存宋廷纪年与处置；战果、数字、地方执行及对方动机须以编年、会要、对方史料或地方材料互校。 因此，宋金战事、金代地方治理与碑志研究 只能扩展问题，不能代替未保存的细节。

\subsection{1173年（南宋乾道九年、金大定13年；公元换算据两国年号对照）}
\noindent\eventanchor{SYD-1173-JIN-054}{遣韓彥直等賀金主生辰}\textbf{南宋与金。}
1173年（南宋乾道九年、金大定13年；公元换算据两国年号对照） 的记载将 南宋与金 与〔遣韓彥直等賀金主生辰〕相连。此处可确认的是这项被记录的行动；材料未将地点细化到城址，不能用中枢所在地替它补足。

它承接的条件是：本纪年条提供日期和中枢可见行动；前因不据单条补定。 随后可追踪的变化为：可回查同年及后续条目，地方影响仍须另证。

关于可说范围，须保留：《宋史》为元末官修，本纪保存宋廷纪年与处置；战果、数字、地方执行及对方动机须以编年、会要、对方史料或地方材料互校。 进一步讨论可参照 宋金战事、金代地方治理与碑志研究

\subsection{1174年（南宋淳熙元年、金大定14年；公元换算据两国年号对照）}
\noindent\eventanchor{SYD-1174-JIN-055}{壬戌，遣吳琚等賀金主生辰}\textbf{南宋与金。}
〔壬戌，遣吳琚等賀金主生辰〕见于 《宋史》卷034〈淳熙元年〉；《金史》本纪同年条（互校） 的 1174年（南宋淳熙元年、金大定14年；公元换算据两国年号对照） 记录。南宋与金 的选择因此有了可回查的时间位置；材料未将地点细化到城址，不能用中枢所在地替它补足。

前一层压力可由以下线索辨认：本纪年条提供日期和中枢可见行动；前因不据单条补定。 这一处置留下的后续条件是：可回查同年及后续条目，地方影响仍须另证。

证据的边界是：《宋史》为元末官修，本纪保存宋廷纪年与处置；战果、数字、地方执行及对方动机须以编年、会要、对方史料或地方材料互校。 宋金战事、金代地方治理与碑志研究 可用于继续检验这一结论。

\subsection{1175年（南宋淳熙二年、金大定15年；公元换算据两国年号对照）}
\noindent\eventanchor{SYD-1175-JIN-056}{遣張宗元等賀金主生辰}\textbf{南宋与金。}
1175年（南宋淳熙二年、金大定15年；公元换算据两国年号对照），《宋史》卷034〈淳熙二年〉；《金史》本纪同年条（互校） 所见的〔遣張宗元等賀金主生辰〕把 南宋与金 的处置留在可复核的年序中；材料未将地点细化到城址，不能用中枢所在地替它补足。

本纪年条提供日期和中枢可见行动；前因不据单条补定。 记录所见的结果则是：可回查同年及后续条目，地方影响仍须另证。 日期相邻不等于因果已经证实。

《宋史》为元末官修，本纪保存宋廷纪年与处置；战果、数字、地方执行及对方动机须以编年、会要、对方史料或地方材料互校。 因此，宋金战事、金代地方治理与碑志研究 只能扩展问题，不能代替未保存的细节。

\subsection{1176年（南宋淳熙三年、金大定16年；公元换算据两国年号对照）}
\noindent\eventanchor{SYD-1176-JIN-057}{庚午，遣張子正等賀金主生辰}\textbf{南宋与金。}
南宋与金 在 1176年（南宋淳熙三年、金大定16年；公元换算据两国年号对照） 的材料痕迹是〔庚午，遣張子正等賀金主生辰〕，出处为 《宋史》卷034〈淳熙三年〉；《金史》本纪同年条（互校）。材料未将地点细化到城址，不能用中枢所在地替它补足。

把本项放回事件链，能够看到的前提为：本纪年条提供日期和中枢可见行动；前因不据单条补定。 而它所打开或收紧的下一步是：可回查同年及后续条目，地方影响仍须另证。

关于可说范围，须保留：《宋史》为元末官修，本纪保存宋廷纪年与处置；战果、数字、地方执行及对方动机须以编年、会要、对方史料或地方材料互校。 进一步讨论可参照 宋金战事、金代地方治理与碑志研究

\subsection{1177年（南宋淳熙四年、金大定17年；公元换算据两国年号对照）}
\noindent\eventanchor{SYD-1177-JIN-058}{癸亥，遣趙思等賀金主生辰}\textbf{南宋与金。}
从 《宋史》卷034〈淳熙四年〉；《金史》本纪同年条（互校） 的 1177年（南宋淳熙四年、金大定17年；公元换算据两国年号对照） 条目看，〔癸亥，遣趙思等賀金主生辰〕使 南宋与金 的一个具体行动进入叙事；材料未将地点细化到城址，不能用中枢所在地替它补足。

它承接的条件是：本纪年条提供日期和中枢可见行动；前因不据单条补定。 随后可追踪的变化为：可回查同年及后续条目，地方影响仍须另证。

证据的边界是：《宋史》为元末官修，本纪保存宋廷纪年与处置；战果、数字、地方执行及对方动机须以编年、会要、对方史料或地方材料互校。 宋金战事、金代地方治理与碑志研究 可用于继续检验这一结论。

\subsection{1178年（南宋淳熙五年、金大定18年；公元换算据两国年号对照）}
\noindent\eventanchor{SYD-1178-JIN-059}{辛卯，遣錢沖之等賀金主生辰}\textbf{南宋与金。}
1178年（南宋淳熙五年、金大定18年；公元换算据两国年号对照） 的记载将 南宋与金 与〔辛卯，遣錢沖之等賀金主生辰〕相连。此处可确认的是这项被记录的行动；材料未将地点细化到城址，不能用中枢所在地替它补足。

前一层压力可由以下线索辨认：本纪所记财用、赈济或征发属于特定报告场景。 这一处置留下的后续条件是：可与食货志、文书和物质材料比较，不外推为全境总量。

《宋史》为元末官修，本纪保存宋廷纪年与处置；战果、数字、地方执行及对方动机须以编年、会要、对方史料或地方材料互校。 因此，宋金战事、金代地方治理与碑志研究 只能扩展问题，不能代替未保存的细节。

\subsection{1179年（南宋淳熙六年、金大定19年；公元换算据两国年号对照）}
\noindent\eventanchor{SYD-1179-JIN-060}{癸未，遣傅淇等賀金主生辰}\textbf{南宋与金。}
〔癸未，遣傅淇等賀金主生辰〕见于 《宋史》卷035〈淳熙六年〉；《金史》本纪同年条（互校） 的 1179年（南宋淳熙六年、金大定19年；公元换算据两国年号对照） 记录。南宋与金 的选择因此有了可回查的时间位置；材料未将地点细化到城址，不能用中枢所在地替它补足。

本纪年条提供日期和中枢可见行动；前因不据单条补定。 记录所见的结果则是：可回查同年及后续条目，地方影响仍须另证。 日期相邻不等于因果已经证实。

关于可说范围，须保留：《宋史》为元末官修，本纪保存宋廷纪年与处置；战果、数字、地方执行及对方动机须以编年、会要、对方史料或地方材料互校。 进一步讨论可参照 宋金战事、金代地方治理与碑志研究

\subsection{1181年（南宋淳熙八年、金大定21年；公元换算据两国年号对照）}
\noindent\eventanchor{SYD-1181-JIN-061}{丁酉，遣燕世良賀金主生辰}\textbf{南宋与金。}
1181年（南宋淳熙八年、金大定21年；公元换算据两国年号对照），《宋史》卷035〈淳熙八年〉；《金史》本纪同年条（互校） 所见的〔丁酉，遣燕世良賀金主生辰〕把 南宋与金 的处置留在可复核的年序中；题名保留的城、路、水道或机构名称，是目前可以落地的空间线索。

把本项放回事件链，能够看到的前提为：本纪年条提供日期和中枢可见行动；前因不据单条补定。 而它所打开或收紧的下一步是：可回查同年及后续条目，地方影响仍须另证。

证据的边界是：《宋史》为元末官修，本纪保存宋廷纪年与处置；战果、数字、地方执行及对方动机须以编年、会要、对方史料或地方材料互校。 宋金战事、金代地方治理与碑志研究 可用于继续检验这一结论。

\subsection{1182年（南宋淳熙九年、金大定22年；公元换算据两国年号对照）}
\noindent\eventanchor{SYD-1182-JIN-062}{丙戌，遣賈選等賀金主生辰}\textbf{南宋与金。}
南宋与金 在 1182年（南宋淳熙九年、金大定22年；公元换算据两国年号对照） 的材料痕迹是〔丙戌，遣賈選等賀金主生辰〕，出处为 《宋史》卷035〈淳熙九年〉；《金史》本纪同年条（互校）。材料未将地点细化到城址，不能用中枢所在地替它补足。

它承接的条件是：本纪年条提供日期和中枢可见行动；前因不据单条补定。 随后可追踪的变化为：可回查同年及后续条目，地方影响仍须另证。

《宋史》为元末官修，本纪保存宋廷纪年与处置；战果、数字、地方执行及对方动机须以编年、会要、对方史料或地方材料互校。 因此，宋金战事、金代地方治理与碑志研究 只能扩展问题，不能代替未保存的细节。

\subsection{1183年（南宋淳熙十年、金大定23年；公元换算据两国年号对照）}
\noindent\eventanchor{SYD-1183-JIN-063}{丁巳，遣陳居仁等賀金主生辰}\textbf{南宋与金。}
从 《宋史》卷035〈淳熙十年〉；《金史》本纪同年条（互校） 的 1183年（南宋淳熙十年、金大定23年；公元换算据两国年号对照） 条目看，〔丁巳，遣陳居仁等賀金主生辰〕使 南宋与金 的一个具体行动进入叙事；材料未将地点细化到城址，不能用中枢所在地替它补足。

前一层压力可由以下线索辨认：本纪年条提供日期和中枢可见行动；前因不据单条补定。 这一处置留下的后续条件是：可回查同年及后续条目，地方影响仍须另证。

关于可说范围，须保留：《宋史》为元末官修，本纪保存宋廷纪年与处置；战果、数字、地方执行及对方动机须以编年、会要、对方史料或地方材料互校。 进一步讨论可参照 宋金战事、金代地方治理与碑志研究

\subsection{1184年（南宋淳熙十一年、金大定24年；公元换算据两国年号对照）}
\noindent\eventanchor{SYD-1184-JIN-064}{癸巳，命利路三都統吳挺、郭鈞、彭杲密陳出師進取利害，以備金人}\textbf{南宋与金。}
1184年（南宋淳熙十一年、金大定24年；公元换算据两国年号对照） 的记载将 南宋与金 与〔癸巳，命利路三都統吳挺、郭鈞、彭杲密陳出師進取利害，以備金人〕相连。此处可确认的是这项被记录的行动；题名保留的城、路、水道或机构名称，是目前可以落地的空间线索。

本纪年条记录政令或机构处置；实施背景须参照制度材料。 记录所见的结果则是：可在同年及后续政令中追踪；条文不单独证明地方执行。 日期相邻不等于因果已经证实。

证据的边界是：《宋史》为元末官修，本纪保存宋廷纪年与处置；战果、数字、地方执行及对方动机须以编年、会要、对方史料或地方材料互校。 宋金战事、金代地方治理与碑志研究 可用于继续检验这一结论。

\noindent\eventanchor{SYD-1184-JIN-142}{盱眙軍言得金人牒，以上京地寒，來歲正旦、生辰人使權止一年}\textbf{南宋与金。}
〔盱眙軍言得金人牒，以上京地寒，來歲正旦、生辰人使權止一年〕见于 《宋史》卷035〈淳熙十一年〉；《金史》本纪同年条（互校） 的 1184年（南宋淳熙十一年、金大定24年；公元换算据两国年号对照） 记录。南宋与金 的选择因此有了可回查的时间位置；题名保留的城、路、水道或机构名称，是目前可以落地的空间线索。

把本项放回事件链，能够看到的前提为：本纪年条提供日期和中枢可见行动；前因不据单条补定。 而它所打开或收紧的下一步是：可回查同年及后续条目，地方影响仍须另证。

《宋史》为元末官修，本纪保存宋廷纪年与处置；战果、数字、地方执行及对方动机须以编年、会要、对方史料或地方材料互校。 因此，宋金战事、金代地方治理与碑志研究 只能扩展问题，不能代替未保存的细节。

\noindent\eventanchor{SYD-1184-JIN-185}{乙丑，遣王信等賀金主生辰}\textbf{南宋与金。}
1184年（南宋淳熙十一年、金大定24年；公元换算据两国年号对照），《宋史》卷035〈淳熙十一年〉；《金史》本纪同年条（互校） 所见的〔乙丑，遣王信等賀金主生辰〕把 南宋与金 的处置留在可复核的年序中；材料未将地点细化到城址，不能用中枢所在地替它补足。

它承接的条件是：本纪年条提供日期和中枢可见行动；前因不据单条补定。 随后可追踪的变化为：可回查同年及后续条目，地方影响仍须另证。

关于可说范围，须保留：《宋史》为元末官修，本纪保存宋廷纪年与处置；战果、数字、地方执行及对方动机须以编年、会要、对方史料或地方材料互校。 进一步讨论可参照 宋金战事、金代地方治理与碑志研究

\subsection{1185年（南宋淳熙十二年、金大定25年；公元换算据两国年号对照）}
\noindent\eventanchor{SYD-1185-JIN-065}{壬辰，遣章森等賀金主生辰}\textbf{南宋与金。}
南宋与金 在 1185年（南宋淳熙十二年、金大定25年；公元换算据两国年号对照） 的材料痕迹是〔壬辰，遣章森等賀金主生辰〕，出处为 《宋史》卷035〈淳熙十二年〉；《金史》本纪同年条（互校）。材料未将地点细化到城址，不能用中枢所在地替它补足。

前一层压力可由以下线索辨认：本纪年条提供日期和中枢可见行动；前因不据单条补定。 这一处置留下的后续条件是：可回查同年及后续条目，地方影响仍须另证。

证据的边界是：《宋史》为元末官修，本纪保存宋廷纪年与处置；战果、数字、地方执行及对方动机须以编年、会要、对方史料或地方材料互校。 宋金战事、金代地方治理与碑志研究 可用于继续检验这一结论。

\subsection{1186年（南宋淳熙十三年、金大定26年；公元换算据两国年号对照）}
\noindent\eventanchor{SYD-1186-JIN-066}{庚申，遣張叔椿等賀金主生辰}\textbf{南宋与金。}
从 《宋史》卷035〈淳熙十三年〉；《金史》本纪同年条（互校） 的 1186年（南宋淳熙十三年、金大定26年；公元换算据两国年号对照） 条目看，〔庚申，遣張叔椿等賀金主生辰〕使 南宋与金 的一个具体行动进入叙事；材料未将地点细化到城址，不能用中枢所在地替它补足。

本纪年条提供日期和中枢可见行动；前因不据单条补定。 记录所见的结果则是：可回查同年及后续条目，地方影响仍须另证。 日期相邻不等于因果已经证实。

《宋史》为元末官修，本纪保存宋廷纪年与处置；战果、数字、地方执行及对方动机须以编年、会要、对方史料或地方材料互校。 因此，宋金战事、金代地方治理与碑志研究 只能扩展问题，不能代替未保存的细节。

\subsection{1187年（南宋淳熙十四年、金大定27年；公元换算据两国年号对照）}
\noindent\eventanchor{SYD-1187-JIN-067}{遣胡晉臣等賀金主生辰}\textbf{南宋与金。}
1187年（南宋淳熙十四年、金大定27年；公元换算据两国年号对照） 的记载将 南宋与金 与〔遣胡晉臣等賀金主生辰〕相连。此处可确认的是这项被记录的行动；材料未将地点细化到城址，不能用中枢所在地替它补足。

把本项放回事件链，能够看到的前提为：本纪年条提供日期和中枢可见行动；前因不据单条补定。 而它所打开或收紧的下一步是：可回查同年及后续条目，地方影响仍须另证。

关于可说范围，须保留：《宋史》为元末官修，本纪保存宋廷纪年与处置；战果、数字、地方执行及对方动机须以编年、会要、对方史料或地方材料互校。 进一步讨论可参照 宋金战事、金代地方治理与碑志研究

\noindent\eventanchor{SYD-1187-JIN-143}{以顏師魯等充金國遣留國信使}\textbf{南宋与金。}
〔以顏師魯等充金國遣留國信使〕见于 《宋史》卷035〈淳熙十四年〉；《金史》本纪同年条（互校） 的 1187年（南宋淳熙十四年、金大定27年；公元换算据两国年号对照） 记录。南宋与金 的选择因此有了可回查的时间位置；材料未将地点细化到城址，不能用中枢所在地替它补足。

它承接的条件是：本纪记录使节、贡赐或往来，不等于稳定统属关系。 随后可追踪的变化为：可与对方编年和交通、文书材料核对其后续落实。

证据的边界是：《宋史》为元末官修，本纪保存宋廷纪年与处置；战果、数字、地方执行及对方动机须以编年、会要、对方史料或地方材料互校。 宋金战事、金代地方治理与碑志研究 可用于继续检验这一结论。

\subsection{1189年}
\noindent\eventanchor{JIN-1189-EVENT-099}{金世宗去世，章宗即位}\textbf{金。}
1189年，《金史》相关本纪；《宋史》或《三朝北盟会编》《建炎以来系年要录》对应年条；金代碑志与蒙古、高丽材料（据事核） 所见的〔金世宗去世，章宗即位〕把 金 的处置留在可复核的年序中；材料未将地点细化到城址，不能用中枢所在地替它补足。

前一层压力可由以下线索辨认：金世宗至章宗时期的制度与边境压力中前任统治的结束与宫廷、部族或官僚的权力安排相互交织，促成“金世宗去世，章宗即位”这一继承节点。 这一处置留下的后续条件是：继承并不自动消除既有矛盾；“金世宗去世，章宗即位”后的任官、军权和对外策略需由后续条目追踪。

《金史》为元修，宋金战报和官修合法性语言各有宣传性；军数、俘获、控制范围和地方执行须以对方及物质材料限缩。 因此，金代国家形成、宋金边境、都城与金蒙转换研究 只能扩展问题，不能代替未保存的细节。

\subsection{1189年（南宋淳熙十六年、金大定29年；公元换算据两国年号对照）}
\noindent\eventanchor{SYD-1189-JIN-068}{十六年春正月癸巳，金主雍殂，孫璟立}\textbf{南宋与金。}
南宋与金 在 1189年（南宋淳熙十六年、金大定29年；公元换算据两国年号对照） 的材料痕迹是〔十六年春正月癸巳，金主雍殂，孫璟立〕，出处为 《宋史》卷035〈淳熙十六年〉；《金史》本纪同年条（互校）。材料未将地点细化到城址，不能用中枢所在地替它补足。

本纪年条提供日期和中枢可见行动；前因不据单条补定。 记录所见的结果则是：可回查同年及后续条目，地方影响仍须另证。 日期相邻不等于因果已经证实。

关于可说范围，须保留：《宋史》为元末官修，本纪保存宋廷纪年与处置；战果、数字、地方执行及对方动机须以编年、会要、对方史料或地方材料互校。 进一步讨论可参照 宋金战事、金代地方治理与碑志研究

\noindent\eventanchor{SYD-1189-JIN-144}{戊辰，遣謝深甫等賀金主生辰}\textbf{南宋与金。}
从 《宋史》卷036〈淳熙十六年〉；《金史》本纪同年条（互校） 的 1189年（南宋淳熙十六年、金大定29年；公元换算据两国年号对照） 条目看，〔戊辰，遣謝深甫等賀金主生辰〕使 南宋与金 的一个具体行动进入叙事；材料未将地点细化到城址，不能用中枢所在地替它补足。

把本项放回事件链，能够看到的前提为：本纪年条提供日期和中枢可见行动；前因不据单条补定。 而它所打开或收紧的下一步是：可回查同年及后续条目，地方影响仍须另证。

证据的边界是：《宋史》为元末官修，本纪保存宋廷纪年与处置；战果、数字、地方执行及对方动机须以编年、会要、对方史料或地方材料互校。 宋金战事、金代地方治理与碑志研究 可用于继续检验这一结论。

\subsection{1190年（南宋紹熙元年、金明昌1年；公元换算据两国年号对照）}
\noindent\eventanchor{SYD-1190-JIN-069}{六月丁亥，遣丘崈等賀金主生辰}\textbf{南宋与金。}
1190年（南宋紹熙元年、金明昌1年；公元换算据两国年号对照） 的记载将 南宋与金 与〔六月丁亥，遣丘崈等賀金主生辰〕相连。此处可确认的是这项被记录的行动；材料未将地点细化到城址，不能用中枢所在地替它补足。

它承接的条件是：本纪年条提供日期和中枢可见行动；前因不据单条补定。 随后可追踪的变化为：可回查同年及后续条目，地方影响仍须另证。

《宋史》为元末官修，本纪保存宋廷纪年与处置；战果、数字、地方执行及对方动机须以编年、会要、对方史料或地方材料互校。 因此，宋金战事、金代地方治理与碑志研究 只能扩展问题，不能代替未保存的细节。

\subsection{1191年（南宋紹熙二年、金明昌2年；公元换算据两国年号对照）}
\noindent\eventanchor{SYD-1191-JIN-070}{庚辰，遣趙廱等賀金主生辰}\textbf{南宋与金。}
〔庚辰，遣趙廱等賀金主生辰〕见于 《宋史》卷036〈紹熙二年〉；《金史》本纪同年条（互校） 的 1191年（南宋紹熙二年、金明昌2年；公元换算据两国年号对照） 记录。南宋与金 的选择因此有了可回查的时间位置；材料未将地点细化到城址，不能用中枢所在地替它补足。

前一层压力可由以下线索辨认：本纪年条提供日期和中枢可见行动；前因不据单条补定。 这一处置留下的后续条件是：可回查同年及后续条目，地方影响仍须另证。

关于可说范围，须保留：《宋史》为元末官修，本纪保存宋廷纪年与处置；战果、数字、地方执行及对方动机须以编年、会要、对方史料或地方材料互校。 进一步讨论可参照 宋金战事、金代地方治理与碑志研究

\subsection{1192年（南宋紹熙三年、金明昌3年；公元换算据两国年号对照）}
\noindent\eventanchor{SYD-1192-JIN-071}{以權禮部尚書陳騤同知樞密院事，甲辰，遣錢之望等賀金主生辰}\textbf{南宋与金。}
1192年（南宋紹熙三年、金明昌3年；公元换算据两国年号对照），《宋史》卷036〈紹熙三年〉；《金史》本纪同年条（互校） 所见的〔以權禮部尚書陳騤同知樞密院事，甲辰，遣錢之望等賀金主生辰〕把 南宋与金 的处置留在可复核的年序中；材料未将地点细化到城址，不能用中枢所在地替它补足。

本纪所记财用、赈济或征发属于特定报告场景。 记录所见的结果则是：可与食货志、文书和物质材料比较，不外推为全境总量。 日期相邻不等于因果已经证实。

证据的边界是：《宋史》为元末官修，本纪保存宋廷纪年与处置；战果、数字、地方执行及对方动机须以编年、会要、对方史料或地方材料互校。 宋金战事、金代地方治理与碑志研究 可用于继续检验这一结论。

\subsection{1193年（南宋紹熙四年、金明昌4年；公元换算据两国年号对照）}
\noindent\eventanchor{SYD-1193-JIN-072}{己亥，遣許及之等賀金主生辰}\textbf{南宋与金。}
南宋与金 在 1193年（南宋紹熙四年、金明昌4年；公元换算据两国年号对照） 的材料痕迹是〔己亥，遣許及之等賀金主生辰〕，出处为 《宋史》卷036〈紹熙四年〉；《金史》本纪同年条（互校）。材料未将地点细化到城址，不能用中枢所在地替它补足。

把本项放回事件链，能够看到的前提为：本纪年条提供日期和中枢可见行动；前因不据单条补定。 而它所打开或收紧的下一步是：可回查同年及后续条目，地方影响仍须另证。

《宋史》为元末官修，本纪保存宋廷纪年与处置；战果、数字、地方执行及对方动机须以编年、会要、对方史料或地方材料互校。 因此，宋金战事、金代地方治理与碑志研究 只能扩展问题，不能代替未保存的细节。

\subsection{1194年（南宋紹熙五年、金明昌5年；公元换算据两国年号对照）}
\noindent\eventanchor{SYD-1194-JIN-073}{六月，遣梁總等賀金主生辰}\textbf{南宋与金。}
从 《宋史》卷036〈紹熙五年〉；《金史》本纪同年条（互校） 的 1194年（南宋紹熙五年、金明昌5年；公元换算据两国年号对照） 条目看，〔六月，遣梁總等賀金主生辰〕使 南宋与金 的一个具体行动进入叙事；材料未将地点细化到城址，不能用中枢所在地替它补足。

它承接的条件是：本纪年条提供日期和中枢可见行动；前因不据单条补定。 随后可追踪的变化为：可回查同年及后续条目，地方影响仍须另证。

关于可说范围，须保留：《宋史》为元末官修，本纪保存宋廷纪年与处置；战果、数字、地方执行及对方动机须以编年、会要、对方史料或地方材料互校。 进一步讨论可参照 宋金战事、金代地方治理与碑志研究

\subsection{1195年（南宋慶元元年、金明昌6年；公元换算据两国年号对照）}
\noindent\eventanchor{SYD-1195-JIN-074}{己未，遣汪義端賀金主生辰}\textbf{南宋与金。}
1195年（南宋慶元元年、金明昌6年；公元换算据两国年号对照） 的记载将 南宋与金 与〔己未，遣汪義端賀金主生辰〕相连。此处可确认的是这项被记录的行动；材料未将地点细化到城址，不能用中枢所在地替它补足。

前一层压力可由以下线索辨认：本纪年条提供日期和中枢可见行动；前因不据单条补定。 这一处置留下的后续条件是：可回查同年及后续条目，地方影响仍须另证。

证据的边界是：《宋史》为元末官修，本纪保存宋廷纪年与处置；战果、数字、地方执行及对方动机须以编年、会要、对方史料或地方材料互校。 宋金战事、金代地方治理与碑志研究 可用于继续检验这一结论。

\subsection{1196年（南宋慶元二年、金承安1年；公元换算据两国年号对照）}
\noindent\eventanchor{SYD-1196-JIN-075}{六月庚戌，遣吳宗旦賀金主生辰，乙丑，命監司、帥守臧否縣令，分三等，丙子，子埈生}\textbf{南宋与金。}
〔六月庚戌，遣吳宗旦賀金主生辰，乙丑，命監司、帥守臧否縣令，分三等，丙子，子埈生〕见于 《宋史》卷037〈慶元二年〉；《金史》本纪同年条（互校） 的 1196年（南宋慶元二年、金承安1年；公元换算据两国年号对照） 记录。南宋与金 的选择因此有了可回查的时间位置；材料未将地点细化到城址，不能用中枢所在地替它补足。

本纪年条记录政令或机构处置；实施背景须参照制度材料。 记录所见的结果则是：可在同年及后续政令中追踪；条文不单独证明地方执行。 日期相邻不等于因果已经证实。

《宋史》为元末官修，本纪保存宋廷纪年与处置；战果、数字、地方执行及对方动机须以编年、会要、对方史料或地方材料互校。 因此，宋金战事、金代地方治理与碑志研究 只能扩展问题，不能代替未保存的细节。

\subsection{1197年（南宋慶元三年、金承安2年；公元换算据两国年号对照）}
\noindent\eventanchor{SYD-1197-JIN-076}{乙亥，遣衛涇賀金主生辰}\textbf{南宋与金。}
1197年（南宋慶元三年、金承安2年；公元换算据两国年号对照），《宋史》卷037〈慶元三年〉；《金史》本纪同年条（互校） 所见的〔乙亥，遣衛涇賀金主生辰〕把 南宋与金 的处置留在可复核的年序中；材料未将地点细化到城址，不能用中枢所在地替它补足。

把本项放回事件链，能够看到的前提为：本纪年条提供日期和中枢可见行动；前因不据单条补定。 而它所打开或收紧的下一步是：可回查同年及后续条目，地方影响仍须另证。

关于可说范围，须保留：《宋史》为元末官修，本纪保存宋廷纪年与处置；战果、数字、地方执行及对方动机须以编年、会要、对方史料或地方材料互校。 进一步讨论可参照 宋金战事、金代地方治理与碑志研究

\subsection{1198年（南宋慶元四年、金承安3年；公元换算据两国年号对照）}
\noindent\eventanchor{SYD-1198-JIN-077}{六月己巳，遣楊王休賀金主生辰}\textbf{南宋与金。}
南宋与金 在 1198年（南宋慶元四年、金承安3年；公元换算据两国年号对照） 的材料痕迹是〔六月己巳，遣楊王休賀金主生辰〕，出处为 《宋史》卷037〈慶元四年〉；《金史》本纪同年条（互校）。材料未将地点细化到城址，不能用中枢所在地替它补足。

它承接的条件是：本纪年条提供日期和中枢可见行动；前因不据单条补定。 随后可追踪的变化为：可回查同年及后续条目，地方影响仍须另证。

证据的边界是：《宋史》为元末官修，本纪保存宋廷纪年与处置；战果、数字、地方执行及对方动机须以编年、会要、对方史料或地方材料互校。 宋金战事、金代地方治理与碑志研究 可用于继续检验这一结论。

\subsection{1199年（南宋慶元五年、金承安4年；公元换算据两国年号对照）}
\noindent\eventanchor{SYD-1199-JIN-078}{六月癸亥，遣李大性賀金主生辰}\textbf{南宋与金。}
从 《宋史》卷037〈慶元五年〉；《金史》本纪同年条（互校） 的 1199年（南宋慶元五年、金承安4年；公元换算据两国年号对照） 条目看，〔六月癸亥，遣李大性賀金主生辰〕使 南宋与金 的一个具体行动进入叙事；材料未将地点细化到城址，不能用中枢所在地替它补足。

前一层压力可由以下线索辨认：本纪年条提供日期和中枢可见行动；前因不据单条补定。 这一处置留下的后续条件是：可回查同年及后续条目，地方影响仍须另证。

《宋史》为元末官修，本纪保存宋廷纪年与处置；战果、数字、地方执行及对方动机须以编年、会要、对方史料或地方材料互校。 因此，宋金战事、金代地方治理与碑志研究 只能扩展问题，不能代替未保存的细节。

\subsection{1200年（南宋慶元六年、金承安5年；公元换算据两国年号对照）}
\noindent\eventanchor{SYD-1200-JIN-079}{丙子，遣丁常任爲金國遺留國信使}\textbf{南宋与金。}
1200年（南宋慶元六年、金承安5年；公元换算据两国年号对照） 的记载将 南宋与金 与〔丙子，遣丁常任爲金國遺留國信使〕相连。此处可确认的是这项被记录的行动；材料未将地点细化到城址，不能用中枢所在地替它补足。

本纪记录使节、贡赐或往来，不等于稳定统属关系。 记录所见的结果则是：可与对方编年和交通、文书材料核对其后续落实。 日期相邻不等于因果已经证实。

关于可说范围，须保留：《宋史》为元末官修，本纪保存宋廷纪年与处置；战果、数字、地方执行及对方动机须以编年、会要、对方史料或地方材料互校。 进一步讨论可参照 宋金战事、金代地方治理与碑志研究

\noindent\eventanchor{SYD-1200-JIN-145}{壬辰，遣趙善義賀金主生辰，吳旴使金告哀}\textbf{南宋与金。}
〔壬辰，遣趙善義賀金主生辰，吳旴使金告哀〕见于 《宋史》卷037〈慶元六年〉；《金史》本纪同年条（互校） 的 1200年（南宋慶元六年、金承安5年；公元换算据两国年号对照） 记录。南宋与金 的选择因此有了可回查的时间位置；材料未将地点细化到城址，不能用中枢所在地替它补足。

把本项放回事件链，能够看到的前提为：本纪记录使节、贡赐或往来，不等于稳定统属关系。 而它所打开或收紧的下一步是：可与对方编年和交通、文书材料核对其后续落实。

证据的边界是：《宋史》为元末官修，本纪保存宋廷纪年与处置；战果、数字、地方执行及对方动机须以编年、会要、对方史料或地方材料互校。 宋金战事、金代地方治理与碑志研究 可用于继续检验这一结论。

\subsection{1201年（南宋嘉泰元年、金泰和1年；公元换算据两国年号对照）}
\noindent\eventanchor{SYD-1201-JIN-080}{六月辛巳，遣陳宗召賀金主生辰}\textbf{南宋与金。}
1201年（南宋嘉泰元年、金泰和1年；公元换算据两国年号对照），《宋史》卷038〈嘉泰元年〉；《金史》本纪同年条（互校） 所见的〔六月辛巳，遣陳宗召賀金主生辰〕把 南宋与金 的处置留在可复核的年序中；材料未将地点细化到城址，不能用中枢所在地替它补足。

它承接的条件是：本纪年条提供日期和中枢可见行动；前因不据单条补定。 随后可追踪的变化为：可回查同年及后续条目，地方影响仍须另证。

《宋史》为元末官修，本纪保存宋廷纪年与处置；战果、数字、地方执行及对方动机须以编年、会要、对方史料或地方材料互校。 因此，宋金战事、金代地方治理与碑志研究 只能扩展问题，不能代替未保存的细节。

\subsection{1202年（南宋嘉泰二年、金泰和2年；公元换算据两国年号对照）}
\noindent\eventanchor{SYD-1202-JIN-081}{六月丙子，遣趙不艱賀金主生辰}\textbf{南宋与金。}
南宋与金 在 1202年（南宋嘉泰二年、金泰和2年；公元换算据两国年号对照） 的材料痕迹是〔六月丙子，遣趙不艱賀金主生辰〕，出处为 《宋史》卷038〈嘉泰二年〉；《金史》本纪同年条（互校）。材料未将地点细化到城址，不能用中枢所在地替它补足。

前一层压力可由以下线索辨认：本纪年条提供日期和中枢可见行动；前因不据单条补定。 这一处置留下的后续条件是：可回查同年及后续条目，地方影响仍须另证。

关于可说范围，须保留：《宋史》为元末官修，本纪保存宋廷纪年与处置；战果、数字、地方执行及对方动机须以编年、会要、对方史料或地方材料互校。 进一步讨论可参照 宋金战事、金代地方治理与碑志研究

\subsection{1203年（南宋嘉泰三年、金泰和3年；公元换算据两国年号对照）}
\noindent\eventanchor{SYD-1203-JIN-082}{六月壬寅，遣劉甲賀金主生辰}\textbf{南宋与金。}
从 《宋史》卷038〈嘉泰三年〉；《金史》本纪同年条（互校） 的 1203年（南宋嘉泰三年、金泰和3年；公元换算据两国年号对照） 条目看，〔六月壬寅，遣劉甲賀金主生辰〕使 南宋与金 的一个具体行动进入叙事；材料未将地点细化到城址，不能用中枢所在地替它补足。

本纪年条提供日期和中枢可见行动；前因不据单条补定。 记录所见的结果则是：可回查同年及后续条目，地方影响仍须另证。 日期相邻不等于因果已经证实。

证据的边界是：《宋史》为元末官修，本纪保存宋廷纪年与处置；战果、数字、地方执行及对方动机须以编年、会要、对方史料或地方材料互校。 宋金战事、金代地方治理与碑志研究 可用于继续检验这一结论。

\noindent\eventanchor{SYD-1203-JIN-146}{是冬，金國多難，懼朝廷乘其隙，沿邊聚糧增戍，且禁襄陽榷場}\textbf{南宋与金。}
1203年（南宋嘉泰三年、金泰和3年；公元换算据两国年号对照） 的记载将 南宋与金 与〔是冬，金國多難，懼朝廷乘其隙，沿邊聚糧增戍，且禁襄陽榷場〕相连。此处可确认的是这项被记录的行动；材料未将地点细化到城址，不能用中枢所在地替它补足。

把本项放回事件链，能够看到的前提为：本纪年条记录政令或机构处置；实施背景须参照制度材料。 而它所打开或收紧的下一步是：可在同年及后续政令中追踪；条文不单独证明地方执行。

《宋史》为元末官修，本纪保存宋廷纪年与处置；战果、数字、地方执行及对方动机须以编年、会要、对方史料或地方材料互校。 因此，宋金战事、金代地方治理与碑志研究 只能扩展问题，不能代替未保存的细节。

\subsection{1204年（南宋嘉泰四年、金泰和4年；公元换算据两国年号对照）}
\noindent\eventanchor{SYD-1204-JIN-083}{六月癸巳，遣張嗣古賀金主生辰}\textbf{南宋与金。}
〔六月癸巳，遣張嗣古賀金主生辰〕见于 《宋史》卷038〈嘉泰四年〉；《金史》本纪同年条（互校） 的 1204年（南宋嘉泰四年、金泰和4年；公元换算据两国年号对照） 记录。南宋与金 的选择因此有了可回查的时间位置；材料未将地点细化到城址，不能用中枢所在地替它补足。

它承接的条件是：本纪年条提供日期和中枢可见行动；前因不据单条补定。 随后可追踪的变化为：可回查同年及后续条目，地方影响仍须另证。

关于可说范围，须保留：《宋史》为元末官修，本纪保存宋廷纪年与处置；战果、数字、地方执行及对方动机须以编年、会要、对方史料或地方材料互校。 进一步讨论可参照 宋金战事、金代地方治理与碑志研究

\subsection{1205年（南宋開禧元年、金泰和5年；公元换算据两国年号对照）}
\noindent\eventanchor{SYD-1205-JIN-084}{己亥，遣李壁賀金主生辰}\textbf{南宋与金。}
1205年（南宋開禧元年、金泰和5年；公元换算据两国年号对照），《宋史》卷038〈開禧元年〉；《金史》本纪同年条（互校） 所见的〔己亥，遣李壁賀金主生辰〕把 南宋与金 的处置留在可复核的年序中；材料未将地点细化到城址，不能用中枢所在地替它补足。

前一层压力可由以下线索辨认：本纪年条提供日期和中枢可见行动；前因不据单条补定。 这一处置留下的后续条件是：可回查同年及后续条目，地方影响仍须另证。

证据的边界是：《宋史》为元末官修，本纪保存宋廷纪年与处置；战果、数字、地方执行及对方动机须以编年、会要、对方史料或地方材料互校。 宋金战事、金代地方治理与碑志研究 可用于继续检验这一结论。

\noindent\eventanchor{SYD-1205-JIN-147}{是月，金國以邊民侵掠及增邊戍來責渝盟}\textbf{南宋与金。}
南宋与金 在 1205年（南宋開禧元年、金泰和5年；公元换算据两国年号对照） 的材料痕迹是〔是月，金國以邊民侵掠及增邊戍來責渝盟〕，出处为 《宋史》卷038〈開禧元年〉；《金史》本纪同年条（互校）。材料未将地点细化到城址，不能用中枢所在地替它补足。

本纪记录使节、贡赐或往来，不等于稳定统属关系。 记录所见的结果则是：可与对方编年和交通、文书材料核对其后续落实。 日期相邻不等于因果已经证实。

《宋史》为元末官修，本纪保存宋廷纪年与处置；战果、数字、地方执行及对方动机须以编年、会要、对方史料或地方材料互校。 因此，宋金战事、金代地方治理与碑志研究 只能扩展问题，不能代替未保存的细节。

\subsection{1206年（金泰和六年、宋开禧二年；干支、月日待核；公元换算据双方本纪年号对照）}
\noindent\eventanchor{JIN-1206-EVENT-100}{金军进攻南宋，南宋北伐受挫}\textbf{金与南宋。}
从 《金史》相关本纪；《宋史》或《三朝北盟会编》《建炎以来系年要录》对应年条；金代碑志与蒙古、高丽材料（据事核） 的 1206年（金泰和六年、宋开禧二年；干支、月日待核；公元换算据双方本纪年号对照） 条目看，〔金军进攻南宋，南宋北伐受挫〕使 金与南宋 的一个具体行动进入叙事；材料未将地点细化到城址，不能用中枢所在地替它补足。

把本项放回事件链，能够看到的前提为：金世宗至章宗时期的制度与边境压力中既有的联盟、交通与补给压力，为“金军进攻南宋，南宋北伐受挫”提供了直接的军事或动员背景。 而它所打开或收紧的下一步是：“金军进攻南宋，南宋北伐受挫”改变了相关城镇、道路或政权间的军事选择；伤亡、俘获和控制范围仅可按各方材料分别确认。

关于可说范围，须保留：《金史》为元修，宋金战报和官修合法性语言各有宣传性；军数、俘获、控制范围和地方执行须以对方及物质材料限缩。 进一步讨论可参照 金代国家形成、宋金边境、都城与金蒙转换研究

\subsection{1206年（南宋開禧二年、金泰和6年；公元换算据两国年号对照）}
\noindent\eventanchor{SYD-1206-JIN-085}{庚戌，金人破成州，守臣辛槱之遁去}\textbf{南宋与金。}
1206年（南宋開禧二年、金泰和6年；公元换算据两国年号对照） 的记载将 南宋与金 与〔庚戌，金人破成州，守臣辛槱之遁去〕相连。此处可确认的是这项被记录的行动；题名保留的城、路、水道或机构名称，是目前可以落地的空间线索。

它承接的条件是：本纪从朝廷视角记录军政行动；出兵与战果需分辨。 随后可追踪的变化为：后续路线、控制与损失应与对方记录和遗址材料复核。

证据的边界是：《宋史》为元末官修，本纪保存宋廷纪年与处置；战果、数字、地方执行及对方动机须以编年、会要、对方史料或地方材料互校。 宋金战事、金代地方治理与碑志研究 可用于继续检验这一结论。

\noindent\eventanchor{SYD-1206-JIN-148}{癸卯，郭倬等還至蘄縣，金人追而圍之，倬執馬軍司統制田俊邁以與金人，乃得免}\textbf{南宋与金。}
〔癸卯，郭倬等還至蘄縣，金人追而圍之，倬執馬軍司統制田俊邁以與金人，乃得免〕见于 《宋史》卷038〈開禧二年〉；《金史》本纪同年条（互校） 的 1206年（南宋開禧二年、金泰和6年；公元换算据两国年号对照） 记录。南宋与金 的选择因此有了可回查的时间位置；材料未将地点细化到城址，不能用中枢所在地替它补足。

前一层压力可由以下线索辨认：本纪从朝廷视角记录军政行动；出兵与战果需分辨。 这一处置留下的后续条件是：后续路线、控制与损失应与对方记录和遗址材料复核。

《宋史》为元末官修，本纪保存宋廷纪年与处置；战果、数字、地方执行及对方动机须以编年、会要、对方史料或地方材料互校。 因此，宋金战事、金代地方治理与碑志研究 只能扩展问题，不能代替未保存的细节。

\noindent\eventanchor{SYD-1206-JIN-186}{甲寅，金人攻六合縣，郭倪遣前軍統制郭僎救之，遇於胥浦橋，大敗，倪棄揚州走}\textbf{南宋与金。}
1206年（南宋開禧二年、金泰和6年；公元换算据两国年号对照），《宋史》卷038〈開禧二年〉；《金史》本纪同年条（互校） 所见的〔甲寅，金人攻六合縣，郭倪遣前軍統制郭僎救之，遇於胥浦橋，大敗，倪棄揚州走〕把 南宋与金 的处置留在可复核的年序中；题名保留的城、路、水道或机构名称，是目前可以落地的空间线索。

本纪从朝廷视角记录军政行动；出兵与战果需分辨。 记录所见的结果则是：后续路线、控制与损失应与对方记录和遗址材料复核。 日期相邻不等于因果已经证实。

关于可说范围，须保留：《宋史》为元末官修，本纪保存宋廷纪年与处置；战果、数字、地方执行及对方动机须以编年、会要、对方史料或地方材料互校。 进一步讨论可参照 宋金战事、金代地方治理与碑志研究

\noindent\eventanchor{SYD-1206-JIN-221}{九月壬午，金兵攻奪和尚原}\textbf{南宋与金。}
南宋与金 在 1206年（南宋開禧二年、金泰和6年；公元换算据两国年号对照） 的材料痕迹是〔九月壬午，金兵攻奪和尚原〕，出处为 《宋史》卷038〈開禧二年〉；《金史》本纪同年条（互校）。材料未将地点细化到城址，不能用中枢所在地替它补足。

把本项放回事件链，能够看到的前提为：本纪从朝廷视角记录军政行动；出兵与战果需分辨。 而它所打开或收紧的下一步是：后续路线、控制与损失应与对方记录和遗址材料复核。

证据的边界是：《宋史》为元末官修，本纪保存宋廷纪年与处置；战果、数字、地方执行及对方动机须以编年、会要、对方史料或地方材料互校。 宋金战事、金代地方治理与碑志研究 可用于继续检验这一结论。

\noindent\eventanchor{SYD-1206-JIN-249}{是月，濠州、安豐軍及邊屯皆爲金人所破}\textbf{南宋与金。}
从 《宋史》卷038〈開禧二年〉；《金史》本纪同年条（互校） 的 1206年（南宋開禧二年、金泰和6年；公元换算据两国年号对照） 条目看，〔是月，濠州、安豐軍及邊屯皆爲金人所破〕使 南宋与金 的一个具体行动进入叙事；题名保留的城、路、水道或机构名称，是目前可以落地的空间线索。

它承接的条件是：本纪从朝廷视角记录军政行动；出兵与战果需分辨。 随后可追踪的变化为：后续路线、控制与损失应与对方记录和遗址材料复核。

《宋史》为元末官修，本纪保存宋廷纪年与处置；战果、数字、地方执行及对方动机须以编年、会要、对方史料或地方材料互校。 因此，宋金战事、金代地方治理与碑志研究 只能扩展问题，不能代替未保存的细节。

\subsection{1207年（南宋開禧三年、金泰和7年；公元换算据两国年号对照）}
\noindent\eventanchor{SYD-1207-JIN-086}{癸-{丑}-，金人復破隨州}\textbf{南宋与金。}
1207年（南宋開禧三年、金泰和7年；公元换算据两国年号对照） 的记载将 南宋与金 与〔癸-{丑}-，金人復破隨州〕相连。此处可确认的是这项被记录的行动；题名保留的城、路、水道或机构名称，是目前可以落地的空间线索。

前一层压力可由以下线索辨认：本纪年条记录政令或机构处置；实施背景须参照制度材料。 这一处置留下的后续条件是：可在同年及后续政令中追踪；条文不单独证明地方执行。

关于可说范围，须保留：《宋史》为元末官修，本纪保存宋廷纪年与处置；战果、数字、地方执行及对方动机须以编年、会要、对方史料或地方材料互校。 进一步讨论可参照 宋金战事、金代地方治理与碑志研究

\noindent\eventanchor{SYD-1207-JIN-149}{癸亥，以林拱辰爲金國通謝使，遣富琯使金告哀，劉彌正賀金主生辰}\textbf{南宋与金。}
〔癸亥，以林拱辰爲金國通謝使，遣富琯使金告哀，劉彌正賀金主生辰〕见于 《宋史》卷038〈開禧三年〉；《金史》本纪同年条（互校） 的 1207年（南宋開禧三年、金泰和7年；公元换算据两国年号对照） 记录。南宋与金 的选择因此有了可回查的时间位置；材料未将地点细化到城址，不能用中枢所在地替它补足。

本纪记录使节、贡赐或往来，不等于稳定统属关系。 记录所见的结果则是：可与对方编年和交通、文书材料核对其后续落实。 日期相邻不等于因果已经证实。

证据的边界是：《宋史》为元末官修，本纪保存宋廷纪年与处置；战果、数字、地方执行及对方动机须以编年、会要、对方史料或地方材料互校。 宋金战事、金代地方治理与碑志研究 可用于继续检验这一结论。

\noindent\eventanchor{SYD-1207-JIN-187}{癸酉，金人復破大散關}\textbf{南宋与金。}
1207年（南宋開禧三年、金泰和7年；公元换算据两国年号对照），《宋史》卷038〈開禧三年〉；《金史》本纪同年条（互校） 所见的〔癸酉，金人復破大散關〕把 南宋与金 的处置留在可复核的年序中；题名保留的城、路、水道或机构名称，是目前可以落地的空间线索。

把本项放回事件链，能够看到的前提为：本纪年条记录政令或机构处置；实施背景须参照制度材料。 而它所打开或收紧的下一步是：可在同年及后续政令中追踪；条文不单独证明地方执行。

《宋史》为元末官修，本纪保存宋廷纪年与处置；战果、数字、地方执行及对方动机须以编年、会要、对方史料或地方材料互校。 因此，宋金战事、金代地方治理与碑志研究 只能扩展问题，不能代替未保存的细节。

\noindent\eventanchor{SYD-1207-JIN-222}{己未，奉使金國通謝、國信所參議官方信孺發行在}\textbf{南宋与金。}
南宋与金 在 1207年（南宋開禧三年、金泰和7年；公元换算据两国年号对照） 的材料痕迹是〔己未，奉使金國通謝、國信所參議官方信孺發行在〕，出处为 《宋史》卷038〈開禧三年〉；《金史》本纪同年条（互校）。材料未将地点细化到城址，不能用中枢所在地替它补足。

它承接的条件是：本纪年条提供日期和中枢可见行动；前因不据单条补定。 随后可追踪的变化为：可回查同年及后续条目，地方影响仍须另证。

关于可说范围，须保留：《宋史》为元末官修，本纪保存宋廷纪年与处置；战果、数字、地方执行及对方动机须以编年、会要、对方史料或地方材料互校。 进一步讨论可参照 宋金战事、金代地方治理与碑志研究

\noindent\eventanchor{SYD-1207-JIN-250}{四川宣撫副使司參贊軍事楊巨源與金人戰於長橋，敗績}\textbf{南宋与金。}
从 《宋史》卷038〈開禧三年〉；《金史》本纪同年条（互校） 的 1207年（南宋開禧三年、金泰和7年；公元换算据两国年号对照） 条目看，〔四川宣撫副使司參贊軍事楊巨源與金人戰於長橋，敗績〕使 南宋与金 的一个具体行动进入叙事；题名保留的城、路、水道或机构名称，是目前可以落地的空间线索。

前一层压力可由以下线索辨认：本纪从朝廷视角记录军政行动；出兵与战果需分辨。 这一处置留下的后续条件是：后续路线、控制与损失应与对方记录和遗址材料复核。

证据的边界是：《宋史》为元末官修，本纪保存宋廷纪年与处置；战果、数字、地方执行及对方动机须以编年、会要、对方史料或地方材料互校。 宋金战事、金代地方治理与碑志研究 可用于继续检验这一结论。

\subsection{1208年（金泰和八年、宋嘉定元年；干支、月日待核；公元换算据双方本纪年号对照）}
\noindent\eventanchor{JIN-1208-EVENT-101}{金宋以嘉定和议结束战事}\textbf{金与南宋。}
1208年（金泰和八年、宋嘉定元年；干支、月日待核；公元换算据双方本纪年号对照） 的记载将 金与南宋 与〔金宋以嘉定和议结束战事〕相连。此处可确认的是这项被记录的行动；材料未将地点细化到城址，不能用中枢所在地替它补足。

金世宗至章宗时期的制度与边境压力里的军事消耗、边境供给与使节往来构成谈判条件，促成“金宋以嘉定和议结束战事”所涉的协议或名义关系。 记录所见的结果则是：“金宋以嘉定和议结束战事”重排了贡赐、使节或停战的制度接口；条款是否执行及边地实际控制不能由盟约文字直接推定。 日期相邻不等于因果已经证实。

《金史》为元修，宋金战报和官修合法性语言各有宣传性；军数、俘获、控制范围和地方执行须以对方及物质材料限缩。 因此，金代国家形成、宋金边境、都城与金蒙转换研究 只能扩展问题，不能代替未保存的细节。

\subsection{1208年（南宋嘉定元年、金泰和8年；公元换算据两国年号对照）}
\noindent\eventanchor{SYD-1208-JIN-087}{丙子，遣鄒應龍賀金主生辰}\textbf{南宋与金。}
〔丙子，遣鄒應龍賀金主生辰〕见于 《宋史》卷039〈嘉定元年〉；《金史》本纪同年条（互校） 的 1208年（南宋嘉定元年、金泰和8年；公元换算据两国年号对照） 记录。南宋与金 的选择因此有了可回查的时间位置；材料未将地点细化到城址，不能用中枢所在地替它补足。

把本项放回事件链，能够看到的前提为：本纪年条提供日期和中枢可见行动；前因不据单条补定。 而它所打开或收紧的下一步是：可回查同年及后续条目，地方影响仍须另证。

关于可说范围，须保留：《宋史》为元末官修，本纪保存宋廷纪年与处置；战果、数字、地方执行及对方动机须以编年、会要、对方史料或地方材料互校。 进一步讨论可参照 宋金战事、金代地方治理与碑志研究

\subsection{1209年（南宋嘉定二年、金大安1年；公元换算据两国年号对照）}
\noindent\eventanchor{SYD-1209-JIN-088}{己巳，遣俞應符賀金主生辰}\textbf{南宋与金。}
1209年（南宋嘉定二年、金大安1年；公元换算据两国年号对照），《宋史》卷039〈嘉定二年〉；《金史》本纪同年条（互校） 所见的〔己巳，遣俞應符賀金主生辰〕把 南宋与金 的处置留在可复核的年序中；材料未将地点细化到城址，不能用中枢所在地替它补足。

它承接的条件是：本纪年条提供日期和中枢可见行动；前因不据单条补定。 随后可追踪的变化为：可回查同年及后续条目，地方影响仍须另证。

证据的边界是：《宋史》为元末官修，本纪保存宋廷纪年与处置；战果、数字、地方执行及对方动机须以编年、会要、对方史料或地方材料互校。 宋金战事、金代地方治理与碑志研究 可用于继续检验这一结论。

\subsection{1210年（南宋嘉定三年、金大安2年；公元换算据两国年号对照）}
\noindent\eventanchor{SYD-1210-JIN-089}{癸亥，遣黃中賀金主生辰}\textbf{南宋与金。}
南宋与金 在 1210年（南宋嘉定三年、金大安2年；公元换算据两国年号对照） 的材料痕迹是〔癸亥，遣黃中賀金主生辰〕，出处为 《宋史》卷039〈嘉定三年〉；《金史》本纪同年条（互校）。材料未将地点细化到城址，不能用中枢所在地替它补足。

前一层压力可由以下线索辨认：本纪年条提供日期和中枢可见行动；前因不据单条补定。 这一处置留下的后续条件是：可回查同年及后续条目，地方影响仍须另证。

《宋史》为元末官修，本纪保存宋廷纪年与处置；战果、数字、地方执行及对方动机须以编年、会要、对方史料或地方材料互校。 因此，宋金战事、金代地方治理与碑志研究 只能扩展问题，不能代替未保存的细节。

\subsection{1211年（金大安三年、宋嘉定四年、蒙古成吉思汗六年；干支、月日待核；公元换算据各方本纪年号对照）}
\noindent\eventanchor{JIN-1211-EVENT-102}{蒙古大举攻金，金北方防线受冲击}\textbf{金与蒙古。}
从 《金史》相关本纪；《宋史》或《三朝北盟会编》《建炎以来系年要录》对应年条；金代碑志与蒙古、高丽材料（据事核） 的 1211年（金大安三年、宋嘉定四年、蒙古成吉思汗六年；干支、月日待核；公元换算据各方本纪年号对照） 条目看，〔蒙古大举攻金，金北方防线受冲击〕使 金与蒙古 的一个具体行动进入叙事；材料未将地点细化到城址，不能用中枢所在地替它补足。

金蒙战争及金宋关系转折中既有的联盟、交通与补给压力，为“蒙古大举攻金，金北方防线受冲击”提供了直接的军事或动员背景。 记录所见的结果则是：“蒙古大举攻金，金北方防线受冲击”改变了相关城镇、道路或政权间的军事选择；伤亡、俘获和控制范围仅可按各方材料分别确认。 日期相邻不等于因果已经证实。

关于可说范围，须保留：《金史》为元修，宋金战报和官修合法性语言各有宣传性；军数、俘获、控制范围和地方执行须以对方及物质材料限缩。 进一步讨论可参照 金代国家形成、宋金边境、都城与金蒙转换研究

\subsection{1211年（南宋嘉定四年、金大安3年；公元换算据两国年号对照）}
\noindent\eventanchor{SYD-1211-JIN-090}{冬十月甲辰，以金國有難，命江淮、京湖、四川制置司謹邊備}\textbf{南宋与金。}
1211年（南宋嘉定四年、金大安3年；公元换算据两国年号对照） 的记载将 南宋与金 与〔冬十月甲辰，以金國有難，命江淮、京湖、四川制置司謹邊備〕相连。此处可确认的是这项被记录的行动；题名保留的城、路、水道或机构名称，是目前可以落地的空间线索。

把本项放回事件链，能够看到的前提为：本纪年条记录政令或机构处置；实施背景须参照制度材料。 而它所打开或收紧的下一步是：可在同年及后续政令中追踪；条文不单独证明地方执行。

证据的边界是：《宋史》为元末官修，本纪保存宋廷纪年与处置；战果、数字、地方执行及对方动机须以编年、会要、对方史料或地方材料互校。 宋金战事、金代地方治理与碑志研究 可用于继续检验这一结论。

\noindent\eventanchor{SYD-1211-JIN-150}{六月丁亥，遣余嶸賀金主生辰，會金國有難，不至而還}\textbf{南宋与金。}
〔六月丁亥，遣余嶸賀金主生辰，會金國有難，不至而還〕见于 《宋史》卷039〈嘉定四年〉；《金史》本纪同年条（互校） 的 1211年（南宋嘉定四年、金大安3年；公元换算据两国年号对照） 记录。南宋与金 的选择因此有了可回查的时间位置；材料未将地点细化到城址，不能用中枢所在地替它补足。

它承接的条件是：本纪年条提供日期和中枢可见行动；前因不据单条补定。 随后可追踪的变化为：可回查同年及后续条目，地方影响仍须另证。

《宋史》为元末官修，本纪保存宋廷纪年与处置；战果、数字、地方执行及对方动机须以编年、会要、对方史料或地方材料互校。 因此，宋金战事、金代地方治理与碑志研究 只能扩展问题，不能代替未保存的细节。

\subsection{1212年（南宋嘉定五年、金崇庆1年；公元换算据两国年号对照）}
\noindent\eventanchor{SYD-1212-JIN-091}{六月癸未，遣傅誠賀金主生辰}\textbf{南宋与金。}
1212年（南宋嘉定五年、金崇庆1年；公元换算据两国年号对照），《宋史》卷039〈嘉定五年〉；《金史》本纪同年条（互校） 所见的〔六月癸未，遣傅誠賀金主生辰〕把 南宋与金 的处置留在可复核的年序中；材料未将地点细化到城址，不能用中枢所在地替它补足。

前一层压力可由以下线索辨认：本纪年条提供日期和中枢可见行动；前因不据单条补定。 这一处置留下的后续条件是：可回查同年及后续条目，地方影响仍须另证。

关于可说范围，须保留：《宋史》为元末官修，本纪保存宋廷纪年与处置；战果、数字、地方执行及对方动机须以编年、会要、对方史料或地方材料互校。 进一步讨论可参照 宋金战事、金代地方治理与碑志研究

\subsection{1213年}
\noindent\eventanchor{JIN-1213-EVENT-103}{金宣宗即位，金廷开始调整军政策略}\textbf{金。}
金 在 1213年 的材料痕迹是〔金宣宗即位，金廷开始调整军政策略〕，出处为 《金史》相关本纪；《宋史》或《三朝北盟会编》《建炎以来系年要录》对应年条；金代碑志与蒙古、高丽材料（据事核）。材料未将地点细化到城址，不能用中枢所在地替它补足。

金蒙战争及金宋关系转折中前任统治的结束与宫廷、部族或官僚的权力安排相互交织，促成“金宣宗即位，金廷开始调整军政策略”这一继承节点。 记录所见的结果则是：继承并不自动消除既有矛盾；“金宣宗即位，金廷开始调整军政策略”后的任官、军权和对外策略需由后续条目追踪。 日期相邻不等于因果已经证实。

证据的边界是：《金史》为元修，宋金战报和官修合法性语言各有宣传性；军数、俘获、控制范围和地方执行须以对方及物质材料限缩。 金代国家形成、宋金边境、都城与金蒙转换研究 可用于继续检验这一结论。

\subsection{1213年（南宋嘉定六年、金贞祐1年；公元换算据两国年号对照）}
\noindent\eventanchor{SYD-1213-JIN-092}{丙戌，以金主新立，命四川謹邊備}\textbf{南宋与金。}
从 《宋史》卷039〈嘉定六年〉；《金史》本纪同年条（互校） 的 1213年（南宋嘉定六年、金贞祐1年；公元换算据两国年号对照） 条目看，〔丙戌，以金主新立，命四川謹邊備〕使 南宋与金 的一个具体行动进入叙事；题名保留的城、路、水道或机构名称，是目前可以落地的空间线索。

把本项放回事件链，能够看到的前提为：本纪年条记录政令或机构处置；实施背景须参照制度材料。 而它所打开或收紧的下一步是：可在同年及后续政令中追踪；条文不单独证明地方执行。

《宋史》为元末官修，本纪保存宋廷纪年与处置；战果、数字、地方执行及对方动机须以编年、会要、对方史料或地方材料互校。 因此，宋金战事、金代地方治理与碑志研究 只能扩展问题，不能代替未保存的细节。

\noindent\eventanchor{SYD-1213-JIN-151}{丁-{丑}-，遣董居誼賀金主生辰，會金國亂，不至而還}\textbf{南宋与金。}
1213年（南宋嘉定六年、金贞祐1年；公元换算据两国年号对照） 的记载将 南宋与金 与〔丁-{丑}-，遣董居誼賀金主生辰，會金國亂，不至而還〕相连。此处可确认的是这项被记录的行动；材料未将地点细化到城址，不能用中枢所在地替它补足。

它承接的条件是：本纪从朝廷视角记录军政行动；出兵与战果需分辨。 随后可追踪的变化为：后续路线、控制与损失应与对方记录和遗址材料复核。

关于可说范围，须保留：《宋史》为元末官修，本纪保存宋廷纪年与处置；战果、数字、地方执行及对方动机须以编年、会要、对方史料或地方材料互校。 进一步讨论可参照 宋金战事、金代地方治理与碑志研究

\noindent\eventanchor{SYD-1213-JIN-188}{戊申，遣真德秀賀金主即位，會金國亂，不至而還}\textbf{南宋与金。}
〔戊申，遣真德秀賀金主即位，會金國亂，不至而還〕见于 《宋史》卷039〈嘉定六年〉；《金史》本纪同年条（互校） 的 1213年（南宋嘉定六年、金贞祐1年；公元换算据两国年号对照） 记录。南宋与金 的选择因此有了可回查的时间位置；材料未将地点细化到城址，不能用中枢所在地替它补足。

前一层压力可由以下线索辨认：本纪从朝廷视角记录军政行动；出兵与战果需分辨。 这一处置留下的后续条件是：后续路线、控制与损失应与对方记录和遗址材料复核。

证据的边界是：《宋史》为元末官修，本纪保存宋廷纪年与处置；战果、数字、地方执行及对方动机须以编年、会要、对方史料或地方材料互校。 宋金战事、金代地方治理与碑志研究 可用于继续检验这一结论。

\subsection{1214年（南宋嘉定七年、金贞祐2年；公元换算据两国年号对照）}
\noindent\eventanchor{SYD-1214-JIN-093}{庚辰，金國來督二年歲幣}\textbf{南宋与金。}
1214年（南宋嘉定七年、金贞祐2年；公元换算据两国年号对照），《宋史》卷039〈嘉定七年〉；《金史》本纪同年条（互校） 所见的〔庚辰，金國來督二年歲幣〕把 南宋与金 的处置留在可复核的年序中；材料未将地点细化到城址，不能用中枢所在地替它补足。

本纪年条提供日期和中枢可见行动；前因不据单条补定。 记录所见的结果则是：可回查同年及后续条目，地方影响仍须另证。 日期相邻不等于因果已经证实。

《宋史》为元末官修，本纪保存宋廷纪年与处置；战果、数字、地方执行及对方动机须以编年、会要、对方史料或地方材料互校。 因此，宋金战事、金代地方治理与碑志研究 只能扩展问题，不能代替未保存的细节。

\noindent\eventanchor{SYD-1214-JIN-152}{庚寅，以起居舍人真德秀奏，罷金國歲幣}\textbf{南宋与金。}
南宋与金 在 1214年（南宋嘉定七年、金贞祐2年；公元换算据两国年号对照） 的材料痕迹是〔庚寅，以起居舍人真德秀奏，罷金國歲幣〕，出处为 《宋史》卷039〈嘉定七年〉；《金史》本纪同年条（互校）。材料未将地点细化到城址，不能用中枢所在地替它补足。

把本项放回事件链，能够看到的前提为：本纪年条记录政令或机构处置；实施背景须参照制度材料。 而它所打开或收紧的下一步是：可在同年及后续政令中追踪；条文不单独证明地方执行。

关于可说范围，须保留：《宋史》为元末官修，本纪保存宋廷纪年与处置；战果、数字、地方执行及对方动机须以编年、会要、对方史料或地方材料互校。 进一步讨论可参照 宋金战事、金代地方治理与碑志研究

\noindent\eventanchor{SYD-1214-JIN-189}{七年春正月丁卯朔，四川制置司遣提舉皂郊博馬務何九齡率諸將及金人戰于秦州城下，敗還}\textbf{南宋与金。}
从 《宋史》卷039〈嘉定七年〉；《金史》本纪同年条（互校） 的 1214年（南宋嘉定七年、金贞祐2年；公元换算据两国年号对照） 条目看，〔七年春正月丁卯朔，四川制置司遣提舉皂郊博馬務何九齡率諸將及金人戰于秦州城下，敗還〕使 南宋与金 的一个具体行动进入叙事；题名保留的城、路、水道或机构名称，是目前可以落地的空间线索。

它承接的条件是：本纪年条记录政令或机构处置；实施背景须参照制度材料。 随后可追踪的变化为：可在同年及后续政令中追踪；条文不单独证明地方执行。

证据的边界是：《宋史》为元末官修，本纪保存宋廷纪年与处置；战果、数字、地方执行及对方动机须以编年、会要、对方史料或地方材料互校。 宋金战事、金代地方治理与碑志研究 可用于继续检验这一结论。

\noindent\eventanchor{SYD-1214-JIN-223}{是月，夏人以書來四川，議夾攻金人，不報}\textbf{南宋与金。}
1214年（南宋嘉定七年、金贞祐2年；公元换算据两国年号对照） 的记载将 南宋与金 与〔是月，夏人以書來四川，議夾攻金人，不報〕相连。此处可确认的是这项被记录的行动；题名保留的城、路、水道或机构名称，是目前可以落地的空间线索。

前一层压力可由以下线索辨认：本纪从朝廷视角记录军政行动；出兵与战果需分辨。 这一处置留下的后续条件是：后续路线、控制与损失应与对方记录和遗址材料复核。

《宋史》为元末官修，本纪保存宋廷纪年与处置；战果、数字、地方执行及对方动机须以编年、会要、对方史料或地方材料互校。 因此，宋金战事、金代地方治理与碑志研究 只能扩展问题，不能代替未保存的细节。

\subsection{1215年}
\noindent\eventanchor{JIN-1215-EVENT-104}{蒙古军攻取中都}\textbf{金与蒙古。}
〔蒙古军攻取中都〕见于 《金史》相关本纪；《宋史》或《三朝北盟会编》《建炎以来系年要录》对应年条；金代碑志与蒙古、高丽材料（据事核） 的 1215年 记录。金与蒙古 的选择因此有了可回查的时间位置；题名保留的城、路、水道或机构名称，是目前可以落地的空间线索。

金蒙战争及金宋关系转折中既有的联盟、交通与补给压力，为“蒙古军攻取中都”提供了直接的军事或动员背景。 记录所见的结果则是：“蒙古军攻取中都”改变了相关城镇、道路或政权间的军事选择；伤亡、俘获和控制范围仅可按各方材料分别确认。 日期相邻不等于因果已经证实。

关于可说范围，须保留：《金史》为元修，宋金战报和官修合法性语言各有宣传性；军数、俘获、控制范围和地方执行须以对方及物质材料限缩。 进一步讨论可参照 金代国家形成、宋金边境、都城与金蒙转换研究

\subsection{1215年（南宋嘉定八年、金贞祐3年；公元换算据两国年号对照）}
\noindent\eventanchor{SYD-1215-JIN-094}{己卯，遣丁焴賀金主生辰}\textbf{南宋与金。}
1215年（南宋嘉定八年、金贞祐3年；公元换算据两国年号对照），《宋史》卷039〈嘉定八年〉；《金史》本纪同年条（互校） 所见的〔己卯，遣丁焴賀金主生辰〕把 南宋与金 的处置留在可复核的年序中；材料未将地点细化到城址，不能用中枢所在地替它补足。

把本项放回事件链，能够看到的前提为：本纪年条提供日期和中枢可见行动；前因不据单条补定。 而它所打开或收紧的下一步是：可回查同年及后续条目，地方影响仍须另证。

证据的边界是：《宋史》为元末官修，本纪保存宋廷纪年与处置；战果、数字、地方执行及对方动机须以编年、会要、对方史料或地方材料互校。 宋金战事、金代地方治理与碑志研究 可用于继续检验这一结论。

\subsection{1216年（南宋嘉定九年、金贞祐4年；公元换算据两国年号对照）}
\noindent\eventanchor{SYD-1216-JIN-095}{乙亥，遣留筠賀金主生辰}\textbf{南宋与金。}
南宋与金 在 1216年（南宋嘉定九年、金贞祐4年；公元换算据两国年号对照） 的材料痕迹是〔乙亥，遣留筠賀金主生辰〕，出处为 《宋史》卷039〈嘉定九年〉；《金史》本纪同年条（互校）。材料未将地点细化到城址，不能用中枢所在地替它补足。

它承接的条件是：本纪年条提供日期和中枢可见行动；前因不据单条补定。 随后可追踪的变化为：可回查同年及后续条目，地方影响仍须另证。

《宋史》为元末官修，本纪保存宋廷纪年与处置；战果、数字、地方执行及对方动机须以编年、会要、对方史料或地方材料互校。 因此，宋金战事、金代地方治理与碑志研究 只能扩展问题，不能代替未保存的细节。

\subsection{1217年}
\noindent\eventanchor{JIN-1217-EVENT-105}{金转而进攻南宋边境}\textbf{金与南宋。}
从 《金史》相关本纪；《宋史》或《三朝北盟会编》《建炎以来系年要录》对应年条；金代碑志与蒙古、高丽材料（据事核） 的 1217年 条目看，〔金转而进攻南宋边境〕使 金与南宋 的一个具体行动进入叙事；材料未将地点细化到城址，不能用中枢所在地替它补足。

前一层压力可由以下线索辨认：金蒙战争及金宋关系转折中既有的联盟、交通与补给压力，为“金转而进攻南宋边境”提供了直接的军事或动员背景。 这一处置留下的后续条件是：“金转而进攻南宋边境”改变了相关城镇、道路或政权间的军事选择；伤亡、俘获和控制范围仅可按各方材料分别确认。

关于可说范围，须保留：《金史》为元修，宋金战报和官修合法性语言各有宣传性；军数、俘获、控制范围和地方执行须以对方及物质材料限缩。 进一步讨论可参照 金代国家形成、宋金边境、都城与金蒙转换研究

\subsection{1217年（南宋嘉定十年、金兴定1年；公元换算据两国年号对照）}
\noindent\eventanchor{SYD-1217-JIN-096}{庚子，遣錢撫賀金主生辰}\textbf{南宋与金。}
1217年（南宋嘉定十年、金兴定1年；公元换算据两国年号对照） 的记载将 南宋与金 与〔庚子，遣錢撫賀金主生辰〕相连。此处可确认的是这项被记录的行动；材料未将地点细化到城址，不能用中枢所在地替它补足。

本纪所记财用、赈济或征发属于特定报告场景。 记录所见的结果则是：可与食货志、文书和物质材料比较，不外推为全境总量。 日期相邻不等于因果已经证实。

证据的边界是：《宋史》为元末官修，本纪保存宋廷纪年与处置；战果、数字、地方执行及对方动机须以编年、会要、对方史料或地方材料互校。 宋金战事、金代地方治理与碑志研究 可用于继续检验这一结论。

\noindent\eventanchor{SYD-1217-JIN-153}{夏四月丁未朔，金人犯光州中渡鎮，執榷場官盛允升殺之，遂分兵犯樊城}\textbf{南宋与金。}
〔夏四月丁未朔，金人犯光州中渡鎮，執榷場官盛允升殺之，遂分兵犯樊城〕见于 《宋史》卷040〈嘉定十年〉；《金史》本纪同年条（互校） 的 1217年（南宋嘉定十年、金兴定1年；公元换算据两国年号对照） 记录。南宋与金 的选择因此有了可回查的时间位置；题名保留的城、路、水道或机构名称，是目前可以落地的空间线索。

把本项放回事件链，能够看到的前提为：本纪年条记录政令或机构处置；实施背景须参照制度材料。 而它所打开或收紧的下一步是：可在同年及后续政令中追踪；条文不单独证明地方执行。

《宋史》为元末官修，本纪保存宋廷纪年与处置；战果、数字、地方执行及对方动机须以编年、会要、对方史料或地方材料互校。 因此，宋金战事、金代地方治理与碑志研究 只能扩展问题，不能代替未保存的细节。

\noindent\eventanchor{SYD-1217-JIN-190}{辛酉，廬州鈐轄王辛敗金人於光山縣之安昌砦，殺其統軍完顏掩}\textbf{南宋与金。}
1217年（南宋嘉定十年、金兴定1年；公元换算据两国年号对照），《宋史》卷040〈嘉定十年〉；《金史》本纪同年条（互校） 所见的〔辛酉，廬州鈐轄王辛敗金人於光山縣之安昌砦，殺其統軍完顏掩〕把 南宋与金 的处置留在可复核的年序中；题名保留的城、路、水道或机构名称，是目前可以落地的空间线索。

它承接的条件是：本纪从朝廷视角记录军政行动；出兵与战果需分辨。 随后可追踪的变化为：后续路线、控制与损失应与对方记录和遗址材料复核。

关于可说范围，须保留：《宋史》为元末官修，本纪保存宋廷纪年与处置；战果、数字、地方执行及对方动机须以编年、会要、对方史料或地方材料互校。 进一步讨论可参照 宋金战事、金代地方治理与碑志研究

\subsection{1218年（南宋嘉定十一年、金兴定2年；公元换算据两国年号对照）}
\noindent\eventanchor{SYD-1218-JIN-097}{丙午，金人破皂郊，死者五萬人}\textbf{南宋与金。}
南宋与金 在 1218年（南宋嘉定十一年、金兴定2年；公元换算据两国年号对照） 的材料痕迹是〔丙午，金人破皂郊，死者五萬人〕，出处为 《宋史》卷040〈嘉定十一年〉；《金史》本纪同年条（互校）。材料未将地点细化到城址，不能用中枢所在地替它补足。

前一层压力可由以下线索辨认：本纪从朝廷视角记录军政行动；出兵与战果需分辨。 这一处置留下的后续条件是：后续路线、控制与损失应与对方记录和遗址材料复核。

证据的边界是：《宋史》为元末官修，本纪保存宋廷纪年与处置；战果、数字、地方执行及对方动机须以编年、会要、对方史料或地方材料互校。 宋金战事、金代地方治理与碑志研究 可用于继续检验这一结论。

\noindent\eventanchor{SYD-1218-JIN-154}{楚州鈐轄梁昭祖焚金人糧舟於大清河，京東忠義副都統沈鐸遣兵助之}\textbf{南宋与金。}
从 《宋史》卷040〈嘉定十一年〉；《金史》本纪同年条（互校） 的 1218年（南宋嘉定十一年、金兴定2年；公元换算据两国年号对照） 条目看，〔楚州鈐轄梁昭祖焚金人糧舟於大清河，京東忠義副都統沈鐸遣兵助之〕使 南宋与金 的一个具体行动进入叙事；题名保留的城、路、水道或机构名称，是目前可以落地的空间线索。

本纪保留灾害或工程的报告地点与朝廷处置；灾情范围需另证。 记录所见的结果则是：可与地方文书、河工和环境材料比较，不将单点记录外推为全域因果。 日期相邻不等于因果已经证实。

《宋史》为元末官修，本纪保存宋廷纪年与处置；战果、数字、地方执行及对方动机须以编年、会要、对方史料或地方材料互校。 因此，宋金战事、金代地方治理与碑志研究 只能扩展问题，不能代替未保存的细节。

\noindent\eventanchor{SYD-1218-JIN-191}{十一月壬申，金人攻安豐軍之黃口灘}\textbf{南宋与金。}
1218年（南宋嘉定十一年、金兴定2年；公元换算据两国年号对照） 的记载将 南宋与金 与〔十一月壬申，金人攻安豐軍之黃口灘〕相连。此处可确认的是这项被记录的行动；材料未将地点细化到城址，不能用中枢所在地替它补足。

把本项放回事件链，能够看到的前提为：本纪从朝廷视角记录军政行动；出兵与战果需分辨。 而它所打开或收紧的下一步是：后续路线、控制与损失应与对方记录和遗址材料复核。

关于可说范围，须保留：《宋史》为元末官修，本纪保存宋廷纪年与处置；战果、数字、地方执行及对方动机须以编年、会要、对方史料或地方材料互校。 进一步讨论可参照 宋金战事、金代地方治理与碑志研究

\noindent\eventanchor{SYD-1218-JIN-224}{戊午，金人復犯大散關，守將王立遁}\textbf{南宋与金。}
〔戊午，金人復犯大散關，守將王立遁〕见于 《宋史》卷040〈嘉定十一年〉；《金史》本纪同年条（互校） 的 1218年（南宋嘉定十一年、金兴定2年；公元换算据两国年号对照） 记录。南宋与金 的选择因此有了可回查的时间位置；题名保留的城、路、水道或机构名称，是目前可以落地的空间线索。

它承接的条件是：本纪年条记录政令或机构处置；实施背景须参照制度材料。 随后可追踪的变化为：可在同年及后续政令中追踪；条文不单独证明地方执行。

证据的边界是：《宋史》为元末官修，本纪保存宋廷纪年与处置；战果、数字、地方执行及对方动机须以编年、会要、对方史料或地方材料互校。 宋金战事、金代地方治理与碑志研究 可用于继续检验这一结论。

\subsection{1219年（南宋嘉定十二年、金兴定3年；公元换算据两国年号对照）}
\noindent\eventanchor{SYD-1219-JIN-098}{二月戊戌朔，金人破光山縣}\textbf{南宋与金。}
1219年（南宋嘉定十二年、金兴定3年；公元换算据两国年号对照），《宋史》卷040〈嘉定十二年〉；《金史》本纪同年条（互校） 所见的〔二月戊戌朔，金人破光山縣〕把 南宋与金 的处置留在可复核的年序中；题名保留的城、路、水道或机构名称，是目前可以落地的空间线索。

前一层压力可由以下线索辨认：本纪从朝廷视角记录军政行动；出兵与战果需分辨。 这一处置留下的后续条件是：后续路线、控制与损失应与对方记录和遗址材料复核。

《宋史》为元末官修，本纪保存宋廷纪年与处置；战果、数字、地方执行及对方动机须以编年、会要、对方史料或地方材料互校。 因此，宋金战事、金代地方治理与碑志研究 只能扩展问题，不能代替未保存的细节。

\noindent\eventanchor{SYD-1219-JIN-155}{庚寅，金人犯隨州、棗陽軍，又破信陽軍之二砦，京西諸將引兵拒之}\textbf{南宋与金。}
南宋与金 在 1219年（南宋嘉定十二年、金兴定3年；公元换算据两国年号对照） 的材料痕迹是〔庚寅，金人犯隨州、棗陽軍，又破信陽軍之二砦，京西諸將引兵拒之〕，出处为 《宋史》卷040〈嘉定十二年〉；《金史》本纪同年条（互校）。题名保留的城、路、水道或机构名称，是目前可以落地的空间线索。

本纪从朝廷视角记录军政行动；出兵与战果需分辨。 记录所见的结果则是：后续路线、控制与损失应与对方记录和遗址材料复核。 日期相邻不等于因果已经证实。

关于可说范围，须保留：《宋史》为元末官修，本纪保存宋廷纪年与处置；战果、数字、地方执行及对方动机须以编年、会要、对方史料或地方材料互校。 进一步讨论可参照 宋金战事、金代地方治理与碑志研究

\noindent\eventanchor{SYD-1219-JIN-192}{癸卯，金人破武休關，興元都統李貴遁還，利州路提刑、權興元府事趙希昔棄城去}\textbf{南宋与金。}
从 《宋史》卷040〈嘉定十二年〉；《金史》本纪同年条（互校） 的 1219年（南宋嘉定十二年、金兴定3年；公元换算据两国年号对照） 条目看，〔癸卯，金人破武休關，興元都統李貴遁還，利州路提刑、權興元府事趙希昔棄城去〕使 南宋与金 的一个具体行动进入叙事；题名保留的城、路、水道或机构名称，是目前可以落地的空间线索。

把本项放回事件链，能够看到的前提为：本纪从朝廷视角记录军政行动；出兵与战果需分辨。 而它所打开或收紧的下一步是：后续路线、控制与损失应与对方记录和遗址材料复核。

证据的边界是：《宋史》为元末官修，本纪保存宋廷纪年与处置；战果、数字、地方执行及对方动机须以编年、会要、对方史料或地方材料互校。 宋金战事、金代地方治理与碑志研究 可用于继续检验这一结论。

\noindent\eventanchor{SYD-1219-JIN-225}{癸巳，丁焴復以書約夏國攻金人}\textbf{南宋与金。}
1219年（南宋嘉定十二年、金兴定3年；公元换算据两国年号对照） 的记载将 南宋与金 与〔癸巳，丁焴復以書約夏國攻金人〕相连。此处可确认的是这项被记录的行动；材料未将地点细化到城址，不能用中枢所在地替它补足。

它承接的条件是：本纪年条记录政令或机构处置；实施背景须参照制度材料。 随后可追踪的变化为：可在同年及后续政令中追踪；条文不单独证明地方执行。

《宋史》为元末官修，本纪保存宋廷纪年与处置；战果、数字、地方执行及对方动机须以编年、会要、对方史料或地方材料互校。 因此，宋金战事、金代地方治理与碑志研究 只能扩展问题，不能代替未保存的细节。

\subsection{1220年（南宋嘉定十三年、金兴定4年；公元换算据两国年号对照）}
\noindent\eventanchor{SYD-1220-JIN-099}{庚戌，金人犯皂郊堡，沔州統制董炤等與戰，大敗}\textbf{南宋与金。}
〔庚戌，金人犯皂郊堡，沔州統制董炤等與戰，大敗〕见于 《宋史》卷040〈嘉定十三年〉；《金史》本纪同年条（互校） 的 1220年（南宋嘉定十三年、金兴定4年；公元换算据两国年号对照） 记录。南宋与金 的选择因此有了可回查的时间位置；题名保留的城、路、水道或机构名称，是目前可以落地的空间线索。

前一层压力可由以下线索辨认：本纪从朝廷视角记录军政行动；出兵与战果需分辨。 这一处置留下的后续条件是：后续路线、控制与损失应与对方记录和遗址材料复核。

关于可说范围，须保留：《宋史》为元末官修，本纪保存宋廷纪年与处置；战果、数字、地方执行及对方动机须以编年、会要、对方史料或地方材料互校。 进一步讨论可参照 宋金战事、金代地方治理与碑志研究

\noindent\eventanchor{SYD-1220-JIN-156}{金人追之，遂攻樊城，趙方督諸將拒退之}\textbf{南宋与金。}
1220年（南宋嘉定十三年、金兴定4年；公元换算据两国年号对照），《宋史》卷040〈嘉定十三年〉；《金史》本纪同年条（互校） 所见的〔金人追之，遂攻樊城，趙方督諸將拒退之〕把 南宋与金 的处置留在可复核的年序中；题名保留的城、路、水道或机构名称，是目前可以落地的空间线索。

本纪从朝廷视角记录军政行动；出兵与战果需分辨。 记录所见的结果则是：后续路线、控制与损失应与对方记录和遗址材料复核。 日期相邻不等于因果已经证实。

证据的边界是：《宋史》为元末官修，本纪保存宋廷纪年与处置；战果、数字、地方执行及对方动机须以编年、会要、对方史料或地方材料互校。 宋金战事、金代地方治理与碑志研究 可用于继续检验这一结论。

\noindent\eventanchor{SYD-1220-JIN-193}{壬申，安丙遺夏人書，定議夾攻金人}\textbf{南宋与金。}
南宋与金 在 1220年（南宋嘉定十三年、金兴定4年；公元换算据两国年号对照） 的材料痕迹是〔壬申，安丙遺夏人書，定議夾攻金人〕，出处为 《宋史》卷040〈嘉定十三年〉；《金史》本纪同年条（互校）。材料未将地点细化到城址，不能用中枢所在地替它补足。

把本项放回事件链，能够看到的前提为：本纪从朝廷视角记录军政行动；出兵与战果需分辨。 而它所打开或收紧的下一步是：后续路线、控制与损失应与对方记录和遗址材料复核。

《宋史》为元末官修，本纪保存宋廷纪年与处置；战果、数字、地方执行及对方动机须以编年、会要、对方史料或地方材料互校。 因此，宋金战事、金代地方治理与碑志研究 只能扩展问题，不能代替未保存的细节。

\noindent\eventanchor{SYD-1220-JIN-226}{壬寅，質俊等自來遠鎮進攻定邊城，金人來救，俊等擊破之}\textbf{南宋与金。}
从 《宋史》卷040〈嘉定十三年〉；《金史》本纪同年条（互校） 的 1220年（南宋嘉定十三年、金兴定4年；公元换算据两国年号对照） 条目看，〔壬寅，質俊等自來遠鎮進攻定邊城，金人來救，俊等擊破之〕使 南宋与金 的一个具体行动进入叙事；题名保留的城、路、水道或机构名称，是目前可以落地的空间线索。

它承接的条件是：本纪从朝廷视角记录军政行动；出兵与战果需分辨。 随后可追踪的变化为：后续路线、控制与损失应与对方记录和遗址材料复核。

关于可说范围，须保留：《宋史》为元末官修，本纪保存宋廷纪年与处置；战果、数字、地方执行及对方动机须以编年、会要、对方史料或地方材料互校。 进一步讨论可参照 宋金战事、金代地方治理与碑志研究

\subsection{1221年（南宋嘉定十四年、金兴定5年；公元换算据两国年号对照）}
\noindent\eventanchor{SYD-1221-JIN-100}{丁亥，金人破黃州，淮西提刑、知州事何大節棄城遁死}\textbf{南宋与金。}
1221年（南宋嘉定十四年、金兴定5年；公元换算据两国年号对照） 的记载将 南宋与金 与〔丁亥，金人破黃州，淮西提刑、知州事何大節棄城遁死〕相连。此处可确认的是这项被记录的行动；题名保留的城、路、水道或机构名称，是目前可以落地的空间线索。

前一层压力可由以下线索辨认：本纪从朝廷视角记录军政行动；出兵与战果需分辨。 这一处置留下的后续条件是：后续路线、控制与损失应与对方记录和遗址材料复核。

证据的边界是：《宋史》为元末官修，本纪保存宋廷纪年与处置；战果、数字、地方执行及对方动机须以编年、会要、对方史料或地方材料互校。 宋金战事、金代地方治理与碑志研究 可用于继续检验这一结论。

\noindent\eventanchor{SYD-1221-JIN-157}{癸-{丑}-，金人退師，扈再興邀擊，敗之於天長鎮，甲寅晦，又敗之}\textbf{南宋与金。}
〔癸-{丑}-，金人退師，扈再興邀擊，敗之於天長鎮，甲寅晦，又敗之〕见于 《宋史》卷040〈嘉定十四年〉；《金史》本纪同年条（互校） 的 1221年（南宋嘉定十四年、金兴定5年；公元换算据两国年号对照） 记录。南宋与金 的选择因此有了可回查的时间位置；材料未将地点细化到城址，不能用中枢所在地替它补足。

本纪从朝廷视角记录军政行动；出兵与战果需分辨。 记录所见的结果则是：后续路线、控制与损失应与对方记录和遗址材料复核。 日期相邻不等于因果已经证实。

《宋史》为元末官修，本纪保存宋廷纪年与处置；战果、数字、地方执行及对方动机须以编年、会要、对方史料或地方材料互校。 因此，宋金战事、金代地方治理与碑志研究 只能扩展问题，不能代替未保存的细节。

\noindent\eventanchor{SYD-1221-JIN-194}{己亥，金人陷蘄州，知州事李誠之及其家人、官屬皆死之}\textbf{南宋与金。}
1221年（南宋嘉定十四年、金兴定5年；公元换算据两国年号对照），《宋史》卷040〈嘉定十四年〉；《金史》本纪同年条（互校） 所见的〔己亥，金人陷蘄州，知州事李誠之及其家人、官屬皆死之〕把 南宋与金 的处置留在可复核的年序中；题名保留的城、路、水道或机构名称，是目前可以落地的空间线索。

把本项放回事件链，能够看到的前提为：本纪年条提供日期和中枢可见行动；前因不据单条补定。 而它所打开或收紧的下一步是：可回查同年及后续条目，地方影响仍须另证。

关于可说范围，须保留：《宋史》为元末官修，本纪保存宋廷纪年与处置；战果、数字、地方执行及对方动机须以编年、会要、对方史料或地方材料互校。 进一步讨论可参照 宋金战事、金代地方治理与碑志研究

\noindent\eventanchor{SYD-1221-JIN-227}{壬申，金人治舟於團風，弗克濟，遂圍黃州，分兵破諸縣，又遣別將犯漢陽軍}\textbf{南宋与金。}
南宋与金 在 1221年（南宋嘉定十四年、金兴定5年；公元换算据两国年号对照） 的材料痕迹是〔壬申，金人治舟於團風，弗克濟，遂圍黃州，分兵破諸縣，又遣別將犯漢陽軍〕，出处为 《宋史》卷040〈嘉定十四年〉；《金史》本纪同年条（互校）。题名保留的城、路、水道或机构名称，是目前可以落地的空间线索。

它承接的条件是：本纪从朝廷视角记录军政行动；出兵与战果需分辨。 随后可追踪的变化为：后续路线、控制与损失应与对方记录和遗址材料复核。

证据的边界是：《宋史》为元末官修，本纪保存宋廷纪年与处置；战果、数字、地方执行及对方动机须以编年、会要、对方史料或地方材料互校。 宋金战事、金代地方治理与碑志研究 可用于继续检验这一结论。

\subsection{1224年}
\noindent\eventanchor{JIN-1224-EVENT-106}{金哀宗即位}\textbf{金。}
从 《金史》相关本纪；《宋史》或《三朝北盟会编》《建炎以来系年要录》对应年条；金代碑志与蒙古、高丽材料（据事核） 的 1224年 条目看，〔金哀宗即位〕使 金 的一个具体行动进入叙事；材料未将地点细化到城址，不能用中枢所在地替它补足。

前一层压力可由以下线索辨认：金蒙战争及金宋关系转折中前任统治的结束与宫廷、部族或官僚的权力安排相互交织，促成“金哀宗即位”这一继承节点。 这一处置留下的后续条件是：继承并不自动消除既有矛盾；“金哀宗即位”后的任官、军权和对外策略需由后续条目追踪。

《金史》为元修，宋金战报和官修合法性语言各有宣传性；军数、俘获、控制范围和地方执行须以对方及物质材料限缩。 因此，金代国家形成、宋金边境、都城与金蒙转换研究 只能扩展问题，不能代替未保存的细节。

\subsection{1224年（南宋嘉定十七年、金正大1年；公元换算据两国年号对照）}
\noindent\eventanchor{SYD-1224-JIN-101}{癸酉，知西和州尚震午坐金兵至謀遁，奪三官、岳州居住}\textbf{南宋与金。}
1224年（南宋嘉定十七年、金正大1年；公元换算据两国年号对照） 的记载将 南宋与金 与〔癸酉，知西和州尚震午坐金兵至謀遁，奪三官、岳州居住〕相连。此处可确认的是这项被记录的行动；题名保留的城、路、水道或机构名称，是目前可以落地的空间线索。

本纪保留灾害或工程的报告地点与朝廷处置；灾情范围需另证。 记录所见的结果则是：可与地方文书、河工和环境材料比较，不将单点记录外推为全域因果。 日期相邻不等于因果已经证实。

关于可说范围，须保留：《宋史》为元末官修，本纪保存宋廷纪年与处置；战果、数字、地方执行及对方动机须以编年、会要、对方史料或地方材料互校。 进一步讨论可参照 宋金战事、金代地方治理与碑志研究

\subsection{1230年}
\noindent\eventanchor{JIN-1230-EVENT-107}{蒙古重新集中兵力攻金}\textbf{金与蒙古。}
〔蒙古重新集中兵力攻金〕见于 《金史》相关本纪；《宋史》或《三朝北盟会编》《建炎以来系年要录》对应年条；金代碑志与蒙古、高丽材料（据事核） 的 1230年 记录。金与蒙古 的选择因此有了可回查的时间位置；材料未将地点细化到城址，不能用中枢所在地替它补足。

把本项放回事件链，能够看到的前提为：金蒙战争及金宋关系转折中既有的联盟、交通与补给压力，为“蒙古重新集中兵力攻金”提供了直接的军事或动员背景。 而它所打开或收紧的下一步是：“蒙古重新集中兵力攻金”改变了相关城镇、道路或政权间的军事选择；伤亡、俘获和控制范围仅可按各方材料分别确认。

证据的边界是：《金史》为元修，宋金战报和官修合法性语言各有宣传性；军数、俘获、控制范围和地方执行须以对方及物质材料限缩。 金代国家形成、宋金边境、都城与金蒙转换研究 可用于继续检验这一结论。

\subsection{1232年}
\noindent\eventanchor{JIN-1232-EVENT-108}{蒙古围攻汴京}\textbf{金与蒙古。}
1232年，《金史》相关本纪；《宋史》或《三朝北盟会编》《建炎以来系年要录》对应年条；金代碑志与蒙古、高丽材料（据事核） 所见的〔蒙古围攻汴京〕把 金与蒙古 的处置留在可复核的年序中；题名保留的城、路、水道或机构名称，是目前可以落地的空间线索。

它承接的条件是：金蒙战争及金宋关系转折中既有的联盟、交通与补给压力，为“蒙古围攻汴京”提供了直接的军事或动员背景。 随后可追踪的变化为：“蒙古围攻汴京”改变了相关城镇、道路或政权间的军事选择；伤亡、俘获和控制范围仅可按各方材料分别确认。

《金史》为元修，宋金战报和官修合法性语言各有宣传性；军数、俘获、控制范围和地方执行须以对方及物质材料限缩。 因此，金代国家形成、宋金边境、都城与金蒙转换研究 只能扩展问题，不能代替未保存的细节。

\subsection{1232年（南宋紹定五年、金天兴1年；公元换算据两国年号对照）}
\noindent\eventanchor{SYD-1232-JIN-102}{時宋與大元兵合圍汴京，金主奔歸德府，尋奔蔡州，大元再遣使議攻金，史嵩之以鄒伸之報謝}\textbf{南宋与金。}
南宋与金 在 1232年（南宋紹定五年、金天兴1年；公元换算据两国年号对照） 的材料痕迹是〔時宋與大元兵合圍汴京，金主奔歸德府，尋奔蔡州，大元再遣使議攻金，史嵩之以鄒伸之報謝〕，出处为 《宋史》卷041〈紹定五年〉；《金史》本纪同年条（互校）。题名保留的城、路、水道或机构名称，是目前可以落地的空间线索。

前一层压力可由以下线索辨认：本纪从朝廷视角记录军政行动；出兵与战果需分辨。 这一处置留下的后续条件是：后续路线、控制与损失应与对方记录和遗址材料复核。

关于可说范围，须保留：《宋史》为元末官修，本纪保存宋廷纪年与处置；战果、数字、地方执行及对方动机须以编年、会要、对方史料或地方材料互校。 进一步讨论可参照 宋金战事、金代地方治理与碑志研究

\subsection{1233年}
\noindent\eventanchor{JIN-1233-EVENT-109}{金哀宗退至蔡州}\textbf{金。}
从 《金史》相关本纪；《宋史》或《三朝北盟会编》《建炎以来系年要录》对应年条；金代碑志与蒙古、高丽材料（据事核） 的 1233年 条目看，〔金哀宗退至蔡州〕使 金 的一个具体行动进入叙事；题名保留的城、路、水道或机构名称，是目前可以落地的空间线索。

“金哀宗退至蔡州”发生在金蒙战争及金宋关系转折的具体压力与机会之中，相关行动者的选择不能脱离此前事件链解释。 记录所见的结果则是：“金哀宗退至蔡州”为后续政治、边防或资源配置留下条件；其社会影响与实际覆盖范围仍须以异类材料限定。 日期相邻不等于因果已经证实。

证据的边界是：《金史》为元修，宋金战报和官修合法性语言各有宣传性；军数、俘获、控制范围和地方执行须以对方及物质材料限缩。 金代国家形成、宋金边境、都城与金蒙转换研究 可用于继续检验这一结论。

\subsection{1234年（金天兴三年、宋端平元年、蒙古窝阔台六年；干支、月日待核；公元换算据各方本纪年号对照）}
\noindent\eventanchor{JIN-1234-EVENT-110}{蒙古与南宋军攻破蔡州，金朝灭亡}\textbf{金南宋蒙古。}
1234年（金天兴三年、宋端平元年、蒙古窝阔台六年；干支、月日待核；公元换算据各方本纪年号对照） 的记载将 金南宋蒙古 与〔蒙古与南宋军攻破蔡州，金朝灭亡〕相连。此处可确认的是这项被记录的行动；题名保留的城、路、水道或机构名称，是目前可以落地的空间线索。

把本项放回事件链，能够看到的前提为：金蒙战争及金宋关系转折中既有的联盟、交通与补给压力，为“蒙古与南宋军攻破蔡州，金朝灭亡”提供了直接的军事或动员背景。 而它所打开或收紧的下一步是：“蒙古与南宋军攻破蔡州，金朝灭亡”改变了相关城镇、道路或政权间的军事选择；伤亡、俘获和控制范围仅可按各方材料分别确认。

《金史》为元修，宋金战报和官修合法性语言各有宣传性；军数、俘获、控制范围和地方执行须以对方及物质材料限缩。 因此，金代国家形成、宋金边境、都城与金蒙转换研究 只能扩展问题，不能代替未保存的细节。
